% ===========================================================================
%  C 语言入门：从计算机组成到指针
%  电控算法组新生培训教程
%  用 ElegantBook 模板排版（lang=cn, scheme=chinese, color=blue）
%  编译：xelatex main.tex （连编两次，目录与交叉引用才正确）
% ===========================================================================
\documentclass[lang=cn,newtx,10pt,scheme=chinese,color=blue]{elegantbook}

\title{C 语言入门}
\subtitle{从计算机组成到指针与算法\\[1ex]\large 电控算法组新生培训教程}
\author{电控算法组}
\institute{电控算法组 $\cdot$ 新生培训系列}
\date{2026 年 10 月}
\version{1.0}

\extrainfo{先想清楚「数据长什么样」，再想清楚「步骤怎么走」。}

\cover{assets/cover.pdf}
\logo{assets/logo.pdf}

% ------------------------- 宏包 -------------------------------------------
\usepackage{tikz}
\usetikzlibrary{arrows.meta,positioning,calc,fit,backgrounds,shapes.geometric,
  shapes.misc,decorations.pathreplacing,decorations.markings,patterns,
  chains,matrix}
\usepackage{listings}
\usepackage{float}
\usepackage{array}
\usepackage{booktabs}
\usepackage{multirow}
\usepackage{makecell}
\usepackage{longtable}

% tcolorbox：代码盒、提示盒、拔高盒
\tcbuselibrary{listings,breakable,skins}

% ------------------------- 颜色 -------------------------------------------
% 与模板主题蓝色协调：代码高亮 / 内存 / 指针 / 结论 / 提醒
\definecolor{codekw}{RGB}{0,82,155}     % 关键字
\definecolor{codestr}{RGB}{163,66,20}   % 字符串
\definecolor{codecom}{RGB}{56,118,74}   % 注释
\definecolor{codetype}{RGB}{120,60,140} % 类型名
\definecolor{codebg}{RGB}{247,249,252}
\definecolor{codeframe}{RGB}{205,214,228}
\definecolor{memfill}{RGB}{232,241,252}
\definecolor{memline}{RGB}{70,110,170}
\definecolor{ptrfill}{RGB}{253,233,232}
\definecolor{ptrline}{RGB}{190,70,60}
\definecolor{okfill}{RGB}{232,246,236}
\definecolor{okline}{RGB}{52,130,82}
\definecolor{warnfill}{RGB}{255,246,228}
\definecolor{warnline}{RGB}{198,130,30}
\definecolor{deepfill}{RGB}{243,238,250}
\definecolor{deepline}{RGB}{120,90,160}
\definecolor{c89fill}{RGB}{235,244,247}
\definecolor{c89line}{RGB}{60,120,140}

% 封面分隔色带与封面图呼应
\definecolor{covercyan}{RGB}{70,200,220}
\colorlet{coverlinecolor}{covercyan}

% 代码字体：Windows 上的 Consolas，接近大家在编辑器里看到的样子
\IfFontExistsTF{Consolas}{%
  \setmonofont{Consolas}[Scale=0.92]%
}{}

% ------------------------- 代码排版 ---------------------------------------
\lstdefinestyle{c89}{
  language=C,
  basicstyle=\ttfamily\small,
  keywordstyle=\color{codekw}\bfseries,
  commentstyle=\color{codecom}\itshape,
  stringstyle=\color{codestr},
  numberstyle=\ttfamily\scriptsize\color{gray},
  numbers=left,
  numbersep=7pt,
  stepnumber=1,
  showstringspaces=false,
  numberblanklines=true,
  breaklines=true,
  breakatwhitespace=false,
  postbreak=\mbox{\textcolor{gray}{$\hookrightarrow$}\space},
  tabsize=2,
  columns=fullflexible,
  keepspaces=true,
  upquote=true,
  extendedchars=true,
  escapeinside={(*@}{@*)},
  morekeywords={size_t,ssize_t,int8_t,int16_t,int32_t,uint8_t,uint16_t,uint32_t},
}
\lstset{style=c89}

% 代码盒：\begin{codebox}{标题} ... \end{codebox}
\newtcblisting{codebox}[1]{
  enhanced, breakable, listing only,
  colback=codebg, colframe=codeframe,
  boxrule=0.5pt, arc=1.5mm,
  left=1mm, right=2mm, top=1.2mm, bottom=1.2mm,
  title={#1},
  fonttitle=\small\bfseries,
  coltitle=black,
  colbacktitle=codeframe!55,
  attach boxed title to top left={xshift=2mm,yshift=-2.2mm},
  boxed title style={arc=1mm,boxrule=0pt},
  listing options={style=c89}
}

% 运行结果盒：\begin{outbox}{标题} ... \end{outbox}
\newtcblisting{outbox}[1]{
  enhanced, breakable, listing only,
  colback=white, colframe=black!25,
  boxrule=0.4pt, arc=1mm,
  left=2mm, right=2mm, top=1mm, bottom=1mm,
  title={#1},
  fonttitle=\small\bfseries,
  coltitle=black,
  colbacktitle=black!8,
  attach boxed title to top left={xshift=2mm,yshift=-2mm},
  boxed title style={arc=0.8mm,boxrule=0pt},
  listing options={language=,basicstyle=\ttfamily\small,numbers=none,
    keywordstyle=,commentstyle=,stringstyle=,columns=fullflexible,
    showstringspaces=false,breaklines=true}
}

% ------------------------- 提示盒子 ---------------------------------------
\newtcolorbox{tipbox}[1][小贴士]{
  enhanced, breakable,
  colback=okfill, colframe=okline,
  fonttitle=\bfseries, title={#1},
  boxrule=0.5pt, arc=1.5mm, left=2mm, right=2mm, top=1.5mm, bottom=1.5mm,
  attach boxed title to top left={xshift=3mm,yshift=-2.2mm},
  boxed title style={colback=okline,arc=1mm},
  coltitle=white,
}
\newtcolorbox{pitfall}[1][小心有坑]{
  enhanced, breakable,
  colback=warnfill, colframe=warnline,
  fonttitle=\bfseries, title={#1},
  boxrule=0.5pt, arc=1.5mm, left=2mm, right=2mm, top=1.5mm, bottom=1.5mm,
  attach boxed title to top left={xshift=3mm,yshift=-2.2mm},
  boxed title style={colback=warnline,arc=1mm},
  coltitle=white,
}
\newtcolorbox{deepbox}[1][拔高阅读：可以以后再回来看]{
  enhanced, breakable,
  colback=deepfill, colframe=deepline,
  fonttitle=\bfseries, title={#1},
  boxrule=0.5pt, arc=1.5mm, left=2mm, right=2mm, top=1.5mm, bottom=1.5mm,
  attach boxed title to top left={xshift=3mm,yshift=-2.2mm},
  boxed title style={colback=deepline,arc=1mm},
  coltitle=white,
}
\newtcolorbox{c89box}[1][C89 约定]{
  enhanced, breakable,
  colback=c89fill, colframe=c89line,
  fonttitle=\bfseries, title={#1},
  boxrule=0.5pt, arc=1.5mm, left=2mm, right=2mm, top=1.5mm, bottom=1.5mm,
  attach boxed title to top left={xshift=3mm,yshift=-2.2mm},
  boxed title style={colback=c89line,arc=1mm},
  coltitle=white,
}

% ------------------------- 记号与命令 -------------------------------------
\newcommand{\code}[1]{\texttt{\detokenize{#1}}}      % 行内代码（自动处理反斜杠、下划线等）
\newcommand{\en}[1]{\textit{#1}}                   % 英文术语
\newcommand{\term}[1]{\textbf{#1}}                 % 首次出现的术语
\newcommand{\file}[1]{\texttt{#1}}                 % 文件名
\newcommand{\key}[1]{\textbf{#1}}                  % 强调
\newcommand{\cnum}[1]{{\CJKfamily{zhkai}#1}}       % 圈码/楷体小字
\newcommand{\ptr}{\textcolor{ptrline}{$\bullet$}}  % 指针记号

% 章首一句话
\newcommand{\chaptermotto}[1]{%
  \begin{flushright}
    \kaishu\small\color{black!55}#1
  \end{flushright}
  \vspace{-0.6em}
}

% TikZ 通用样式
\tikzset{
  box/.style={draw=structurecolor,fill=structurecolor!8,rounded corners=2pt,
    align=center,minimum height=8mm,inner sep=3pt,font=\small},
  cpubox/.style={draw=structurecolor!80!black,fill=structurecolor!12,
    rounded corners=2pt,align=center,minimum height=10mm,
    minimum width=26mm,font=\small\bfseries},
  membox/.style={draw=memline,fill=memfill,minimum width=12mm,
    minimum height=7mm,font=\ttfamily\small,inner sep=1pt,align=center},
  emptybox/.style={draw=memline,dashed,fill=white,minimum width=12mm,
    minimum height=7mm,font=\ttfamily\small,inner sep=1pt,align=center},
  ptrbox/.style={draw=ptrline,fill=ptrfill,minimum width=14mm,
    minimum height=7mm,font=\ttfamily\small,inner sep=1pt,align=center},
  bus/.style={draw=structurecolor!70,line width=1.2pt},
  ptr/.style={-{Stealth[length=2.2mm]},draw=ptrline,line width=0.9pt},
  flow/.style={-{Stealth[length=2mm]},draw=structurecolor!80,line width=0.7pt},
  decision/.style={diamond,draw=structurecolor!80!black,fill=structurecolor!8,
    aspect=2,align=center,inner sep=1pt,font=\small},
  process/.style={rectangle,draw=structurecolor!80!black,fill=structurecolor!8,
    rounded corners=2pt,align=center,minimum width=28mm,minimum height=8mm,
    font=\small},
  io/.style={trapezium,trapezium left angle=70,trapezium right angle=110,
    draw=structurecolor!80!black,fill=structurecolor!8,align=center,
    minimum width=30mm,inner sep=2pt,font=\small},
  terminator/.style={rounded rectangle,draw=structurecolor!80!black,
    fill=structurecolor!12,align=center,minimum width=20mm,minimum height=7mm,
    font=\small},
  lbl/.style={font=\scriptsize,color=black!60},
  arr/.style={-{Stealth[length=2mm]},draw=black!60,line width=0.6pt},
}

\hypersetup{
  pdftitle={C 语言入门：从计算机组成到指针},
  pdfauthor={电控算法组},
  pdfsubject={电控算法组新生培训教程（C89/C90）},
  pdfkeywords={C 语言, C89, 指针, 数据结构, 排序, 新生培训}
}
\setcounter{tocdepth}{1}
\setcounter{secnumdepth}{2}

\begin{document}

\maketitle
\frontmatter
\chapter*{写在前面}
\markboth{写在前面}{写在前面}
\addcontentsline{toc}{chapter}{写在前面}

\section*{这本书是写给谁的}

这本书写给刚加入电控算法组的你。

也许你之前从未写过一行代码；也许你在高中信息课上看过 \code{print("hello")}，但不知道它为什么能运行；也许你用过 Python，却对“指针”“内存”这些词心存畏惧。都没关系。这本书假设你\textbf{只会用电脑，不会编程}，从“计算机里到底发生了什么”讲起，一步一步把你带到能独立写出一个像样程序的位置。

我们不追求“三天速成”，也不打算把你训练成“语法词典”。我们想让你建立三样东西：

\begin{enumerate}
  \item \textbf{一个正确的心智模型}：知道变量住在内存的哪里、数组为什么从 0 开始、指针到底指向什么；
  \item \textbf{一套可靠的编程习惯}：先想清楚再动手、先声明再使用、写完就编译、出错会调试；
  \item \textbf{一种工程化的思维方式}：把大问题拆成小问题，把重复劳动交给循环和函数，把“能跑”变成“可读、可改、可维护”。
\end{enumerate}

\section*{为什么电控算法组要学 C}

你以后会接触电机、电池、传感器、嵌入式控制器，会写让无人机悬停、让小车循迹、让机械臂抓取的算法。这些场景有一个共同点：\textbf{程序要直接和硬件打交道，而且要在有限的时间、有限的内存里稳定运行}。

C 语言正好站在“硬件”和“算法”的交界处：

\begin{itemize}
  \item 它\textbf{贴近硬件}：一个 \code{int} 占几个字节、数组在内存里怎么排布、指针里存的是哪个地址，你都能说得清；
  \item 它\textbf{足够快}：没有隐藏的运行时开销，适合实时控制；
  \item 它\textbf{足够小}：从 8 位单片机到服务器，几乎处处都有 C 编译器；
  \item 它\textbf{是很多语言的地基}：C++、Rust、Python 解释器、操作系统内核，都能看到它的影子。
\end{itemize}

学会 C，不只是学会一门语言，更是学会“站在机器的角度思考”。以后你再学 C++、学数据结构、学嵌入式，都会轻松很多。

\section*{关于“C98”和本书的代码标准}

你可能会听到“C98”这个说法。严格来说，C 语言的标准版本叫 \textbf{C89}（1989 年 ANSI 标准，也叫 C90），之后还有 C99、C11、C17 等。本书所说的“C98 风格”，我们理解为：\textbf{按 C89/C90 的规矩写代码}。

这带来几条你必须从一开始就养成的习惯：

\begin{c89box}[本书的 C89/C90 硬性约定]
\begin{enumerate}
  \item \textbf{先声明，后使用}：所有变量都写在语句块（\texttt{\{}）的开头，声明之前不能执行任何语句。
  \item \textbf{循环变量要在外面声明}：写成
        \code{int i; for (i = 0; i < n; i++)}，
        \textbf{不能}写成 \code{for (int i = 0; ...)}——后者是 C99 才允许的。
  \item \textbf{只用 \code{/* */} 注释}，不用 \code{//}（\code{//} 是 C99 引入的）。
  \item \textbf{函数必须写返回类型}，\code{main} 结尾写 \code{return 0;}。
  \item \textbf{不使用变长数组}（\code{int a[n];} 且 \code{n} 是变量），数组长度必须是编译期常量。
\end{enumerate}
\end{c89box}

为什么这么“老派”？因为电控领域大量使用老编译器、老代码库和严格的编码规范（例如汽车电子的 MISRA C）。先把规矩学正，以后换到任何环境都不会被“惯坏”。而且，这些限制会逼你把“声明”和“逻辑”分开，代码反而更清楚。

\begin{tipbox}[如果你的编译器报错……]
如果你的环境默认使用 C99 或更新标准，用 \code{for (int i = ...)} 可能不会报错，但这不代表它符合本书约定。请主动按 C89 的写法练习。附录~\ref{app:toolchain} 给出了推荐的编译命令，其中包含 \code{-std=c89 -pedantic} 选项，可以帮你检查出不合规的写法。
\end{tipbox}

\section*{怎么读这本书}

\subsection*{第一遍：顺着读，别跳}

本书的章节顺序是精心安排的：先认识机器（第 1 章），再学“一条一条往下走”的顺序结构（第 2 章），然后才学判断（第 3、4 章）和循环（第 5 章）。数组、结构体、指针、字符串这四章是“数据怎么组织”的递进关系，排序算法则是前四章的综合运用，宏定义放在最后，讲的是“写代码的代码”。

请不要因为“这章好像很简单”就跳过。很多同学最后卡在指针，根源往往是前面数组和内存模型没学扎实。

\subsection*{第二遍：把代码敲一遍}

本书所有完整程序都可以直接编译运行。请\textbf{亲手敲}，不要复制粘贴。敲代码时你会被迫注意分号、括号、大小写，这些“低级错误”是每个程序员都要经历的。

每敲完一个程序，做三件事：

\begin{enumerate}
  \item 预测输出：先别运行，想一想它会打印什么；
  \item 运行对照：和你的预测比较，不一样就追问“为什么”；
  \item 改一改：改一个数字、删一个分号、把 \code{<} 改成 \code{<=}，看看会发生什么。
\end{enumerate}

\subsection*{第三遍：做章末习题}

每章末尾都有习题，附录~\ref{app:answers} 给出参考答案与提示。建议先自己写，卡住 20 分钟再看答案。看懂答案不等于会写，合上答案再默写一遍，才算真的会。

\section*{本书的排版约定}

为了让你在阅读时快速抓住重点，本书用了几种固定的“视觉语言”：

\begin{center}
\begin{tabular}{@{}p{0.22\textwidth}p{0.68\textwidth}@{}}
\toprule
\textbf{版式} & \textbf{含义} \\
\midrule
蓝色标题的“定义”盒 & 重要概念的正式定义，第一次出现时务必读懂 \\
绿色“小贴士” & 经验、技巧、可以让你少走弯路的建议 \\
橙色“小心有坑” & 常见错误、易混淆点，考试和调试中的高频雷区 \\
蓝灰色“C89 约定” & 与代码标准相关的硬性规定 \\
紫色“拔高阅读” & 选读内容，第一次读不懂可以先跳过，不影响后续学习 \\
带行号的浅色代码盒 & 完整或半完整的 C 程序，可直接编译 \\
“运行结果”盒 & 程序在典型环境下的输出 \\
“例题 / 解答” & 带完整分析过程的示范题 \\
“本章练习” & 章末习题，答案见附录~\ref{app:answers} \\
\bottomrule
\end{tabular}
\end{center}

正文中的行内代码用等宽字体表示，例如 \code{printf}、\code{&a}、\code{a[i]}。

\section*{学习地图}

图~\ref{fig:map} 是全书的学习路线。可以把它贴在显示器旁边：每学完一章，就在心里点亮一个格子。

\begin{figure}[htbp]
  \centering
  \begin{tikzpicture}[x=1cm,y=1cm,>=Stealth]
    % 三个阶段背景
    \fill[structurecolor!6,rounded corners=3mm] (-0.4,-0.2) rectangle (5.2,5.4);
    \fill[main!6,rounded corners=3mm] (5.4,-0.2) rectangle (11.0,5.4);
    \fill[second!6,rounded corners=3mm] (11.2,-0.2) rectangle (16.4,5.4);
    \node[font=\small\bfseries,color=structurecolor] at (2.4,5.0) {第一阶段：认识机器与基本语句};
    \node[font=\small\bfseries,color=main!80!black] at (8.2,5.0) {第二阶段：组织数据};
    \node[font=\small\bfseries,color=second!80!black] at (13.8,5.0) {第三阶段：算法与工程};

    % 节点
    \node[box,minimum width=2.6cm] (c1) at (1.2,4.0) {1 计算机组成};
    \node[box,minimum width=2.6cm] (c2) at (3.9,3.1) {2 顺序结构};
    \node[box,minimum width=2.6cm] (c3) at (2.4,2.0) {3 运算符};
    \node[box,minimum width=2.6cm] (c4) at (3.9,0.9) {4 分支};
    \node[box,minimum width=2.6cm] (c5) at (1.4,0.0) {5 循环};

    \node[box,minimum width=2.6cm] (c6) at (6.4,4.0) {6 数组};
    \node[box,minimum width=2.6cm] (c7) at (9.6,3.4) {7 结构体};
    \node[box,minimum width=2.6cm] (c8) at (7.0,2.4) {8 指针};
    \node[box,minimum width=2.6cm] (c9) at (9.8,1.4) {9 字符串};

    \node[box,minimum width=2.6cm] (c10) at (12.8,2.6) {10 排序算法};
    \node[box,minimum width=2.6cm] (c11) at (14.6,1.2) {11 宏定义};

    % 连接
    \draw[flow] (c1) -- (c2);
    \draw[flow] (c2) -- (c3);
    \draw[flow] (c3) -- (c4);
    \draw[flow] (c4) -- (c5);
    \draw[flow] (c5.east) to[bend left=12] (c6.west);
    \draw[flow] (c6) -- (c7);
    \draw[flow] (c7) -- (c8);
    \draw[flow] (c8) -- (c9);
    \draw[flow] (c9) -- (c10);
    \draw[flow] (c10) -- (c11);

    % 注释
    \node[lbl,align=center] at (3.0,2.75) {先会走路};
    \node[lbl,align=center] at (8.5,0.7) {数据是一切的载体};
    \node[lbl,align=center] at (13.4,0.2) {把学过的用起来};
  \end{tikzpicture}
  \caption{全书学习路线：先建立机器模型，再学控制流程，然后组织数据，最后进入算法与工程}
  \label{fig:map}
\end{figure}

\begin{tipbox}[三个“不要”]
不要只看不敲；不要一遇到报错就问人（先自己读错误信息）；不要背语法，要理解“这条语句让内存发生了什么”。
\end{tipbox}

\section*{开发环境}

本书不绑定任何编辑器。你只需要一个能编译 C89 的编译器（推荐 GCC/MinGW），以及一个顺手的编辑器（VS Code、Notepad++、Vim 都可以）。安装步骤、编译命令和常见报错见附录~\ref{app:toolchain}。

第一个程序永远是：

\begin{codebox}{hello.c —— 一切的开始}
#include <stdio.h>

int main(void)
{
    printf("Hello, world!\n");
    return 0;
}
\end{codebox}

\begin{outbox}{运行结果}
Hello, world!
\end{outbox}

如果它能跑起来，恭喜你，你已经迈出了第一步。接下来，我们先把“计算机到底是怎么执行这行字的”弄清楚——这会让你后面少走很多弯路。

\vspace{1em}
\begin{flushright}
\kaishu 电控算法组\\
2026 年 10 月
\end{flushright}


\tableofcontents

\mainmatter

\chapter{计算机组成基础}

\begin{introduction}
  \item 计算机由哪几个部分组成
  \item 程序到底是什么
  \item 计算机为什么只认识 0 和 1
  \item 内存、地址和变量
  \item 从写代码到程序运行
\end{introduction}

\chaptermotto{「计算机很笨，但它很听话。」}

同学们，先别急着写代码。我们先用一节课的时间，认识一下计算机这个“大玩具”。

你可能会觉得计算机很聪明：它会算题、会放视频、会下棋。其实它一点都不聪明。它只会做很简单的事情：\textbf{算一算、比一比、把东西搬来搬去}。它之所以看起来聪明，是因为人把做事情的步骤一条一条写给了它。

这一章不写程序，只讲五件事：计算机由什么组成、程序是什么、为什么它只认识 0 和 1、内存是什么、以及你写的代码是怎么变成能运行的程序的。

\section{计算机能做什么}

我们用计算机做事情，通常都是这样一个过程：

\begin{center}
\begin{tikzpicture}[x=1cm,y=1cm,>=Stealth,node distance=12mm]
  \node[box] (in) {输入};
  \node[box,right=of in] (do) {处理};
  \node[box,right=of do] (out) {输出};
  \draw[flow] (in) -- (do);
  \draw[flow] (do) -- (out);
\end{tikzpicture}
\end{center}

举几个例子：

\begin{itemize}
  \item 用计算器算 $3+5$：你按的“3、+、5”是\textbf{输入}；计算器在内部做加法是\textbf{处理}；屏幕上出现“8”是\textbf{输出}。
  \item 用手机拍照：镜头看到画面是输入；手机把画面变成照片是处理；屏幕上显示照片是输出。
  \item 让小车前进：传感器测到距离是输入；单片机判断“前面没有障碍”是处理；电机转动是输出。
\end{itemize}

\textbf{输入 → 处理 → 输出}，这就是计算机做事的三个步骤。我们写程序，主要就是在告诉计算机“怎么处理”。

\section{计算机的四样东西}

一台计算机看起来零件很多，但真正重要的只有四样：

\begin{figure}[htbp]
  \centering
  \begin{tikzpicture}[x=1cm,y=1cm,>=Stealth]
    \node[cpubox,minimum width=3cm,minimum height=1.4cm] (cpu) at (0,1.8) {CPU\\（大脑）};
    \node[box,minimum width=2.8cm] (mem) at (-4,-0.8) {内存\\（草稿纸）};
    \node[box,minimum width=2.8cm] (disk) at (4,-0.8) {硬盘\\（书包）};
    \node[box,minimum width=2.8cm] (io) at (0,-2.8) {输入 / 输出\\（眼睛、嘴巴）};

    \draw[<->,>=Stealth,draw=structurecolor,line width=0.8pt] (cpu.west) -| (mem.north);
    \draw[<->,>=Stealth,draw=structurecolor,line width=0.8pt] (cpu.east) -| (disk.north);
    \draw[<->,>=Stealth,draw=structurecolor,line width=0.8pt] (cpu.south) -- (io.north);
    \draw[<->,>=Stealth,draw=structurecolor,line width=0.8pt] (mem.east) -- (io.west);
    \draw[<->,>=Stealth,draw=structurecolor,line width=0.8pt] (disk.west) -- (io.east);
  \end{tikzpicture}
  \caption{计算机的四样核心部件：CPU、内存、硬盘、输入输出}
  \label{fig:four-parts}
\end{figure}

\subsection{CPU：计算机的大脑}

CPU 是计算机的“大脑”，全名叫\textbf{中央处理器}。

它负责两件事：

\begin{enumerate}
  \item \textbf{算}：做加减乘除、比较大小、判断真假；
  \item \textbf{指挥}：告诉其他部件“现在该做什么”。
\end{enumerate}

CPU 算得很快，一秒钟能做几十亿次运算。但它只会做很简单的运算，复杂的事情要拆成很多小步骤。

\subsection{内存：计算机的草稿纸}

内存是计算机的“草稿纸”，用来临时存放正在使用的数据。

它的特点是：

\begin{itemize}
  \item 速度比较快；
  \item \textbf{一断电，里面的东西就全没了}。
\end{itemize}

所以内存只适合放“正在用”的东西。比如你正在写的文档、正在玩的游戏，数据都在内存里。突然断电，没保存的内容就丢了。

\subsection{硬盘：计算机的书包}

硬盘是计算机的“书包”，用来长期存放东西。

它的特点是：

\begin{itemize}
  \item 速度比内存慢；
  \item \textbf{断电以后东西还在}。
\end{itemize}

你的照片、作业、程序文件，平时都放在硬盘里。用的时候，计算机会把它们从硬盘搬到内存里。

\begin{tipbox}[手机也是一样的]
手机的“8 GB + 256 GB”里，8 GB 通常指运行内存（相当于电脑的内存），256 GB 是存储空间（相当于硬盘）。内存越大，能同时开的程序越多；存储越大，能存的东西越多。
\end{tipbox}

\subsection{输入和输出：眼睛、耳朵和嘴巴}

输入设备让计算机“知道外面发生了什么”，例如键盘、鼠标、麦克风、摄像头、温度传感器。

输出设备让计算机“把结果告诉外面”，例如屏幕、音箱、打印机、电机、指示灯。

\begin{center}
\begin{tabular}{@{}lll@{}}
\toprule
部件 & 像人的什么 & 作用 \\
\midrule
CPU & 大脑 & 计算、指挥 \\
内存 & 草稿纸 & 临时存放正在用的数据 \\
硬盘 & 书包 & 长期存放数据 \\
输入设备 & 眼睛、耳朵 & 把信息送进计算机 \\
输出设备 & 嘴巴、手 & 把结果送出去 \\
\bottomrule
\end{tabular}
\end{center}

\section{程序是什么}

\subsection{程序就是一张步骤清单}

程序不是魔法，它只是一张\textbf{步骤清单}。就像泡面的做法：

\begin{enumerate}
  \item 烧一壶水；
  \item 把面饼放进碗里；
  \item 倒入开水；
  \item 等三分钟；
  \item 放入调料，搅拌。
\end{enumerate}

计算机做事情也是这样：你给它一张步骤清单，它就从第一条开始，一条一条做下去。

\subsection{计算机不会猜}

人听别人说话，可以猜意思。计算机不会猜。你写“把水烧开”，它不知道多少度算开、用什么烧、要等多久。你必须写清楚：

\begin{itemize}
  \item 用哪个锅？
  \item 加热到多少度？
  \item 到了温度以后做什么？
\end{itemize}

写程序也是一样：\textbf{每一步都要明确，不能有歧义}。这就是为什么程序设计要学很多规则。

\begin{tipbox}[把计算机当成一个非常听话、但非常死板的同学]
你说“往前走”，他就一直往前走，撞墙也不停，因为你没说“遇到墙就停”。写程序时，要把所有情况都想到。
\end{tipbox}

\section{计算机为什么只认识 0 和 1}

\subsection{电灯只有开和关}

计算机里面是电子电路。电路最简单的两种状态是：

\begin{itemize}
  \item 通电；
  \item 断电。
\end{itemize}

就像电灯：要么亮，要么灭。我们用数字表示这两种状态：

\[
\text{通电} = 1, \qquad \text{断电} = 0.
\]

只用两种状态，电路最简单、最不容易出错。所以计算机选择用 0 和 1 来表示所有信息。

\subsection{位和字节}

\begin{itemize}
  \item 一个 0 或一个 1，叫一位（bit，比特）；
  \item 8 位合在一起，叫一个字节（byte）。
\end{itemize}

\begin{center}
\begin{tabular}{@{}cccccccc@{}}
\toprule
1 & 0 & 1 & 1 & 0 & 0 & 1 & 0 \\
\bottomrule
\end{tabular}
\end{center}

这是一个字节，一共 8 位。一个字节可以组合出 $2^8 = 256$ 种不同的样子（从 \texttt{00000000} 到 \texttt{11111111}）。

\subsection{用 0 和 1 数数}

我们平时用十进制：0、1、2、3、4、5、6、7、8、9，到 9 再往上就进位，变成 10。

计算机用二进制：只有 0 和 1，到 1 再往上就进位：

\begin{center}
\begin{tabular}{@{}ll@{}}
\toprule
二进制 & 十进制 \\
\midrule
\texttt{0} & 0 \\
\texttt{1} & 1 \\
\texttt{10} & 2 \\
\texttt{11} & 3 \\
\texttt{100} & 4 \\
\texttt{101} & 5 \\
\texttt{110} & 6 \\
\texttt{111} & 7 \\
\texttt{1000} & 8 \\
\bottomrule
\end{tabular}
\end{center}

规律是：\textbf{满二进一}。和我们数数“满十进一”是一样的，只是计算机只有两个数字。

\begin{example}
二进制 \texttt{1011} 表示十进制多少？
\end{example}
\begin{solution}
从右往左，每一位的“份量”分别是 1、2、4、8：

\[
\texttt{1011} = 8 + 0 + 2 + 1 = 11.
\]

所以二进制 \texttt{1011} 就是十进制的 11。
\end{solution}

\subsection{字符怎么表示}

计算机里没有“字母”，只有数字。那怎么存字母呢？

人们约定了一张“密码表”，给每个字符编一个数字。最常用的是 \textbf{ASCII 码}。例如：

\begin{center}
\begin{tabular}{@{}lll@{}}
\toprule
字符 & ASCII 码 & 说明 \\
\midrule
\texttt{'0'} & 48 & 数字字符 0 \\
\texttt{'A'} & 65 & 大写字母 A \\
\texttt{'a'} & 97 & 小写字母 a \\
空格 & 32 & 空格 \\
\texttt{'\textbackslash n'} & 10 & 换行 \\
\bottomrule
\end{tabular}
\end{center}

所以计算机里存的不是字母 \texttt{A}，而是数字 65。程序输出时，再按照密码表把 65 变回字母 \texttt{A}。

汉字也有编码表，例如 GBK、UTF-8。一个汉字通常要占 2 个或 3 个字节。这部分以后用到再了解。

\begin{tipbox}[负数和小数怎么办？]
负数和小数在计算机里也有专门的表示方法，但比正数复杂。初学阶段先不用管，只要知道“计算机有办法表示”就可以了。等你把基本程序写熟，再回头看也不迟。
\end{tipbox}

\section{内存像一排小盒子}

\subsection{地址就是盒子编号}

内存可以想象成一排小盒子，每个盒子装 1 个字节。每个盒子都有一个编号，这个编号叫\textbf{地址}。

\begin{figure}[htbp]
  \centering
  \begin{tikzpicture}[x=1cm,y=1cm]
    \foreach \i/\addr/\v in {0/1000/?,1/1001/?,2/1002/?,3/1003/?,4/1004/?} {
      \node[font=\ttfamily\scriptsize,color=black!55,anchor=east] at (-0.2,-\i*0.9) {\addr};
      \node[membox,minimum width=2.2cm] at (1.6,-\i*0.9) {\v};
    }
    \draw[decorate,decoration={brace,amplitude=4pt},structurecolor!70]
      (3.0,0.3) -- (3.0,-2.4);
    \node[anchor=west,font=\small] at (3.2,-1.05) {从 1000 号盒子开始};
  \end{tikzpicture}
  \caption{内存像一排有编号的小盒子，每个盒子有地址}
  \label{fig:memory-boxes}
\end{figure}

地址就像酒店的房间号：房间号不是房间本身，但拿着房间号就能找到房间。

\subsection{变量就是贴了标签的盒子}

写程序时，我们不需要每次都记“1000”这样的地址。我们给盒子贴一个标签，起一个名字，例如：

\begin{codebox}{variable-box.c —— 变量就是一个盒子}
int a = 10;
\end{codebox}

这句话的意思是：

\begin{itemize}
  \item 准备一个能装整数的盒子；
  \item 在盒子外面贴上标签 \texttt{a}；
  \item 往盒子里放 10。
\end{itemize}

以后写 \texttt{a}，就是在说“那个贴了 \texttt{a} 标签的盒子”。盒子里的东西可以换，但标签不变。

\begin{tipbox}[变量三件事]
每个变量都有：\textbf{名字}（标签）、\textbf{类型}（盒子能装什么）、\textbf{值}（盒子里现在装什么）。地址是盒子的编号，通常由计算机管理，不需要你手动记。
\end{tipbox}

\section{从写代码到程序运行}

\subsection{人话和机器话}

你写的 C 代码，长这样：

\begin{codebox}{hello.c —— 给人看的代码}
#include <stdio.h>

int main(void)
{
    printf("Hello!\n");
    return 0;
}
\end{codebox}

但 CPU 看不懂英文，它只认识 0 和 1 组成的机器指令。所以中间需要一个“翻译官”。

\subsection{编译器就是翻译官}

把 C 代码翻译成机器能懂的程序，这个工具叫\textbf{编译器}。

\begin{figure}[htbp]
  \centering
  \begin{tikzpicture}[x=1cm,y=1cm,>=Stealth,node distance=13mm]
    \node[box,minimum width=2.8cm] (src) {你写的代码\\\texttt{hello.c}};
    \node[box,minimum width=2.8cm,right=of src] (cc) {编译器\\（翻译官）};
    \node[box,minimum width=2.8cm,right=of cc] (exe) {可执行程序\\\texttt{hello.exe}};
    \node[box,minimum width=2.8cm,right=of exe] (run) {双击运行\\屏幕上出字};
    \draw[flow] (src) -- (cc);
    \draw[flow] (cc) -- (exe);
    \draw[flow] (exe) -- (run);
  \end{tikzpicture}
  \caption{写代码 → 编译 → 运行}
  \label{fig:compile-simple}
\end{figure}

过程是这样的：

\begin{enumerate}
  \item 你在电脑上写一个 \texttt{hello.c} 文件；
  \item 用编译器把它翻译成 \texttt{hello.exe}；
  \item 双击 \texttt{hello.exe}，计算机就按里面的指令做事。
\end{enumerate}

\subsection{两种错误}

写程序会遇到两种错误：

\begin{itemize}
  \item \textbf{编译错误}：翻译官看不懂你写的话。比如少了一个分号、括号不配对。这时程序根本生成不了。
  \item \textbf{运行错误}：程序能生成，但运行的时候做错了事。比如让计算机除以 0，或者读到了不该读的地方。
\end{itemize}

\begin{tipbox}[报错不可怕]
编译器报错时，会告诉你“第几行出了问题”。先看第一行错误，通常改掉它，后面的错误就消失了。
\end{tipbox}

\section{C 语言为什么适合电控}

你以后要接触电机、传感器、单片机、机器人。这些设备有三个特点：

\begin{itemize}
  \item 要\textbf{快}：电机控制一秒钟要算很多次；
  \item 要\textbf{省}：单片机的内存很小；
  \item 要\textbf{直接控制硬件}：读写寄存器、引脚、串口。
\end{itemize}

C 语言正好适合：

\begin{enumerate}
  \item 它很\textbf{接近硬件}，能直接操作内存和寄存器；
  \item 它\textbf{很快}，没有太多额外的开销；
  \item 它\textbf{通用}，从单片机到电脑都能用；
  \item 学好 C，再学 C++、Python、数据结构都会更容易。
\end{enumerate}

\section{本章小结}

\begin{center}
\begin{tabular}{@{}p{0.28\textwidth}p{0.6\textwidth}@{}}
\toprule
问题 & 一句话答案 \\
\midrule
计算机做什么 & 输入 → 处理 → 输出 \\
有哪四样核心部件 & CPU、内存、硬盘、输入输出 \\
CPU 是什么 & 大脑，负责计算和指挥 \\
内存和硬盘的区别 & 内存快但断电就丢；硬盘慢但能长期保存 \\
程序是什么 & 给计算机的步骤清单，必须一步一步写清楚 \\
计算机为什么用二进制 & 电路只有通/断两种稳定状态 \\
位和字节 & 1 位是 0 或 1；8 位是 1 字节 \\
内存是什么 & 一排有编号的小盒子，编号叫地址 \\
变量是什么 & 贴了名字标签的盒子 \\
代码怎么运行 & 写代码 → 编译器翻译 → 可执行程序 → 运行 \\
\bottomrule
\end{tabular}
\end{center}

\begin{tipbox}[下一章预告]
认识完计算机，下一章我们就写第一个真正的 C 程序：让计算机在屏幕上说话，再让它听我们输入数字、做计算。
\end{tipbox}

\begin{problemset}
  \item 用自己的话说一说：计算机做一件事要经过哪三个步骤？
  \item CPU、内存、硬盘各像什么？它们分别负责什么？
  \item 内存和硬盘最大的区别是什么？为什么手机内存越大越贵？
  \item 程序为什么必须一步一步写清楚？举一个“没说清楚”的例子。
  \item 二进制 \texttt{1010}、\texttt{1100}、\texttt{11111} 分别等于十进制多少？
  \item 1 个字节等于多少位？一个字节最多能表示多少种不同的组合？
  \item 字符 \texttt{'A'} 在计算机里存的是什么？为什么不是直接存字母？
  \item 把“写代码 → 运行”的过程画出来，并说明编译器做了什么。
  \item 编译错误和运行错误有什么区别？各举一个例子。
  \item 思考题：为什么计算机不直接用十进制，而是用二进制？
\end{problemset}

\chapter{正式上路：顺序结构程序设计}

\begin{introduction}
  \item 从第一个程序开始，逐行读懂 C 程序的骨架
  \item 变量：名字、类型、值、地址四件事
  \item 基本数据类型与字面量的写法
  \item \texttt{printf} 与 \texttt{scanf}：程序如何说话、如何倾听
  \item 表达式、语句、注释与 C89 的声明规则
  \item 顺序结构的综合例题与常见错误
\end{introduction}

\chaptermotto{「所有伟大的程序，都是从一行 \texttt{printf} 开始的。」}

上一章我们把舞台搭好了：CPU 执行指令，数据存在内存里，变量有地址，整数用补码、小数用浮点。现在，是时候让程序真正跑起来了。

这一章只讲一件事：\textbf{顺序结构}——程序从上到下，一条语句接一条语句地执行。没有判断，没有循环，就像照着菜谱一步步做菜。听起来简单，但它是后面所有结构的基础。请把这一章的每个例子都亲手敲一遍。

\section{第一个程序：Hello, world!}

\subsection{完整代码}

\begin{codebox}{hello.c —— 第一个 C 程序}
#include <stdio.h>

int main(void)
{
    printf("Hello, world!\n");
    return 0;
}
\end{codebox}

\begin{outbox}{运行结果}
Hello, world!
\end{outbox}

短短 5 行代码，包含了 C 程序的所有基本骨架。我们一行一行拆开看。

\subsection{逐行解剖}

\subsubsection*{第 1 行：\texttt{\#include <stdio.h>}}

这是一条\textbf{预处理指令}。它以 \code{#} 开头，在真正的编译开始之前，由预处理器处理。

\begin{itemize}
  \item \code{stdio} 是 \en{standard input/output}（标准输入输出）的缩写；
  \item \code{.h} 表示头文件（header）；
  \item 这一行的意思是：把 \file{stdio.h} 这个文件的内容“复制粘贴”到当前位置。
\end{itemize}

为什么要包含它？因为 \code{printf} 这个函数的“说明书”就写在 \file{stdio.h} 里。没有这条指令，编译器就不认识 \code{printf}，会报“隐式声明”的警告或错误。

\begin{tipbox}[头文件不是库]
\code{#include <stdio.h>} 只是把\textbf{声明}拿进来，告诉你“有这么个函数，参数长这样”。函数真正的实现（机器指令）在链接阶段才从 C 标准库里接进来。这就是第 1 章讲的“编译”和“链接”的分工。
\end{tipbox}

\subsubsection*{第 3 行：\texttt{int main(void)}}

这是程序的\textbf{入口}。操作系统加载你的程序后，会先找到 \code{main} 函数，从它的第一行开始执行。

\begin{itemize}
  \item \code{int}：这个函数执行完后，会返回一个整数给操作系统；
  \item \code{main}：函数名，固定写法，不能改；
  \item \code{void}：表示这个函数不接收参数。也可以写成 \code{int main()}，但 \code{void} 更明确；
  \item 括号和后面的花括号：函数体，也就是这个函数要执行的语句。
\end{itemize}

\begin{pitfall}[\texttt{void main()} 是错的]
很多老教材写 \code{void main()}，这在标准 C 里是不规范的。请始终写 \code{int main(void)}，并在结尾写 \code{return 0;}。\code{return 0;} 告诉操作系统“程序正常结束”。
\end{pitfall}

\subsubsection*{第 5 行：\texttt{printf("Hello, world!\textbackslash n");}}

这是一条\textbf{函数调用语句}，拆成四部分：

\begin{center}
\begin{tikzpicture}[x=1cm,y=1cm,>=Stealth]
  \node[font=\ttfamily\small] (p) at (0,0) {printf};
  \node[font=\ttfamily\small] (lp) at (1.5,0) {(};
  \node[font=\ttfamily\small,draw=structurecolor!60,fill=structurecolor!6,rounded corners=1mm,inner sep=2pt] (s) at (5.2,0) {"Hello, world!\textbackslash n"};
  \node[font=\ttfamily\small] (rp) at (8.7,0) {);};
  \node[lbl,above=3mm of p] {函数名};
  \node[lbl,above=3mm of s] {格式串（字符串字面量）};
  \node[lbl,above=3mm of rp] {语句结束};
\end{tikzpicture}
\end{center}

\begin{itemize}
  \item \code{printf} 是标准库函数，负责“格式化输出”；
  \item 双引号括起来的部分叫\textbf{字符串字面量}；
  \item \code{\n} 是转义字符，表示“换行”，不是两个普通字符；
  \item 最后的分号 \code{;} 表示一条语句结束。\textbf{C 语言里，语句必须有分号}。
\end{itemize}

\subsubsection*{第 6 行：\texttt{return 0;}}

把整数 0 返回给操作系统。0 通常表示“成功”，非 0 表示“出错”。在 \code{main} 里，\code{return 0;} 是程序正常退出的标志。

\subsection{程序的结构}

把上面的代码抽象一下，一个最简单的 C 程序长这样：

\begin{center}
\begin{tikzpicture}[x=1cm,y=1cm,>=Stealth,node distance=4mm]
  \node[box,minimum width=9cm,align=left,font=\ttfamily\small] (inc) {\#include <stdio.h>};
  \node[box,minimum width=9cm,align=left,font=\ttfamily\small,below=of inc] (main) {int main(void)};
  \node[box,minimum width=9cm,align=left,font=\ttfamily\small,below=of main] (body) {\{ \quad 语句 1；语句 2；…… \quad \}};
  \node[lbl,right=3mm of inc] {预处理指令};
  \node[lbl,right=3mm of main] {函数头};
  \node[lbl,right=3mm of body] {函数体};
\end{tikzpicture}
\end{center}

\begin{c89box}
从第一个程序开始就养成习惯：
\begin{enumerate}
  \item \code{main} 写成 \code{int main(void)}；
  \item 左花括号 \texttt{\{} 单独占一行（也可以跟在函数头后面，但本书统一单独一行）；
  \item 花括号内的语句统一缩进 4 个空格；
  \item 每条语句以分号结尾；
  \item 最后一个 \code{return 0;} 不要忘。
\end{enumerate}
\end{c89box}

\section{变量：给内存里的格子起个名字}

\subsection{变量是什么}

第 1 章说过：变量是内存中一块有名字、有类型、有地址的存储空间。用一张图表示：

\begin{figure}[htbp]
  \centering
  \begin{tikzpicture}[x=1cm,y=1cm]
    \node[draw=structurecolor!70!black,fill=structurecolor!7,rounded corners=2mm,
          minimum width=8.2cm,minimum height=3.4cm] (box) at (0,0) {};
    \node[anchor=west,font=\small] at (-3.8,1.1) {\textbf{名字}：\texttt{a}};
    \node[anchor=west,font=\small] at (-3.8,0.35) {\textbf{类型}：\texttt{int}（4 字节）};
    \node[anchor=west,font=\small] at (-3.8,-0.4) {\textbf{值}：\texttt{10}};
    \node[anchor=west,font=\small] at (-3.8,-1.15) {\textbf{地址}：\texttt{0x1000}};
    \node[font=\Large\ttfamily,color=structurecolor!80!black] at (2.9,0.35) {10};
    \node[font=\scriptsize,color=black!55] at (2.9,-0.45) {占用 4 个字节};
  \end{tikzpicture}
  \caption{一个 \texttt{int} 变量的四个要素：名字、类型、值、地址}
  \label{fig:variable}
\end{figure}

这四件事里，\textbf{名字}和\textbf{类型}是给程序员和编译器看的；\textbf{值}是运行时真正存的东西；\textbf{地址}是内存中的位置。编译之后，变量名通常就不存在了，机器只认地址。

\begin{tipbox}[把变量想成贴了标签的储物格]
储物格有编号（地址）、有大小（类型）、里面放着东西（值），外面贴着标签（名字）。你写 \code{a = 10;}，就是“把 10 放进标签为 \code{a} 的格子里”。
\end{tipbox}

\subsection{声明、赋值、初始化}

\begin{codebox}{declare.c —— 声明与赋值}
#include <stdio.h>

int main(void)
{
    int a;          /* 声明：告诉编译器“有这么个 int 变量” */
    int b = 20;     /* 初始化：声明的同时给初值 */
    int sum;        /* 只声明，暂时没有有意义的值 */

    a = 10;         /* 赋值：把 10 写进 a 的内存 */
    sum = a + b;    /* 用 a 和 b 的值算出结果，写进 sum */

    printf("a = %d, b = %d, sum = %d\n", a, b, sum);
    return 0;
}
\end{codebox}

\begin{outbox}{运行结果}
a = 10, b = 20, sum = 30
\end{outbox}

\begin{itemize}
  \item \textbf{声明}（declaration）：\code{int a;} 告诉编译器“我要一个 \code{int} 类型的变量，名字叫 \code{a}”。
  \item \textbf{赋值}（assignment）：\code{a = 10;} 把右边表达式的值写进左边变量的内存。
  \item \textbf{初始化}（initialization）：\code{int b = 20;} 在声明的同时给初值。
\end{itemize}

\begin{pitfall}[未初始化的局部变量是“垃圾值”]
\code{int sum;} 之后如果没有赋值就直接打印它，输出的是一个不可预测的值。局部变量不会自动清零！这是初学者最常见的 bug 之一。\textbf{变量在使用前一定要先赋值}。
\end{pitfall}

\subsection{C89 的声明规则}

这是本书反复强调的重点。在 C89 里，\textbf{一个语句块内的所有声明必须放在块的开头}，在任何执行语句之前。

\begin{codebox}{order.c —— 正确的声明顺序（C89）}
#include <stdio.h>

int main(void)
{
    int a = 3;
    int b = 4;
    int sum;

    sum = a + b;
    printf("%d\n", sum);
    return 0;
}
\end{codebox}

\begin{codebox}{wrong.c —— 错误示范（C99 才允许）}
#include <stdio.h>

int main(void)
{
    int a = 3;

    printf("a = %d\n", a);

    int b = 4;      /* 错误：声明出现在了执行语句之后 */
    printf("b = %d\n", b);
    return 0;
}
\end{codebox}

在 C89 编译器下，\file{wrong.c} 会报类似这样的错误：

\begin{outbox}{编译错误（GCC，-std=c89）}
wrong.c: In function 'main':
wrong.c:9:5: error: ISO C90 forbids mixed declarations and code
\end{outbox}

\begin{c89box}[正确的写法]
把 \code{int b = 4;} 挪到 \code{printf} 之前、所有执行语句之前：
\begin{codebox}{right.c —— 声明全部提前}
#include <stdio.h>

int main(void)
{
    int a = 3;
    int b = 4;

    printf("a = %d\n", a);
    printf("b = %d\n", b);
    return 0;
}
\end{codebox}
\end{c89box}

\subsection{命名规则与风格}

变量名不是随便起的，必须遵守语法：

\begin{itemize}
  \item 只能包含字母、数字、下划线，且不能以数字开头；
  \item 区分大小写：\code{sum} 和 \code{Sum} 是两个不同的变量；
  \item 不能使用关键字（如 \code{int}、\code{return}、\code{for}）；
  \item 不能和标准库函数重名（例如不要定义叫 \code{printf} 的变量）。
\end{itemize}

风格上，本书推荐：

\begin{center}
\begin{tabular}{@{}lll@{}}
\toprule
场景 & 推荐写法 & 说明 \\
\midrule
普通变量 & \code{sum}、\code{max_value} & 全小写，单词间用下划线 \\
常量 & \code{MAX_SIZE}、\code{PI} & 全大写，突出“不打算改” \\
结构体类型 & \code{Student}、\code{MotorState} & 首字母大写 \\
函数 & \code{compute_sum}、\code{read_sensor} & 小写 + 动词开头 \\
\bottomrule
\end{tabular}
\end{center}

\begin{pitfall}[这些名字不要用]
\code{int}、\code{float}、\code{for}、\code{while}、\code{return} 是关键字；\code{printf}、\code{scanf}、\code{main} 是标准库里的名字。用它们当变量名，轻则编译报错，重则行为诡异。
\end{pitfall}

\section{基本数据类型}

C 语言的基本类型不多，但每一个都有自己的“大小”和“取值范围”。

\subsection{整数类型}

\begin{center}
\begin{tabular}{@{}llll@{}}
\toprule
类型 & 典型大小 & 典型范围 & 用途 \\
\midrule
\code{char} & 1 字节 & $-128 \sim 127$ & 字符、很小的整数 \\
\code{unsigned char} & 1 字节 & $0 \sim 255$ & 字节数据、位操作 \\
\code{short} & 2 字节 & $-32768 \sim 32767$ & 节省内存的小整数 \\
\code{int} & 4 字节 & 约 $\pm 21$ 亿 & 最常用的整数 \\
\code{unsigned int} & 4 字节 & $0 \sim 4294967295$ & 计数、位掩码 \\
\code{long} & 4 或 8 字节 & 平台相关 & 需要更大范围时 \\
\code{long long} & 8 字节 & 约 $\pm 9.2 \times 10^{18}$ & C99 引入，C89 不一定支持 \\
\bottomrule
\end{tabular}
\end{center}

\begin{note}
“典型大小”只是最常见的情况。C 标准只规定了 \code{char} 一定是 1 字节、\code{short} 不小于 \code{char}、\code{int} 不小于 \code{short} 等相对关系，具体大小由编译器和平台决定。写程序时如果需要精确大小，用 \code{sizeof} 或者 C99 的 \code{int32_t}（C89 下需要自己 typedef）。
\end{note}

\subsection{浮点类型}

\begin{center}
\begin{tabular}{@{}llll@{}}
\toprule
类型 & 典型大小 & 有效数字 & 用途 \\
\midrule
\code{float} & 4 字节 & 约 6\textasciitilde7 位 & 精度要求不高的计算 \\
\code{double} & 8 字节 & 约 15\textasciitilde16 位 & 默认选择，精度更高 \\
\code{long double} & 8/16 字节 & 平台相关 & 特殊高精度场景 \\
\bottomrule
\end{tabular}
\end{center}

\begin{tipbox}[新手先用 \texttt{double}]
除非你明确知道自己在省内存（比如单片机），否则浮点数一律用 \code{double}。\code{float} 的精度比你想的低，累加几次就可能出现明显误差。
\end{tipbox}

\subsection{字符类型}

\code{char} 存的是字符的编码（本质是一个小整数）。例如字符 \code{'A'} 的 ASCII 码是 65：

\begin{codebox}{char.c —— 字符与编码}
#include <stdio.h>

int main(void)
{
    char c = 'A';

    printf("c = %c\n", c);
    printf("c 的编码 = %d\n", c);

    c = c + 1;              /* 字符也可以参与运算 */
    printf("c 加 1 后 = %c\n", c);

    return 0;
}
\end{codebox}

\begin{outbox}{运行结果}
c = A
c 的编码 = 65
c 加 1 后 = B
\end{outbox}

注意：\code{'A'} 是字符常量，用单引号；\code{"A"} 是字符串，用双引号，两者完全不同。字符串的细节第 9 章再讲。

\subsection{字面量怎么写}

字面量就是直接写在代码里的值。

\begin{center}
\begin{tabular}{@{}lll@{}}
\toprule
写法 & 例子 & 说明 \\
\midrule
十进制整数 & \code{42}、\code{-7} & 最常用 \\
八进制整数 & \code{052} & 前导 0，等于十进制 42（不推荐） \\
十六进制整数 & \code{0x2A} & 前缀 \code{0x}，等于十进制 42 \\
无符号后缀 & \code{42U} & 表示 \code{unsigned} \\
长整型后缀 & \code{42L} & 表示 \code{long} \\
浮点数 & \code{3.14}、\code{0.5} & 带小数点 \\
科学计数法 & \code{1.5e3} & 等于 1500.0 \\
字符 & \code{'A'}、\code{'0'}、\code{\n} & 单引号，一个字符 \\
字符串 & \code{"hello"}、\code{""} & 双引号，零个或多个字符 \\
\bottomrule
\end{tabular}
\end{center}

\subsection{转义字符}

有些字符没法直接写出来（比如换行、制表符），或者有特殊含义（比如双引号），就用反斜杠加字母表示：

\begin{center}
\begin{tabular}{@{}lll@{}}
\toprule
转义 & 含义 & ASCII 码 \\
\midrule
\code{\n} & 换行（newline） & 10 \\
\code{\t} & 水平制表符（tab） & 9 \\
\code{\r} & 回车（carriage return） & 13 \\
\code{\b} & 退格 & 8 \\
\code{\\} & 一个反斜杠 & 92 \\
\code{\"} & 双引号 & 34 \\
\code{\'} & 单引号 & 39 \\
\code{\0} & 空字符（字符串结束标志） & 0 \\
\bottomrule
\end{tabular}
\end{center}

\begin{example}
下面这条语句会输出什么？
\begin{codebox}{escape.c —— 转义字符}
#include <stdio.h>

int main(void)
{
    printf("姓名\t年龄\n");
    printf("张三\t18\n");
    printf("他说：\"你好！\"\n");
    return 0;
}
\end{codebox}
\end{example}
\begin{solution}
\code{\t} 输出制表符，对齐两列；\code{\n} 换行；\code{\"} 输出普通的双引号。结果：
\begin{outbox}{运行结果}
姓名    年龄
张三    18
他说："你好！"
\end{outbox}
\end{solution}

\section{输出与输入}

程序要和外界交流：把结果告诉人（输出），把人给的数据读进来（输入）。C 语言的标准库提供了一组函数，最常用的是 \code{printf} 和 \code{scanf}。

\subsection{\texttt{printf}：格式化输出}

\code{printf} 的名字来自 \en{print formatted}（格式化打印）。它的第一个参数是\textbf{格式串}，里面可以嵌入\textbf{占位符}，后面的参数依次填进去。

\begin{center}
\begin{tikzpicture}[x=1cm,y=1cm,>=Stealth]
  \node[font=\ttfamily\small,anchor=west] (f) at (0,0) {printf("a = \%d, b = \%d\textbackslash n", a, b);};
  \draw[ptr] (2.4,-0.25) -- (2.4,-1.1);
  \node[font=\scriptsize,color=ptrline,anchor=north] at (2.4,-1.1) {第 1 个占位符填 \texttt{a}};
  \draw[ptr] (4.6,-0.25) -- (4.6,-1.8);
  \node[font=\scriptsize,color=ptrline,anchor=north] at (4.6,-1.8) {第 2 个占位符填 \texttt{b}};
\end{tikzpicture}
\end{center}

常用的格式符：

\begin{center}
\begin{tabular}{@{}lll@{}}
\toprule
格式符 & 类型 & 例子 \\
\midrule
\texttt{\%d} & 有符号十进制整数 & \texttt{printf("\%d", 42);} \\
\texttt{\%u} & 无符号十进制整数 & \texttt{printf("\%u", 42U);} \\
\texttt{\%x} & 十六进制（小写） & \texttt{printf("\%x", 255);} 输出 \texttt{ff} \\
\texttt{\%o} & 八进制 & \texttt{printf("\%o", 8);} 输出 \texttt{10} \\
\texttt{\%c} & 单个字符 & \texttt{printf("\%c", 'A');} \\
\texttt{\%f} & 浮点数（默认 6 位小数） & \texttt{printf("\%f", 3.14);} \\
\texttt{\%.2f} & 保留 2 位小数 & \texttt{printf("\%.2f", 3.14159);} 输出 \texttt{3.14} \\
\texttt{\%s} & 字符串 & \texttt{printf("\%s", "hi");} \\
\texttt{\%p} & 指针（地址） & \texttt{printf("\%p", (void *)\&a);} \\
\texttt{\%\%} & 输出一个百分号 & \texttt{printf("100\%\%");} \\
\bottomrule
\end{tabular}
\end{center}

\begin{pitfall}[格式符和参数类型必须匹配]
用 \texttt{\%d} 打印 \code{double}，或用 \texttt{\%f} 打印 \code{int}，属于未定义行为。编译器通常会给出警告（\code{-Wall} 打开时），但程序可能仍然“看起来能跑”。请务必让格式符和参数类型一一对应。
\end{pitfall}

\subsection{宽度、精度与对齐}

格式串里还可以指定“占多少格、保留几位小数、左对齐还是右对齐”：

\begin{codebox}{format.c —— 格式化输出}
#include <stdio.h>

int main(void)
{
    printf("[%5d]\n", 42);       /* 宽度 5，右对齐 */
    printf("[%-5d]\n", 42);      /* 宽度 5，左对齐 */
    printf("[%05d]\n", 42);      /* 宽度 5，用 0 填充 */
    printf("[%.3f]\n", 3.14159); /* 保留 3 位小数 */
    printf("[%8.2f]\n", 3.14159);/* 总宽 8，保留 2 位小数 */
    printf("[%c]\n", 'A');
    printf("[%s]\n", "hello");
    return 0;
}
\end{codebox}

\begin{outbox}{运行结果}
[   42]
[42   ]
[00042]
[3.142]
[    3.14]
[A]
[hello]
\end{outbox}

\begin{tipbox}[对齐表格的小技巧]
打印成绩单、实验数据时，用 \texttt{\%8.2f} 这样的固定宽度格式，可以让小数点对齐，输出整齐好看。
\end{tipbox}

\subsection{\texttt{scanf}：读入数据}

\code{scanf} 的名字来自 \en{scan formatted}。它按格式串读取输入，并把结果\textbf{存到变量的地址里}。

\begin{codebox}{input.c —— 读入两个整数}
#include <stdio.h>

int main(void)
{
    int a, b;

    printf("请输入两个整数：");
    scanf("%d%d", &a, &b);

    printf("它们的和是 %d\n", a + b);
    return 0;
}
\end{codebox}

\begin{outbox}{运行示例}
请输入两个整数：3 5
它们的和是 8
\end{outbox}

注意 \code{scanf} 的参数是 \texttt{\&a}、\texttt{\&b}，不是 \code{a}、\code{b}。为什么？

\begin{center}
\begin{tikzpicture}[x=1cm,y=1cm,>=Stealth]
  \node[ptrbox,minimum width=2cm] (a) at (0,0) {\texttt{a}};
  \node[font=\ttfamily\small] (addr) at (2.8,0) {0x1000};
  \node[font=\small,anchor=west] at (4.4,0) {变量的地址};
  \draw[ptr] (a.east) -- (addr.west);
  \node[font=\small,anchor=west] at (-1.2,1.2) {告诉 \texttt{scanf}：“把读到的数写到这个地址”};
\end{tikzpicture}
\end{center}

\texttt{\&} 是“取地址”运算符。\code{scanf} 需要知道把数据写到哪里，所以必须给它地址。关于地址和指针，第 8 章会详细讲，现在先记住“\code{scanf} 的变量前面要加 \texttt{\&}”。

\begin{pitfall}[\texttt{scanf} 的五个坑]
\begin{enumerate}
  \item \textbf{忘记 \texttt{\&}}：\texttt{scanf("\%d", a);} 会导致程序崩溃或读到错误位置。这是最经典的错误。
  \item \textbf{格式串里乱加空格}：\texttt{scanf("\%d, \%d", \&a, \&b);} 要求输入时也必须带逗号。初学者常在格式串里加 \code{\n}，这会让程序等待你多输入一些内容。
  \item \textbf{用 \texttt{\%f} 读 \code{double}}：\code{scanf} 读 \code{double} 必须用 \texttt{\%lf}（注意是字母 l 加 f），而 \code{printf} 里 \texttt{\%f} 和 \texttt{\%lf} 都可以。
  \item \textbf{输入缓冲区残留}：先用 \code{scanf} 读数字，再用 \code{getchar} 读字符时，会读到一个换行符。第 9 章会详细处理。
  \item \textbf{不检查返回值}：输入失败时 \code{scanf} 返回成功读到的项数。健壮的代码应该检查它，例如 \texttt{if (scanf("\%d", \&a) != 1) \{ ... \}}。
\end{enumerate}
\end{pitfall}

\subsection{\texttt{getchar} 与 \texttt{putchar}}

如果只需要读/写单个字符，可以用更简单的 \code{getchar} 和 \code{putchar}：

\begin{codebox}{getchar.c —— 读写单个字符}
#include <stdio.h>

int main(void)
{
    int c;

    printf("请输入一个字符：");
    c = getchar();          /* 从标准输入读一个字符，返回它的编码 */

    printf("你输入的是：");
    putchar(c);
    putchar('\n');

    return 0;
}
\end{codebox}

\begin{outbox}{运行示例}
请输入一个字符：A
你输入的是：A
\end{outbox}

注意这里用 \code{int c} 而不是 \code{char c}：因为 \code{getchar} 在遇到文件结束（EOF）时返回一个特殊值（通常是 $-1$），\code{char} 可能存不下。这个细节第 9 章会展开。

\section{表达式、语句与注释}

\subsection{语句}

语句是程序执行的最小单位。常见的语句有：

\begin{itemize}
  \item 声明语句：\code{int a;}
  \item 表达式语句：\code{a = 10;}、\code{printf("hi");}
  \item 复合语句：用花括号括起来的一组语句 \texttt{\{ ... \}}
  \item 空语句：只有一个分号 \code{;}（有时用于占位）
  \item 控制语句：\code{if}、\code{for}、\code{return} 等（后面几章讲）
\end{itemize}

\subsection{表达式}

表达式是由变量、常量、运算符组成的式子，\textbf{有值}。例如：

\begin{center}
\begin{tabular}{@{}ll@{}}
\toprule
表达式 & 值 \\
\midrule
\code{3 + 5} & 8 \\
\code{a * 2} & \code{a} 的两倍 \\
\code{a = 10} & 赋值表达式的值是 10（赋值本身也有值） \\
\code{printf("hi")} & 打印的字符数（通常用不到） \\
\bottomrule
\end{tabular}
\end{center}

在 C 语言里，\textbf{表达式加上分号就变成语句}。这一点很重要：\code{a + b;} 是合法的语句（虽然没什么用），而 \code{a = b} 是表达式。

\subsection{注释}

注释是写给人看的，编译器会忽略它。C89 只有一种注释：

\begin{codebox}{comment.c —— 注释的写法}
#include <stdio.h>

int main(void)
{
    /* 这是块注释，可以跨行
       也可以写很多行 */
    int a = 10;     /* 行尾注释 */

    /*
     * 这种星号开头的排版只是习惯，不是语法要求
     */
    printf("%d\n", a);
    return 0;
}
\end{codebox}

\begin{c89box}
C89 只支持 \code{/* ... */} 注释，\textbf{不支持} \code{//} 单行注释（那是 C99 引入的）。本书所有代码都遵守这一点。
\end{c89box}

\begin{tipbox}[注释写什么]
好的注释解释“为什么”，而不是“是什么”。\code{a = a + 1;} 不需要注释“让 a 加 1”，但可以注释“记录重试次数”。代码本身应该让人看懂“做什么”，注释用来补充意图、边界条件和坑。
\end{tipbox}

\section{顺序结构综合例题}

顺序结构的程序有一个固定套路：\textbf{输入 → 计算 → 输出}。下面用几个例子把这个套路练熟。

\subsection{例：摄氏度转华氏度}

\begin{example}
输入一个摄氏温度 $C$，输出对应的华氏温度 $F$。公式：$F = C \times \frac{9}{5} + 32$。
\end{example}
\begin{solution}
先看一个想当然的写法：

\begin{codebox}{temperature.c —— 温度转换（错误示范）}
#include <stdio.h>

int main(void)
{
    int c;
    int f;

    printf("请输入摄氏温度：");
    scanf("%d", &c);

    f = c * 9 / 5 + 32;
    printf("华氏温度 = %d\n", f);
    return 0;
}
\end{codebox}

这个程序能跑，但有一个隐藏问题：\code{9 / 5} 是整数除法，结果是 1，不是 1.8。所以 $C=20$ 时算出来是 $20 \times 1 + 32 = 52$，而不是 $68$。

正确写法是把至少一个操作数变成浮点数：
\begin{codebox}{temperature2.c —— 温度转换（正确写法）}
#include <stdio.h>

int main(void)
{
    double c;
    double f;

    printf("请输入摄氏温度：");
    scanf("%lf", &c);

    f = c * 9.0 / 5.0 + 32.0;
    printf("华氏温度 = %.2f\n", f);
    return 0;
}
\end{codebox}

\begin{outbox}{运行示例}
请输入摄氏温度：20
华氏温度 = 68.00
\end{outbox}

要点：\code{9.0 / 5.0} 是浮点除法；\code{scanf} 读 \code{double} 用 \texttt{\%lf}；输出保留两位小数用 \texttt{\%.2f}。
\end{solution}

\subsection{例：交换两个变量}

\begin{example}
输入两个整数 \code{a} 和 \code{b}，交换它们的值后输出。
\end{example}
\begin{solution}
不能直接写 \code{a = b; b = a;}——那样 \code{a} 原来的值会被覆盖。需要借助一个临时变量，就像交换两杯水需要第三个空杯子：

\begin{figure}[htbp]
  \centering
  \begin{tikzpicture}[x=1cm,y=1cm,>=Stealth]
    \node[membox,minimum width=2cm] (a) at (0,0) {a};
    \node[membox,minimum width=2cm] (b) at (2.5,0) {b};
    \node[emptybox,minimum width=2cm] (t) at (5,0) {temp};
    \node[font=\small] at (0,1.0) {3};
    \node[font=\small] at (2.5,1.0) {5};
    \node[font=\scriptsize,color=black!60] at (5,1.0) {?};

    \node[font=\small] at (0,-1.2) {初始};
    \node[font=\small] at (2.5,-1.2) {};
    \node[font=\small] at (5,-1.2) {};

    \draw[ptr] (0,-2.2) to[bend left] node[lbl,below]{① temp = a} (5,-2.2);
    \draw[ptr] (2.5,-2.9) to[bend left] node[lbl,below]{② a = b} (0,-2.9);
    \draw[ptr] (5,-3.6) to[bend left] node[lbl,below]{③ b = temp} (2.5,-3.6);
  \end{tikzpicture}
  \caption{交换两个变量需要借助临时变量}
  \label{fig:swap}
\end{figure}

\begin{codebox}{swap.c —— 交换两个变量}
#include <stdio.h>

int main(void)
{
    int a, b, temp;

    printf("请输入两个整数：");
    scanf("%d%d", &a, &b);

    printf("交换前：a = %d, b = %d\n", a, b);

    temp = a;       /* ① 把 a 存到临时变量 */
    a = b;          /* ② 把 b 写到 a */
    b = temp;       /* ③ 把临时变量写回 b */

    printf("交换后：a = %d, b = %d\n", a, b);
    return 0;
}
\end{codebox}

\begin{outbox}{运行示例}
请输入两个整数：3 5
交换前：a = 3, b = 5
交换后：a = 5, b = 3
\end{outbox}

这个“三步交换”在后面排序算法里会反复出现。
\end{solution}

\subsection{例：三位数反转}

\begin{example}
输入一个三位正整数，输出它的逆序数。例如输入 123，输出 321。
\end{example}
\begin{solution}
利用整数除法和取余：

\begin{itemize}
  \item 个位：\texttt{n \% 10}；
  \item 十位：\texttt{n / 10 \% 10}；
  \item 百位：\code{n / 100}。
\end{itemize}

\begin{codebox}{reverse.c —— 三位数反转}
#include <stdio.h>

int main(void)
{
    int n, ge, shi, bai, rev;

    printf("请输入一个三位数：");
    scanf("%d", &n);

    ge  = n % 10;
    shi = n / 10 % 10;
    bai = n / 100;

    rev = ge * 100 + shi * 10 + bai;
    printf("逆序数 = %d\n", rev);
    return 0;
}
\end{codebox}

\begin{outbox}{运行示例}
请输入一个三位数：123
逆序数 = 321
\end{outbox}

注意 \code{n / 100} 对 123 得到 1（整数除法丢弃小数部分），这正好是百位。
\end{solution}

\subsection{例：求平均值}

\begin{example}
输入三门课的成绩（可能有小数），输出平均分，保留两位小数。
\end{example}
\begin{solution}
思路如下：

\begin{codebox}{average.c —— 求平均分}
#include <stdio.h>

int main(void)
{
    double s1, s2, s3, average;

    printf("请输入三门课成绩：");
    scanf("%lf%lf%lf", &s1, &s2, &s3);

    average = (s1 + s2 + s3) / 3.0;
    printf("平均分 = %.2f\n", average);
    return 0;
}
\end{codebox}

\begin{outbox}{运行示例}
请输入三门课成绩：85.5 90 92.5
平均分 = 89.33
\end{outbox}

注意除以 \code{3.0} 而不是 \code{3}。虽然这里 \code{s1 + s2 + s3} 已经是 \code{double}，除以 \code{3} 也会做浮点除法（因为一个操作数是浮点），但写成 \code{3.0} 更清楚，也更不容易出错。
\end{solution}

\section{常见错误与调试}

初学阶段 90\% 的编译错误都来自下面几类。遇到报错，先对照这张表。

\begin{center}
\begin{tabular}{@{}p{0.34\textwidth}p{0.56\textwidth}@{}}
\toprule
错误现象 & 可能原因 \\
\midrule
\code{expected ';' before ...} & 上一行漏了分号 \\
\texttt{stray '\textbackslash 357' in program} & 用了中文标点（，；：）\\
\code{'printf' undeclared} & 忘了 \code{#include <stdio.h>} \\
\code{undefined reference to 'main'} & 没有 \code{main} 函数，或拼写错误 \\
\code{implicit declaration of function} & 函数没声明/没包含头文件 \\
输出乱码或数字奇怪 & 格式符与类型不匹配，或变量未初始化 \\
程序一闪而过 & 控制台程序运行完自动关闭（正常现象） \\
\bottomrule
\end{tabular}
\end{center}

\begin{tipbox}[读错误信息的顺序]
编译器报错时，\textbf{先看第一条}。后面的错误经常是第一条引起的连锁反应。错误信息里的文件名、行号、列号都要看，尤其是行号。
\end{tipbox}

\begin{deepbox}[拔高：用 \texttt{printf} 调试]
在没有调试器的情况下，最有效的调试手段是“打印中间结果”：
\begin{codebox}{debug.c —— 打印中间变量}
#include <stdio.h>

int main(void)
{
    int a = 10;
    int b = 0;
    int c;

    printf("debug: a=%d, b=%d\n", a, b);
    c = a / b;          /* 这里会出问题 */
    printf("debug: c=%d\n", c);
    return 0;
}
\end{codebox}
运行时会发现程序在除法处崩溃（除零）。通过打印可以看到 \code{b} 是 0。等学了第 4 章的分支，就可以在除法前先判断 \code{b != 0}。
\end{deepbox}

\section{拔高：变量在内存里的细节}

\subsection{声明与定义}

严格来说，“声明”是告诉编译器“有这么个东西”，“定义”是为它分配内存。对于局部变量，\code{int a;} 既是声明也是定义。全局变量和函数的情况稍微复杂一些，等学到多文件编程时再区分。现在你只要知道：\textbf{局部变量在栈上分配，函数返回时自动释放}。

\subsection{为什么局部变量不自动清零}

栈空间会被反复复用。函数 A 用完的栈空间，函数 B 接着用，里面还留着 A 的数据。所以 B 的局部变量如果没有初始化，就会看到“上次留下的垃圾”。这也是为什么“未初始化变量”是危险的。

\subsection{类型转换的初步印象}

当不同类型的值混在一起运算时，C 会自动做类型转换。例如：

\begin{codebox}{convert.c —— 隐式类型转换}
#include <stdio.h>

int main(void)
{
    int a = 7;
    int b = 2;
    double c;

    c = a / b;          /* 先做整数除法，得 3，再转成 3.0 */
    printf("c = %f\n", c);

    c = (double)a / b;  /* 先把 a 转成 double，再做浮点除法 */
    printf("c = %f\n", c);

    return 0;
}
\end{codebox}

\begin{outbox}{运行结果}
c = 3.000000
c = 3.500000
\end{outbox}

\code{(double)a} 叫\textbf{强制类型转换}，第 3 章会系统讲。现在记住：\textbf{整数除法会截断小数部分}，想要精确结果，先把操作数转成浮点。

\section{本章小结}

\begin{center}
\begin{tabular}{@{}p{0.3\textwidth}p{0.6\textwidth}@{}}
\toprule
主题 & 要点 \\
\midrule
程序骨架 & 预处理指令、\code{int main(void)}、函数体、\code{return 0;} \\
语句 & 以分号结尾；花括号组成复合语句 \\
变量 & 名字 + 类型 + 值 + 地址；先声明后使用 \\
C89 声明 & 所有声明必须在块开头，执行语句之前 \\
类型 & \code{char}/\code{int}/\code{float}/\code{double}；大小用 \code{sizeof} 查 \\
输出 & \code{printf} 格式串 + 占位符；格式符与类型必须匹配 \\
输入 & \code{scanf} 要传地址（\texttt{\&a}）；注意返回值与缓冲区 \\
注释 & C89 只用 \code{/* */} \\
顺序结构 & 输入 → 计算 → 输出，自上而下执行 \\
\bottomrule
\end{tabular}
\end{center}

\begin{tipbox}[下一章预告]
顺序结构只能“一条路走到黑”。真实程序需要根据条件做选择（分支），也需要重复做某件事（循环）。在学分支之前，我们先把运算符讲透——尤其是关系运算和逻辑运算，它们是“判断”的基础。
\end{tipbox}

\begin{problemset}
  \item 写一个程序，输出你的姓名、专业和一句座右铭，每行一项。
  \item 输入两个整数，输出它们的和、差、积、商和余数。
  \item 输入一个华氏温度，输出摄氏温度（公式 $C = (F - 32) \times 5 / 9$），保留两位小数。
  \item 输入圆的半径 $r$，输出周长和面积（$\pi$ 取 3.1415926），保留三位小数。
  \item 输入一个四位数，输出它的各位数字之和。例如 2026，输出 10。
  \item 输入两个整数并交换，要求用临时变量实现，并打印交换前后的值。
  \item 下面的程序有几处错误？请找出来并改正。
\begin{codebox}{find-bugs.c —— 找错题}
#include <stdio.h>

int main(void)
{
    int a = 1
    int b = 2;
    sum = a + b;
    printf("%d\n", sum)
    return 0;
}
\end{codebox}
  \item 解释 \texttt{printf("\%d", 3.14);} 为什么是错的，应该怎么写。
  \item 输入一个三位数，输出它的各位数字之和与逆序数。
  \item 思考题：为什么 \texttt{scanf("\%d", \&a);} 需要 \texttt{\&}，而 \texttt{printf("\%d", a);} 不需要？
\end{problemset}

\chapter{写在分支结构之前：运算符}

\begin{introduction}
  \item 算术、赋值、自增自减：最常见的四类运算
  \item 关系与逻辑：为“判断”准备的两把尺子
  \item 位运算：直接操作二进制位
  \item 优先级、结合性、类型转换与求值顺序
  \item 副作用与未定义行为：那些“看起来对”的坑
\end{introduction}

\chaptermotto{「程序 = 数据 + 运算 + 控制。先把运算说清楚，控制才有着落。」}

下一章要学分支，也就是“如果满足某个条件，就做某件事”。条件从哪来？从\textbf{关系运算}和\textbf{逻辑运算}来。所以这一章是分支的“前置技能包”。

除此之外，本章还会把 C 语言里所有常用运算符过一遍。内容看起来多，但核心只有三句话：

\begin{enumerate}
  \item 每个运算符都有\textbf{优先级}和\textbf{结合性}，不确定就加括号；
  \item 不同类型混算会发生\textbf{类型转换}，转换规则要心里有数；
  \item 有些表达式有\textbf{副作用}，求值顺序不对结果就不可预测。
\end{enumerate}

\section{表达式与运算符总览}

\subsection{什么是运算符}

运算符是告诉编译器“做什么运算”的符号，例如 \code{+}、\code{*}、\code{<}、\texttt{\&\&}。参与运算的数据叫\textbf{操作数}。

\begin{center}
\begin{tabular}{@{}lll@{}}
\toprule
分类 & 运算符 & 说明 \\
\midrule
算术 & \code{+} \code{-} \code{*} \code{/} \texttt{\%} & 加减乘除、取余 \\
赋值 & \code{=} \code{+=} \code{-=} \code{*=} \code{/=} \texttt{\%=} & 把值写进变量 \\
自增自减 & \code{++} \code{--} & 加 1 / 减 1 \\
关系 & \code{>} \code{<} \code{>=} \code{<=} \code{==} \code{!=} & 比较大小、判断相等 \\
逻辑 & \texttt{\&\&} \code{||} \code{!} & 与、或、非 \\
位运算 & \texttt{\&} \code{|} \code{^} \code{~} \code{<<} \code{>>} & 按二进制位操作 \\
条件 & \code{?:} & 三目运算符 \\
逗号 & \code{,} & 顺序求值 \\
其他 & \code{sizeof}、强制类型转换 \code{(类型)} & 求大小、转换类型 \\
\bottomrule
\end{tabular}
\end{center}

\subsection{优先级与结合性：先算谁}

一个表达式里出现多个运算符时，先算谁由\textbf{优先级}决定；优先级相同的，由\textbf{结合性}决定（从左往右还是从右往左）。

例如 \code{2 + 3 * 4}，因为 \code{*} 的优先级高于 \code{+}，所以先算 \code{3 * 4}，结果是 14。如果先算加法，结果会是 20。这就是优先级的差别。

\begin{tipbox}[记不住优先级怎么办]
加括号。括号的优先级最高，而且能让读代码的人一眼看懂你的意图。经验法则是：\textbf{只要混合了三种以上运算符，就加括号}。
\end{tipbox}

完整的优先级表（从高到低）放在本章后面，先学具体运算符，再回来看表。

\section{算术运算符}

\subsection{加、减、乘、除}

\begin{codebox}{arith.c —— 四则运算}
#include <stdio.h>

int main(void)
{
    int a = 7;
    int b = 2;

    printf("a + b = %d\n", a + b);
    printf("a - b = %d\n", a - b);
    printf("a * b = %d\n", a * b);
    printf("a / b = %d\n", a / b);
    return 0;
}
\end{codebox}

\begin{outbox}{运行结果}
a + b = 9
a - b = 5
a * b = 14
a / b = 3
\end{outbox}

注意 \code{7 / 2} 的结果是 3，不是 3.5。因为两个操作数都是整数，C 执行的是\textbf{整数除法}：只保留商的整数部分，小数部分直接丢弃（不是四舍五入）。

\begin{pitfall}[整数除法向零截断]
\code{7 / 2} 得 3，\code{-7 / 2} 得 $-3$（向零截断，不是向下取整 $-4$）。如果你需要小数结果，把至少一个操作数变成浮点：\code{7.0 / 2} 得 3.5。
\end{pitfall}

\subsection{取余运算符 \texttt{\%}}

\texttt{\%} 计算除法后的余数，只能用于整数：

\begin{codebox}{mod.c —— 取余}
#include <stdio.h>

int main(void)
{
    printf("7 %% 2 = %d\n", 7 % 2);
    printf("8 %% 2 = %d\n", 8 % 2);
    printf("7 %% 3 = %d\n", 7 % 3);
    return 0;
}
\end{codebox}

\begin{outbox}{运行结果}
7 % 2 = 1
8 % 2 = 0
7 % 3 = 1
\end{outbox}

取余的常见用途：

\begin{itemize}
  \item \textbf{判断奇偶}：\texttt{n \% 2} 为 0 是偶数，非 0 是奇数；
  \item \textbf{取某一位数字}：\texttt{n \% 10} 取个位；
  \item \textbf{循环计数}：让计数器在 $0 \sim N-1$ 之间循环；
  \item \textbf{时间换算}：\texttt{seconds \% 60} 得到秒数。
\end{itemize}

\begin{pitfall}[负数取余要小心]
在多数现代编译器上，\texttt{-7 \% 2} 的结果是 $-1$，不是 1。所以判断奇数不要写 \texttt{n \% 2 == 1}，而要写 \texttt{n \% 2 != 0}，否则负数会被误判为偶数。C89 对负数除法和取余的符号没有强制规定，跨平台代码更要注意。
\end{pitfall}

\subsection{除零与溢出}

整数除以 0 是\textbf{未定义行为}，程序通常会直接崩溃：

\begin{codebox}{divzero.c —— 除零}
#include <stdio.h>

int main(void)
{
    int a = 10;
    int b = 0;

    printf("%d\n", a / b);   /* 运行时错误 */
    return 0;
}
\end{codebox}

浮点数除以 0 不会崩溃，而是得到“无穷大”或“非数”（NaN），但同样不应该依赖。

整数运算还可能溢出。第 1 章已经讲过，有符号溢出是未定义行为。看一个例子：

\begin{codebox}{int-overflow.c —— 整数运算溢出}
#include <stdio.h>

int main(void)
{
    int a = 100000;
    int b = 100000;
    int c;

    c = a * b;       /* 10^10 超出了 int 的范围 */
    printf("c = %d\n", c);
    return 0;
}
\end{codebox}

\begin{outbox}{运行结果（常见实现，不可依赖）}
c = 1410065408
\end{outbox}

\begin{tipbox}[电控中的溢出]
电机编码器计数、定时器计数、ADC 采样值累加都可能溢出。做乘加运算前，先估算最大值是否会超过类型范围；必要时用 \code{long}、\code{double} 或先做除法。
\end{tipbox}

\section{赋值运算符}

\subsection{基本赋值}

\code{=} 是赋值，不是“等于”。它把右边的值写进左边的变量：

\begin{codebox}{assign.c —— 赋值}
#include <stdio.h>

int main(void)
{
    int a;
    int b;

    a = 10;         /* 把 10 写进 a */
    b = a;          /* 把 a 的值复制给 b */
    a = 20;         /* 再改 a，b 不受影响 */

    printf("a = %d, b = %d\n", a, b);
    return 0;
}
\end{codebox}

\begin{outbox}{运行结果}
a = 20, b = 10
\end{outbox}

赋值表达式的值就是被赋的值，所以可以连续赋值：

\begin{codebox}{chain.c —— 连续赋值}
#include <stdio.h>

int main(void)
{
    int a, b, c;

    a = b = c = 5;      /* 从右往左：c=5，b=5，a=5 */
    printf("a=%d b=%d c=%d\n", a, b, c);
    return 0;
}
\end{codebox}

\subsection{复合赋值}

\code{a += b} 等价于 \code{a = a + b}。常见的复合赋值有：

\begin{center}
\begin{tabular}{@{}ll@{}}
\toprule
写法 & 等价于 \\
\midrule
\code{a += b} & \code{a = a + b} \\
\code{a -= b} & \code{a = a - b} \\
\code{a *= b} & \code{a = a * b} \\
\code{a /= b} & \code{a = a / b} \\
\texttt{a \%= b} & \texttt{a = a \% b} \\
\code{a <<= 2} & \code{a = a << 2} \\
\texttt{a \&= mask} & \texttt{a = a \& mask} \\
\code{a |= flag} & \code{a = a | flag} \\
\bottomrule
\end{tabular}
\end{center}

复合赋值的好处是简洁，而且左边只求值一次（对指针和数组下标有意义，后面会讲）。但不要为了简洁牺牲可读性。

\begin{pitfall}[\texttt{a += b * c} 不是 \texttt{(a + b) * c}]
\code{a += b * c} 等价于 \code{a = a + (b * c)}，因为 \code{*} 的优先级高于 \code{+=}。复合赋值运算符的优先级很低，仅高于逗号。
\end{pitfall}

\section{自增与自减}

\code{++} 让变量加 1，\code{--} 让变量减 1。它们有两种写法：

\begin{center}
\begin{tabular}{@{}lll@{}}
\toprule
写法 & 名称 & 含义 \\
\midrule
\code{++a} & 前置 & 先把 \code{a} 加 1，再使用 \code{a} 的新值 \\
\code{a++} & 后置 & 先使用 \code{a} 的旧值，再把 \code{a} 加 1 \\
\code{--a} & 前置 & 先减 1，再使用 \\
\code{a--} & 后置 & 先使用，再减 1 \\
\bottomrule
\end{tabular}
\end{center}

\begin{codebox}{inc.c —— 前置与后置}
#include <stdio.h>

int main(void)
{
    int a = 5;
    int b;

    b = ++a;        /* a 先变 6，再把 6 赋给 b */
    printf("a=%d b=%d\n", a, b);

    a = 5;
    b = a++;        /* 先把 5 赋给 b，a 再变 6 */
    printf("a=%d b=%d\n", a, b);

    return 0;
}
\end{codebox}

\begin{outbox}{运行结果}
a=6 b=6
a=6 b=5
\end{outbox}

\begin{figure}[htbp]
  \centering
  \begin{tikzpicture}[x=1cm,y=1cm,>=Stealth]
    \node[process,minimum width=34mm] (pre) at (0,0) {\texttt{++a}};
    \node[process,minimum width=34mm] (post) at (6,0) {\texttt{a++}};
    \node[box,minimum width=34mm,below=8mm of pre] (pre2) {先加 1};
    \node[box,minimum width=34mm,below=8mm of post] (post2) {先取值};
    \node[box,minimum width=34mm,below=6mm of pre2] (pre3) {再取值使用};
    \node[box,minimum width=34mm,below=6mm of post2] (post3) {再加 1};
    \draw[flow] (pre) -- (pre2);
    \draw[flow] (pre2) -- (pre3);
    \draw[flow] (post) -- (post2);
    \draw[flow] (post2) -- (post3);
  \end{tikzpicture}
  \caption{前置 \texttt{++} 与后置 \texttt{++} 的执行顺序}
  \label{fig:inc}
\end{figure}

单独使用 \code{a++;} 时，前置和后置效果一样，都是让 \code{a} 加 1。区别只在“参与其他运算”时体现。

\begin{pitfall}[不要在同一个表达式里对同一变量多次自增]
\code{i = i++;}、\code{a = i++ + ++i;} 这类写法在 C 标准里属于\textbf{未定义行为}，不同编译器可能给出不同结果。原因是 C 没有规定“先算左边还是先算右边”。记住：\textbf{一个变量在一个表达式里最多被修改一次}。
\end{pitfall}

\begin{deepbox}[拔高：为什么 \texttt{i = i++} 是未定义的]
\code{i++} 包含两个动作：读取 \code{i} 的旧值、把 \code{i} 加 1 写回。\code{i = ...} 又要把结果写回 \code{i}。两次写入之间没有“顺序点”（sequence point），标准无法规定谁先谁后，于是干脆规定这是未定义行为。现代编译器打开优化后，结果可能和你想的完全不同。
\end{deepbox}

\section{关系运算符}

关系运算符用来比较两个值，结果只有两种：\textbf{真（1）}或\textbf{假（0）}。

\begin{center}
\begin{tabular}{@{}lll@{}}
\toprule
运算符 & 含义 & 例子（\code{a=3, b=5}） \\
\midrule
\code{>} & 大于 & \code{a > b} 为 0 \\
\code{<} & 小于 & \code{a < b} 为 1 \\
\code{>=} & 大于等于 & \code{a >= 3} 为 1 \\
\code{<=} & 小于等于 & \code{a <= 2} 为 0 \\
\code{==} & 等于 & \code{a == 3} 为 1 \\
\code{!=} & 不等于 & \code{a != b} 为 1 \\
\bottomrule
\end{tabular}
\end{center}

\begin{codebox}{relational.c —— 关系运算的结果}
#include <stdio.h>

int main(void)
{
    int a = 3;
    int b = 5;

    printf("a < b  = %d\n", a < b);
    printf("a > b  = %d\n", a > b);
    printf("a == b = %d\n", a == b);
    printf("a != b = %d\n", a != b);
    return 0;
}
\end{codebox}

\begin{outbox}{运行结果}
a < b  = 1
a > b  = 0
a == b = 0
a != b = 1
\end{outbox}

\begin{pitfall}[\texttt{=} 和 \texttt{==} 是两回事]
\code{if (a = 5)} 是赋值，把 5 写进 \code{a}，然后判断 5 是否为真（永远为真）；\code{if (a == 5)} 才是比较。这个错误编译器通常只会给警告，一定要看清。
\end{pitfall}

\begin{pitfall}[不要连写比较]
数学里可以写 $3 < x < 10$，C 语言里 \code{3 < x < 10} 的含义完全不同：先算 \code{3 < x} 得到 0 或 1，再拿这个 0/1 和 10 比较，结果永远为真。正确写法是 \texttt{3 < x \&\& x < 10}。
\end{pitfall}

\section{逻辑运算符}

逻辑运算符把多个条件组合起来：

\begin{center}
\begin{tabular}{@{}lll@{}}
\toprule
运算符 & 名称 & 含义 \\
\midrule
\texttt{\&\&} & 逻辑与 & 两个都为真，结果才为真 \\
\code{||} & 逻辑或 & 有一个为真，结果就为真 \\
\code{!} & 逻辑非 & 真变假，假变真 \\
\bottomrule
\end{tabular}
\end{center}

真值表：

\begin{center}
\begin{tabular}{@{}ccccc@{}}
\toprule
\code{a} & \code{b} & \texttt{a \&\& b} & \code{a || b} & \code{!a} \\
\midrule
0 & 0 & 0 & 0 & 1 \\
0 & 非 0 & 0 & 1 & 1 \\
非 0 & 0 & 0 & 1 & 0 \\
非 0 & 非 0 & 1 & 1 & 0 \\
\bottomrule
\end{tabular}
\end{center}

在 C 语言里，\textbf{0 为假，非 0 为真}。任何非零值（包括负数、浮点数）都算真。

\subsection{短路求值}

\texttt{\&\&} 和 \code{||} 有一个非常重要的特性：\textbf{短路}。

\begin{itemize}
  \item \texttt{a \&\& b}：如果 \code{a} 为假，就不再计算 \code{b}，直接得 0；
  \item \code{a || b}：如果 \code{a} 为真，就不再计算 \code{b}，直接得 1。
\end{itemize}

这个特性有两个用途：一是提高效率，二是\textbf{避免危险操作}。

\begin{codebox}{short-circuit.c —— 短路求值避免除零}
#include <stdio.h>

int main(void)
{
    int a = 10;
    int b = 0;

    if (b != 0 && a / b > 2) {
        printf("商大于 2\n");
    } else {
        printf("不满足条件\n");
    }
    return 0;
}
\end{codebox}

\begin{outbox}{运行结果}
不满足条件
\end{outbox}

因为 \code{b != 0} 为假，\texttt{\&\&} 后面的 \code{a / b} 根本不会被计算，所以不会除零崩溃。如果写成 \texttt{a / b > 2 \&\& b != 0}，就会先除零，程序崩溃。

\begin{pitfall}[短路只保证“不被求值”，不保证“顺序好”]
虽然 \texttt{\&\&} 和 \code{||} 保证从左到右求值，但其余运算符（如 \code{+}、\code{*}、函数参数）的求值顺序是未指定的。不要把有副作用的表达式塞进复杂表达式里。
\end{pitfall}

\subsection{德·摩根定律}

逻辑表达式可以互相转换：

\[
!(a \&\& b) \equiv (!a) || (!b), \qquad
!(a || b) \equiv (!a) \&\& (!b).
\]

例如“不是（年龄大于 18 且身高大于 170）”等价于“年龄不大于 18 或身高不大于 170”。化简逻辑表达式时很有用。

\section{位运算符}

位运算符直接操作整数的二进制位。它在电控里非常常用：寄存器配置、标志位、掩码、协议解析，都离不开它。

\begin{center}
\begin{tabular}{@{}lll@{}}
\toprule
运算符 & 名称 & 作用 \\
\midrule
\texttt{\&} & 按位与 & 两位都为 1，结果才为 1 \\
\code{|} & 按位或 & 有一位为 1，结果就为 1 \\
\code{^} & 按位异或 & 两位不同，结果为 1 \\
\code{~} & 按位取反 & 0 变 1，1 变 0 \\
\code{<<} & 左移 & 整体左移，右边补 0 \\
\code{>>} & 右移 & 整体右移 \\
\bottomrule
\end{tabular}
\end{center}

\subsection{按位与、或、异或、取反}

以 8 位为例：

\begin{center}
\begin{tabular}{@{}cccc@{}}
\toprule
运算 & 例子 & 结果 & 说明 \\
\midrule
\texttt{\&} & \texttt{1010 1100 \& 0000 1111} & \code{0000 1100} & 取低 4 位 \\
\code{|} & \code{1010 0000 | 0000 1111} & \code{1010 1111} & 置低 4 位为 1 \\
\code{^} & \code{1010 1100 ^ 1111 0000} & \code{0101 1100} & 翻转高 4 位 \\
\code{~} & \code{~1010 1100} & \code{0101 0011} & 按位取反 \\
\bottomrule
\end{tabular}
\end{center}

\subsection{移位}

\code{x << n} 把 \code{x} 的二进制位整体左移 \code{n} 位，右边补 0；\code{x >> n} 右移。

\begin{codebox}{shift.c —— 移位}
#include <stdio.h>

int main(void)
{
    unsigned int x = 1;

    printf("1 << 0 = %u\n", x << 0);
    printf("1 << 1 = %u\n", x << 1);
    printf("1 << 4 = %u\n", x << 4);
    printf("1 << 8 = %u\n", x << 8);
    return 0;
}
\end{codebox}

\begin{outbox}{运行结果}
1 << 0 = 1
1 << 1 = 2
1 << 4 = 16
1 << 8 = 256
\end{outbox}

左移 $n$ 位相当于乘以 $2^n$，右移 $n$ 位相当于除以 $2^n$（对无符号数）。但\textbf{不要用移位代替乘除来“优化”}，现代编译器比你聪明，而且移位遇到负数或溢出时会出问题。

\begin{pitfall}[移位位数不能超过类型宽度]
对 32 位 \code{int} 左移 32 位是未定义行为。移位位数必须小于类型宽度。另外，有符号数右移时高位补 0 还是补 1 由实现决定（算术右移通常补符号位），所以位操作用 \code{unsigned} 类型最安全。
\end{pitfall}

\subsection{位操作的经典用法}

位操作最常用的四个动作是：\textbf{判断某位、置位、清位、取反某位}。

\begin{codebox}{bits.c —— 位操作的四个动作}
#include <stdio.h>

int main(void)
{
    unsigned int flags = 0;

    /* 置位：把第 3 位置 1 */
    flags = flags | (1u << 3);
    printf("置位后：%u\n", flags);

    /* 判断：第 3 位是否为 1 */
    if (flags & (1u << 3)) {
        printf("第 3 位是 1\n");
    }

    /* 清位：把第 3 位置 0 */
    flags = flags & ~(1u << 3);
    printf("清位后：%u\n", flags);

    /* 翻转：把第 0 位取反 */
    flags = flags ^ (1u << 0);
    printf("翻转后：%u\n", flags);

    return 0;
}
\end{codebox}

\begin{outbox}{运行结果}
置位后：8
第 3 位是 1
清位后：0
翻转后：1
\end{outbox}

\begin{tipbox}[记住这四个式子]
\begin{itemize}
  \item 置位：\code{flags |= (1u << n);}
  \item 清位：\texttt{flags \&= \textasciitilde{}(1u << n);}
  \item 判断：\texttt{if (flags \& (1u << n))}
  \item 翻转：\code{flags ^= (1u << n);}
\end{itemize}
其中 \code{1u} 表示无符号的 1，避免有符号移位的坑。
\end{tipbox}

\begin{deepbox}[拔高：寄存器配置就是这么干的]
假设某个 32 位寄存器里，第 0\textasciitilde3 位控制电机使能，第 4\textasciitilde7 位控制方向，第 8\textasciitilde15 位是速度。你可能会写：
\begin{codebox}{register.c —— 寄存器位域操作（示意）}
unsigned int reg = 0;

reg |= (1u << 0);                 /* 使能 */
reg = (reg & ~(0xFu << 4)) | (3u << 4);   /* 方向设为 3 */
reg = (reg & ~(0xFFu << 8)) | (200u << 8);/* 速度设为 200 */
\end{codebox}
这段代码不需要访问真实硬件，但思想是通用的：先清掉要改的位，再用“或”写入新值。直接赋值 \code{reg = 3 << 4} 会把其他位清零，是初学者常犯的错误。
\end{deepbox}

\section{条件运算符}

\code{?:} 是 C 语言里唯一的三目运算符，形式是：

\begin{center}
\code{条件 ? 表达式1 : 表达式2}
\end{center}

条件为真时取表达式 1 的值，否则取表达式 2 的值。

\begin{codebox}{ternary.c —— 条件运算符}
#include <stdio.h>

int main(void)
{
    int a = 7;
    int b = 3;
    int max;

    max = (a > b) ? a : b;
    printf("较大的数是 %d\n", max);

    printf("a 是%s\n", (a % 2 == 0) ? "偶数" : "奇数");
    return 0;
}
\end{codebox}

\begin{outbox}{运行结果}
较大的数是 7
a 是奇数
\end{outbox}

条件运算符适合“二选一”的简单场景。如果逻辑复杂，还是用 \code{if-else} 更清楚。

\begin{pitfall}[条件运算符的优先级很低]
\code{a > b ? a : b + 1} 会被解析成 \code{a > b ? a : (b + 1)}，而不是 \code{(a > b ? a : b) + 1}。不确定就加括号。
\end{pitfall}

\section{逗号、\texttt{sizeof} 与强制类型转换}

\subsection{逗号运算符}

逗号表达式依次计算每个子表达式，整个表达式的值是\textbf{最后一个子表达式的值}：

\begin{codebox}{comma.c —— 逗号运算符}
#include <stdio.h>

int main(void)
{
    int a;
    int b;

    a = (1, 2, 3);      /* 依次计算 1、2、3，a 得 3 */
    b = (a++, a + 10);  /* 先 a 加 1，再算 a+10 */

    printf("a = %d, b = %d\n", a, b);
    return 0;
}
\end{codebox}

\begin{outbox}{运行结果}
a = 4, b = 14
\end{outbox}

逗号运算符主要用于 \code{for} 循环的初始化和更新部分（第 5 章会看到），平时不要滥用。

\subsection{\texttt{sizeof}}

\code{sizeof} 不是函数，是运算符。它返回某个类型或变量占用的字节数，结果类型是 \code{size_t}（无符号整数）。

\begin{codebox}{sizeof2.c —— sizeof 的用法}
#include <stdio.h>

int main(void)
{
    int a;
    double b;
    int arr[10];

    printf("sizeof(int)    = %d\n", (int)sizeof(int));
    printf("sizeof(a)      = %d\n", (int)sizeof(a));
    printf("sizeof(double) = %d\n", (int)sizeof(double));
    printf("sizeof(b)      = %d\n", (int)sizeof(b));
    printf("sizeof(arr)    = %d\n", (int)sizeof(arr));
    printf("数组元素个数   = %d\n", (int)(sizeof(arr) / sizeof(arr[0])));
    return 0;
}
\end{codebox}

\begin{outbox}{运行结果（典型 32 位 int 环境）}
sizeof(int)    = 4
sizeof(a)      = 4
sizeof(double) = 8
sizeof(b)      = 8
sizeof(arr)    = 40
数组元素个数   = 10
\end{outbox}

\code{sizeof(arr) / sizeof(arr[0])} 是求数组元素个数的经典写法，第 6 章会详细讲。

\subsection{强制类型转换}

强制类型转换的语法是 \code{(类型)表达式}：

\begin{codebox}{cast.c —— 强制类型转换}
#include <stdio.h>

int main(void)
{
    int a = 7;
    int b = 2;
    double c;

    c = a / b;              /* 整数除法，得 3，再转 double */
    printf("c = %f\n", c);

    c = (double)a / b;      /* 先把 a 转成 double，再除 */
    printf("c = %f\n", c);

    printf("(int)3.99 = %d\n", (int)3.99);  /* 截断小数部分 */
    return 0;
}
\end{codebox}

\begin{outbox}{运行结果}
c = 3.000000
c = 3.500000
(int)3.99 = 3
\end{outbox}

\begin{pitfall}[强制转换可能丢失精度]
\code{(int)3.99} 得 3，不是 4；\code{(char)300} 会溢出，结果由实现决定。强制转换是“我确定要这样”，不是“帮我修好”。
\end{pitfall}

\section{优先级与结合性总表}

下面这张表从高到低排列。同一行的运算符优先级相同。

\begin{center}
\small
\begin{tabular}{@{}clll@{}}
\toprule
优先级 & 运算符 & 含义 & 结合性 \\
\midrule
1 & \code{()} \code{[]} \code{->} \code{.} & 函数调用、下标、成员 & 左 \\
2 & \code{!} \code{~} \code{++} \code{--} \code{+} \code{-} \code{*} \texttt{\&} \code{(类型)} \code{sizeof} & 单目运算 & 右 \\
3 & \code{*} \code{/} \texttt{\%} & 乘除、取余 & 左 \\
4 & \code{+} \code{-} & 加减 & 左 \\
5 & \code{<<} \code{>>} & 移位 & 左 \\
6 & \code{<} \code{<=} \code{>} \code{>=} & 关系 & 左 \\
7 & \code{==} \code{!=} & 相等/不等 & 左 \\
8 & \texttt{\&} & 按位与 & 左 \\
9 & \code{^} & 按位异或 & 左 \\
10 & \code{|} & 按位或 & 左 \\
11 & \texttt{\&\&} & 逻辑与 & 左 \\
12 & \code{||} & 逻辑或 & 左 \\
13 & \code{?:} & 条件 & 右 \\
14 & \code{=} \code{+=} \code{-=} \code{*=} \code{/=} \texttt{\%=} \code{<<=} \code{>>=} \texttt{\&=} \code{^=} \code{|=} & 赋值 & 右 \\
15 & \code{,} & 逗号 & 左 \\
\bottomrule
\end{tabular}
\end{center}

\begin{example}
分析表达式 \texttt{a + b * c > d \&\& e} 的求值顺序。
\end{example}
\begin{solution}
按优先级从高到低：

\begin{enumerate}
  \item \code{b * c}（乘法优先级 3）；
  \item \code{a + (b * c)}（加法优先级 4）；
  \item \code{(a + b * c) > d}（关系优先级 6，结果是 0 或 1）；
  \item \texttt{((a + b * c) > d) \&\& e}（逻辑与优先级 11）。
\end{enumerate}

\begin{center}
\begin{tikzpicture}[x=1cm,y=1cm,>=Stealth,level distance=9mm,
  level 1/.style={sibling distance=32mm},
  level 2/.style={sibling distance=18mm},
  level 3/.style={sibling distance=12mm}]
  \node[box,minimum width=8mm] {\texttt{\&\&}}
    child {node[box,minimum width=8mm] {\code{>}}
      child {node[box,minimum width=8mm] {\code{+}}
        child {node[box,minimum width=8mm] {\code{a}}}
        child {node[box,minimum width=8mm] {\code{*}}
          child {node[box,minimum width=8mm] {\code{b}}}
          child {node[box,minimum width=8mm] {\code{c}}}}}
      child {node[box,minimum width=8mm] {\code{d}}}}
    child {node[box,minimum width=8mm] {\code{e}}};
\end{tikzpicture}
\end{center}
\end{solution}

\section{类型转换}

不同类型的数据一起运算时，C 会自动做类型转换。规则可以概括为“向精度更高、范围更大的类型看齐”。

\subsection{整型提升}

\code{char}、\code{short} 在参与运算前，会先被提升为 \code{int}：

\begin{codebox}{promote.c —— 整型提升}
#include <stdio.h>

int main(void)
{
    char a = 100;
    char b = 100;
    int c;

    c = a + b;      /* a、b 先提升为 int，再相加，不会溢出 */
    printf("c = %d\n", c);
    return 0;
}
\end{codebox}

\begin{outbox}{运行结果}
c = 200
\end{outbox}

如果直接用一个 \code{char} 存 200，会溢出（\code{char} 范围通常是 $-128 \sim 127$）。这里因为先提升成 \code{int} 再相加，所以结果正确。

\subsection{通常算术转换}

当 \code{int} 和 \code{double} 运算时，\code{int} 会被转成 \code{double}；当 \code{int} 和 \code{unsigned int} 运算时，\code{int} 会被转成 \code{unsigned int}。后者是很多隐蔽 bug 的来源：

\begin{codebox}{unsigned-trap.c —— 有符号与无符号混算}
#include <stdio.h>

int main(void)
{
    int a = -1;
    unsigned int b = 1;

    if (a < b) {
        printf("a < b\n");
    } else {
        printf("a >= b\n");
    }
    return 0;
}
\end{codebox}

\begin{outbox}{运行结果（常见实现）}
a >= b
\end{outbox}

为什么？因为 \code{a} 被转换成了 \code{unsigned int}，$-1$ 变成 4294967295，当然大于 1。这类问题在比较 \code{sizeof} 的结果时特别常见：

\begin{codebox}{sizeof-trap.c —— sizeof 比较陷阱}
#include <stdio.h>

int main(void)
{
    int i;

    for (i = 10; i >= 0; i--) {
        if (i < sizeof(int)) {     /* i 被转成无符号，比较结果可能反了 */
            printf("%d\n", i);
        }
    }
    return 0;
}
\end{codebox}

\begin{tipbox}[避免有符号/无符号混算]
尽量让参与比较的两个变量类型一致；如果必须混用，先显式转换，或者把 \code{sizeof} 的结果强制转成 \code{int}。
\end{tipbox}

\section{求值顺序与副作用}

\subsection{什么是副作用}

表达式除了产生值，还可能修改环境，这叫\textbf{副作用}。赋值、自增自减、函数调用都可能有副作用。

C 标准只规定了少数几个“顺序点”（sequence point），在顺序点之前的所有副作用必须完成。顺序点包括：

\begin{itemize}
  \item 每条完整语句结束的分号；
  \item \texttt{\&\&}、\code{||}、\code{?:}、逗号运算符的求值处；
  \item 函数调用前（参数求值完成后）。
\end{itemize}

\begin{pitfall}[函数参数的求值顺序不确定]
\texttt{printf("\%d \%d", i++, i++);} 的结果在不同编译器上可能不同。不要写这种代码。同理，\code{f(i++, i)} 也不可预测。
\end{pitfall}

\subsection{一个经典例子}

\begin{codebox}{sequence.c —— 未定义行为示例}
#include <stdio.h>

int main(void)
{
    int i = 1;
    int a;

    a = i++ + i++;      /* 未定义行为！不要这样写 */
    printf("a = %d, i = %d\n", a, i);
    return 0;
}
\end{codebox}

这段代码在不同编译器、不同优化级别下可能输出不同结果。正确的写法是拆成多条语句：

\begin{codebox}{sequence-ok.c —— 拆开写}
#include <stdio.h>

int main(void)
{
    int i = 1;
    int a;

    a = i;
    i++;
    a = a + i;
    i++;
    printf("a = %d, i = %d\n", a, i);
    return 0;
}
\end{codebox}

\begin{tipbox}[一条语句只做一件事]
这不是语法要求，而是可靠性要求。一条语句里修改多个变量、或者对同一变量多次修改，都会让代码难以理解和调试。
\end{tipbox}

\section{综合例题}

\begin{example}
不用 \code{if}，用表达式判断一个整数 \code{n} 是否为偶数。
\end{example}
\begin{solution}
\texttt{n \% 2 == 0}。结果为 1 表示偶数，0 表示奇数。注意负数的情况，用 \code{== 0} 判断比 \code{== 1} 安全。
\end{solution}

\begin{example}
用位运算判断整数 \code{n} 的第 5 位是否为 1。
\end{example}
\begin{solution}
\texttt{(n >> 5) \& 1} 或 \texttt{n \& (1 << 5)}。前者得到 0 或 1，后者得到 0 或 32（非零即真）。
\end{solution}

\begin{example}
把秒数 \code{total} 转换成“时:分:秒”的形式。
\end{example}
\begin{solution}
思路如下：

\begin{codebox}{time.c —— 秒数转换}
#include <stdio.h>

int main(void)
{
    int total = 3725;
    int h, m, s;

    h = total / 3600;
    m = total % 3600 / 60;
    s = total % 60;

    printf("%d:%02d:%02d\n", h, m, s);
    return 0;
}
\end{codebox}

\begin{outbox}{运行结果}
1:02:05
\end{outbox}

\texttt{total \% 3600} 先去掉整小时，再除以 60 得到分钟；\texttt{total \% 60} 直接得到秒。
\end{solution}

\begin{example}
用异或运算交换两个整数（不使用临时变量）。
\end{example}
\begin{solution}
思路如下：

\begin{codebox}{xor-swap.c —— 异或交换}
#include <stdio.h>

int main(void)
{
    int a = 3;
    int b = 5;

    a = a ^ b;
    b = a ^ b;      /* 此时 b 变成原来的 a */
    a = a ^ b;      /* 此时 a 变成原来的 b */

    printf("a = %d, b = %d\n", a, b);
    return 0;
}
\end{codebox}

\begin{outbox}{运行结果}
a = 5, b = 3
\end{outbox}

\begin{pitfall}[异或交换只适合炫技]
这个技巧在面试题里很常见，但实际工程中不推荐：可读性差，而且如果 \code{a} 和 \code{b} 是同一个变量，会把值清零。老老实实用临时变量。
\end{pitfall}
\end{solution}

\begin{example}
写出判断闰年的逻辑表达式（先不考虑 \code{if} 的写法）。
\end{example}
\begin{solution}
闰年的规则是：能被 4 整除但不能被 100 整除，或者能被 400 整除。表达式为：

\begin{center}
\texttt{(year \% 4 == 0 \&\& year \% 100 != 0) || (year \% 400 == 0)}
\end{center}

下一章我们就用这个表达式写完整的判断程序。
\end{solution}

\section{本章小结}

\begin{center}
\begin{tabular}{@{}p{0.24\textwidth}p{0.66\textwidth}@{}}
\toprule
主题 & 要点 \\
\midrule
算术 & 整数除法会截断；\texttt{\%} 只能用于整数；除零危险；溢出危险 \\
赋值 & \code{=} 是赋值，\code{==} 是比较；复合赋值优先级低 \\
自增自减 & 前置先加后用，后置先用后加；同一表达式不要多次修改同一变量 \\
关系 & 结果是 0 或 1；不要连写比较；浮点不要直接用 \code{==} \\
逻辑 & 非 0 即真；\texttt{\&\&}、\code{||} 短路；德·摩根定律 \\
位运算 & 与、或、异或、取反、移位；用 \code{unsigned} 更安全 \\
条件运算符 & 二选一的简写；优先级低，注意括号 \\
类型转换 & 整型提升、有符号/无符号混算的坑 \\
求值顺序 & 只有顺序点保证顺序；避免副作用叠加 \\
\bottomrule
\end{tabular}
\end{center}

\begin{tipbox}[下一章预告]
有了关系运算和逻辑运算，我们就可以写“如果……就……”了。下一章学习 \code{if}、\code{if-else}、\code{else if} 和 \code{switch}，让程序学会“明辨是非”。
\end{tipbox}

\begin{problemset}
  \item 计算 \code{7 / 2}、\texttt{7 \% 2}、\code{-7 / 2}、\texttt{-7 \% 2} 的值，并解释。
  \item 写出判断整数 \code{n} 是奇数的表达式（要求对负数也正确）。
  \item 写出判断整数 \code{n} 能被 3 整除但不能被 5 整除的表达式。
  \item 用位运算实现：判断 \code{n} 的第 0 位是否为 1；把第 2 位置 1；把第 3 位清零。
  \item 分析表达式 \code{a = b = c + 1} 的求值顺序。
  \item 下面代码有什么问题？
\begin{codebox}{q6.c —— 找问题}
#include <stdio.h>

int main(void)
{
    int i = 0;
    int a;
    a = i++ + ++i;
    printf("%d %d\n", a, i);
    return 0;
}
\end{codebox}
  \item 写一个表达式，判断字符 \code{c} 是不是小写字母。
  \item 解释为什么 \code{-1 < 1U} 的结果是假。
  \item 用条件运算符求三个整数中的最大值。
  \item 思考题：为什么 \code{x << 2} 不一定等于 \code{x * 4}？举一个反例。
\end{problemset}

\chapter{分支：明辨是非}

\begin{introduction}
  \item \texttt{if}：满足条件才执行
  \item \texttt{if-else}：二选一
  \item \texttt{else if}：多路分支与分段函数
  \item 嵌套分支与“就近匹配”的 \texttt{else}
  \item \texttt{switch}：按整数常量分情况处理
  \item 分支的常见错误与综合例题
\end{introduction}

\chaptermotto{「计算机的“聪明”，来自你替它写好的每一个“如果”。」}

顺序结构只能一条路走到黑。可现实中的程序要会看情况：温度过高就报警，按键按下就启动，成绩及格就通过，除数不为零才做除法。这就需要\textbf{分支结构}。

C 语言的分支工具只有两类：\code{if} 语句和 \code{switch} 语句。这一章把它们的语法、套路、坑和例题一次讲透。

\section{为什么需要分支}

先看一个没有分支的程序：

\begin{codebox}{no-branch.c —— 只会往前走的程序}
#include <stdio.h>

int main(void)
{
    int score = 85;

    printf("你的成绩是 %d 分\n", score);
    printf("恭喜你，通过了！\n");
    return 0;
}
\end{codebox}

无论 \code{score} 是多少，它都会打印“恭喜你，通过了”。这显然不合理。我们希望程序这样思考：

\begin{center}
\kaishu 如果成绩大于等于 60，就打印“通过”；否则打印“未通过”。
\end{center}

用流程图表示：

\begin{figure}[htbp]
  \centering
  \begin{tikzpicture}[x=1cm,y=1cm,>=Stealth,node distance=7mm]
    \node[terminator] (start) {开始};
    \node[io,below=of start] (in) {读入成绩 score};
    \node[decision,below=9mm of in] (cond) {score $\ge$ 60 ?};
    \node[process,below left=9mm and 4mm of cond] (pass) {输出“通过”};
    \node[process,below right=9mm and 4mm of cond] (fail) {输出“未通过”};
    \node[terminator,below=18mm of cond] (end) {结束};

    \draw[flow] (start) -- (in);
    \draw[flow] (in) -- (cond);
    \draw[flow] (cond) -- node[lbl,left]{是} (pass);
    \draw[flow] (cond) -- node[lbl,right]{否} (fail);
    \draw[flow] (pass) |- (end);
    \draw[flow] (fail) |- (end);
  \end{tikzpicture}
  \caption{带分支的流程图：条件成立走一条路，不成立走另一条路}
  \label{fig:if-flow}
\end{figure}

流程图里的\textbf{菱形}表示判断，有“是/否”两条出口；\textbf{矩形}表示处理；\textbf{圆角矩形}表示开始/结束；\textbf{平行四边形}表示输入输出。后面我们还会用这套符号。

\section{\texttt{if} 语句}

\subsection{基本语法}

\begin{center}
\fbox{\parbox{0.72\textwidth}{
\ttfamily if (条件) \{ \\
\phantom{xxxx}语句; \\
\ttfamily \}
}}
\end{center}

执行逻辑：计算括号里的条件，如果结果为\textbf{真}（非 0），就执行花括号里的语句；否则跳过。

\begin{codebox}{if-basic.c —— 最简单的 if}
#include <stdio.h>

int main(void)
{
    int score = 85;

    if (score >= 60) {
        printf("恭喜你，通过了！\n");
    }

    printf("程序结束\n");
    return 0;
}
\end{codebox}

\begin{outbox}{运行结果}
恭喜你，通过了！
程序结束
\end{outbox}

如果把 \code{score} 改成 50，第一行就不会打印，但“程序结束”仍然会打印——因为 \code{if} 只控制它后面那一组花括号。

\subsection{条件可以是任意表达式}

\code{if} 只关心“是不是 0”：非 0 为真，0 为假。所以下面这些写法都合法：

\begin{codebox}{if-truth.c —— 条件的真真假假}
#include <stdio.h>

int main(void)
{
    int a = 5;

    if (a) {
        printf("a 非零，条件为真\n");
    }

    if (a - 5) {
        printf("这一行不会打印\n");
    }

    if (a > 0 && a < 10) {
        printf("a 在 0 到 10 之间\n");
    }

    if (!a) {
        printf("这一行也不会打印\n");
    }

    return 0;
}
\end{codebox}

\begin{outbox}{运行结果}
a 非零，条件为真
a 在 0 到 10 之间
\end{outbox}

\begin{tipbox}[条件写清楚，别炫技]
\code{if (a)} 虽然合法，但读者要猜“\code{a} 是什么含义”。更好的写法是 \code{if (a != 0)} 或 \code{if (count > 0)}，把意图写出来。
\end{tipbox}

\subsection{单条语句可以省略花括号}

如果 \code{if} 后面只有一条语句，花括号可以省略：

\begin{codebox}{if-nobrace.c —— 省略花括号}
#include <stdio.h>

int main(void)
{
    int a = 10;

    if (a > 0)
        printf("a 是正数\n");

    return 0;
}
\end{codebox}

\begin{pitfall}[省略花括号是 bug 的温床]
\begin{codebox}{if-bug.c —— 少写花括号的后果}
#include <stdio.h>

int main(void)
{
    int a = -5;

    if (a > 0)
        printf("a 是正数\n");
        printf("这一行无论 a 正负都会打印\n");   /* 缩进骗了你 */

    return 0;
}
\end{codebox}
第二行看起来在 \code{if} 里面，实际上不在。\textbf{本书约定：即使只有一条语句，也写花括号}。多打两个字符，少一个通宵调试。
\end{pitfall}

\section{\texttt{if-else} 语句}

\subsection{二选一}

\begin{center}
\fbox{\parbox{0.78\textwidth}{
\ttfamily if (条件) \{ \\
\phantom{xxxx}条件为真时执行; \\
\ttfamily \} else \{ \\
\phantom{xxxx}条件为假时执行; \\
\ttfamily \}
}}
\end{center}

\begin{figure}[htbp]
  \centering
  \begin{tikzpicture}[x=1cm,y=1cm,>=Stealth,node distance=7mm]
    \node[terminator] (start) {开始};
    \node[decision,below=9mm of start] (cond) {条件 ?};
    \node[process,below left=9mm and 8mm of cond] (yes) {语句组 A};
    \node[process,below right=9mm and 8mm of cond] (no) {语句组 B};
    \node[terminator,below=20mm of cond] (end) {结束};

    \draw[flow] (start) -- (cond);
    \draw[flow] (cond) -- node[lbl,left]{真} (yes);
    \draw[flow] (cond) -- node[lbl,right]{假} (no);
    \draw[flow] (yes) |- (end);
    \draw[flow] (no) |- (end);
  \end{tikzpicture}
  \caption{\texttt{if-else} 的执行路径}
  \label{fig:ifelse-flow}
\end{figure}

\begin{codebox}{if-else.c —— 判断正负数}
#include <stdio.h>

int main(void)
{
    int n;

    printf("请输入一个整数：");
    scanf("%d", &n);

    if (n > 0) {
        printf("%d 是正数\n", n);
    } else {
        printf("%d 不是正数\n", n);
    }

    return 0;
}
\end{codebox}

\begin{outbox}{运行示例}
请输入一个整数：-7
-7 不是正数
\end{outbox}

注意“不是正数”包括 0 和负数。如果想区分 0、正数、负数，就需要下一节的多分支。

\subsection{例：求绝对值}

\begin{example}
输入一个整数，输出它的绝对值。
\end{example}
\begin{solution}
思路如下：

\begin{codebox}{abs.c —— 求绝对值}
#include <stdio.h>

int main(void)
{
    int n;
    int abs_n;

    printf("请输入一个整数：");
    scanf("%d", &n);

    if (n < 0) {
        abs_n = -n;
    } else {
        abs_n = n;
    }

    printf("|%d| = %d\n", n, abs_n);
    return 0;
}
\end{codebox}

\begin{outbox}{运行示例}
请输入一个整数：-9
|-9| = 9
\end{outbox}

也可以用条件运算符写成一行：\code{abs_n = (n < 0) ? -n : n;}。
\end{solution}

\section{\texttt{else if}：多路分支}

\subsection{语法与执行顺序}

当情况不止两种时，可以用 \code{else if} 串成一条“阶梯”：

\begin{center}
\fbox{\parbox{0.86\textwidth}{
\ttfamily if (条件1) \{ ... \} \\
\ttfamily else if (条件2) \{ ... \} \\
\ttfamily else if (条件3) \{ ... \} \\
\ttfamily else \{ ... \}
}}
\end{center}

执行逻辑：从上往下依次判断，\textbf{一旦某个条件成立，就执行对应的语句组，然后跳出整个阶梯}；如果所有条件都不成立，执行最后的 \code{else}。

\begin{figure}[htbp]
  \centering
  \begin{tikzpicture}[x=1cm,y=1cm,>=Stealth,node distance=5mm]
    \node[decision] (c1) at (0,0) {条件 1 ?};
    \node[process,right=14mm of c1] (b1) {语句组 1};
    \node[decision,below=8mm of c1] (c2) {条件 2 ?};
    \node[process,right=14mm of c2] (b2) {语句组 2};
    \node[decision,below=8mm of c2] (c3) {条件 3 ?};
    \node[process,right=14mm of c3] (b3) {语句组 3};
    \node[process,below=8mm of c3] (be) {else 语句组};

    \draw[flow] (c1) -- node[lbl,above]{真} (b1);
    \draw[flow] (c1) -- node[lbl,right]{假} (c2);
    \draw[flow] (c2) -- node[lbl,above]{真} (b2);
    \draw[flow] (c2) -- node[lbl,right]{假} (c3);
    \draw[flow] (c3) -- node[lbl,above]{真} (b3);
    \draw[flow] (c3) -- node[lbl,right]{假} (be);
    \node[lbl,right=2mm of b1] {执行完即跳出};
    \node[lbl,right=2mm of b2] {执行完即跳出};
    \node[lbl,right=2mm of b3] {执行完即跳出};
  \end{tikzpicture}
  \caption{\texttt{else if} 阶梯：从上到下，命中即停}
  \label{fig:elseif}
\end{figure}

\subsection{例：成绩等级}

\begin{example}
输入一个 0\textasciitilde100 的成绩，输出等级：90 及以上为 A，80\textasciitilde89 为 B，70\textasciitilde79 为 C，60\textasciitilde69 为 D，60 以下为 E。非法成绩给出提示。
\end{example}
\begin{solution}
思路如下：

\begin{codebox}{grade.c —— 成绩等级}
#include <stdio.h>

int main(void)
{
    int score;

    printf("请输入成绩（0-100）：");
    scanf("%d", &score);

    if (score < 0 || score > 100) {
        printf("成绩不合法\n");
    } else if (score >= 90) {
        printf("等级 A\n");
    } else if (score >= 80) {
        printf("等级 B\n");
    } else if (score >= 70) {
        printf("等级 C\n");
    } else if (score >= 60) {
        printf("等级 D\n");
    } else {
        printf("等级 E\n");
    }

    return 0;
}
\end{codebox}

\begin{outbox}{运行示例}
请输入成绩（0-100）：85
等级 B
\end{outbox}

注意判断顺序：先排除非法值，再从高到低判断。如果把 \code{score >= 60} 放在最前面，85 会先命中 60 这一档，直接输出 D，后面的判断永远不会执行。
\end{solution}

\begin{pitfall}[\texttt{else if} 的顺序很重要]
多分支是“命中即停”的。范围判断要从宽到窄或从窄到宽，保证先命中的是更具体的条件。写完后用边界值（59、60、89、90）测试一遍。
\end{pitfall}

\subsection{例：分段函数}

\begin{example}
输入 $x$，计算
\[
y = \begin{cases}
-1, & x < 0,\\
0, & x = 0,\\
1, & x > 0.
\end{cases}
\]
\end{example}
\begin{solution}
思路如下：

\begin{codebox}{sign.c —— 符号函数}
#include <stdio.h>

int main(void)
{
    double x;
    int y;

    printf("请输入 x：");
    scanf("%lf", &x);

    if (x < 0) {
        y = -1;
    } else if (x == 0) {
        y = 0;
    } else {
        y = 1;
    }

    printf("y = %d\n", y);
    return 0;
}
\end{codebox}

\begin{outbox}{运行示例}
请输入 x：-2.5
y = -1
\end{outbox}

\begin{pitfall}[浮点数不要直接用 \texttt{==} 判断]
\code{x == 0} 对精确输入的 0 可以工作，但如果是计算出来的值（例如 \code{0.1 + 0.2 - 0.3}），它可能是一个极小的非零数，判断就会失败。更稳妥的写法是 \code{fabs(x) < 1e-9}（需要 \file{math.h}）。
\end{pitfall}
\end{solution}

\section{嵌套分支}

\subsection{一个分支里再套一个分支}

分支可以嵌套。例如判断一个年份是否是闰年，再判断这一年的二月有多少天：

\begin{codebox}{nested.c —— 嵌套 if}
#include <stdio.h>

int main(void)
{
    int year;
    int is_leap;

    printf("请输入年份：");
    scanf("%d", &year);

    if (year % 4 == 0) {
        if (year % 100 != 0) {
            is_leap = 1;
        } else {
            if (year % 400 == 0) {
                is_leap = 1;
            } else {
                is_leap = 0;
            }
        }
    } else {
        is_leap = 0;
    }

    if (is_leap) {
        printf("%d 年是闰年，二月有 29 天\n", year);
    } else {
        printf("%d 年不是闰年，二月有 28 天\n", year);
    }

    return 0;
}
\end{codebox}

嵌套让逻辑更接近人的思考方式，但层数太多会难以阅读。一般建议嵌套不超过三层；超过时，考虑用“提前返回”或拆成函数（第 9 章之后会讲）。

\subsection{悬空的 \texttt{else}}

看下面这段代码，\code{else} 到底和哪个 \code{if} 配对？

\begin{codebox}{dangling.c —— 悬空的 else}
#include <stdio.h>

int main(void)
{
    int a = 1;
    int b = -1;

    if (a > 0)
        if (b > 0)
            printf("a、b 都是正数\n");
    else
        printf("b 不是正数\n");

    return 0;
}
\end{codebox}

\begin{outbox}{运行结果}
b 不是正数
\end{outbox}

C 语言的规则是：\textbf{\code{else} 与离它最近的、还没有配对的 \code{if} 配对}。所以这里的 \code{else} 属于内层 \code{if (b > 0)}，而不是外层。缩进在这件事上完全不起作用——编译器只看语法。

\begin{pitfall}[用花括号消灭歧义]
上面的代码如果本意是“\code{a} 不大于 0 时打印”，应该写成：
\begin{codebox}{dangling-fixed.c —— 加花括号}
#include <stdio.h>

int main(void)
{
    int a = 1;
    int b = -1;

    if (a > 0) {
        if (b > 0) {
            printf("a、b 都是正数\n");
        }
    } else {
        printf("a 不是正数\n");
    }

    return 0;
}
\end{codebox}
\textbf{只要 \code{if} 或 \code{else} 后面有嵌套，就一定加花括号}。
\end{pitfall}

\subsection{例：三个数中的最大值}

\begin{example}
输入三个整数，输出最大值。
\end{example}
\begin{solution}
思路如下：

\begin{codebox}{max3.c —— 三个数的最大值}
#include <stdio.h>

int main(void)
{
    int a, b, c, max;

    printf("请输入三个整数：");
    scanf("%d%d%d", &a, &b, &c);

    if (a > b) {
        if (a > c) {
            max = a;
        } else {
            max = c;
        }
    } else {
        if (b > c) {
            max = b;
        } else {
            max = c;
        }
    }

    printf("最大值是 %d\n", max);
    return 0;
}
\end{codebox}

也可以先令 \code{max = a}，再依次比较：
\begin{codebox}{max3-simple.c —— 更简洁的写法}
#include <stdio.h>

int main(void)
{
    int a, b, c, max;

    printf("请输入三个整数：");
    scanf("%d%d%d", &a, &b, &c);

    max = a;
    if (b > max) {
        max = b;
    }
    if (c > max) {
        max = c;
    }

    printf("最大值是 %d\n", max);
    return 0;
}
\end{codebox}
后一种写法更容易扩展到 $n$ 个数，也更不容易出错。第 5 章学循环后，我们会用循环处理任意多个数。
\end{solution}

\section{\texttt{switch} 语句}

\subsection{基本语法}

当分支依据是“一个整数表达式的若干个确定取值”时，\code{switch} 比一长串 \code{else if} 更清楚。

\begin{center}
\fbox{\parbox{0.86\textwidth}{
\ttfamily switch (表达式) \{ \\
\phantom{xxxx}\ttfamily case 常量1: 语句; break; \\
\phantom{xxxx}\ttfamily case 常量2: 语句; break; \\
\phantom{xxxx}\ttfamily default: 语句; break; \\
\ttfamily \}
}}
\end{center}

\begin{itemize}
  \item \code{switch} 后面的表达式必须是\textbf{整数类型}（\code{int}、\code{char}、\code{enum} 等），不能是 \code{double} 或字符串；
  \item \code{case} 后面必须是\textbf{编译期常量}，不能是变量；
  \item \code{break} 表示“跳出 \code{switch}”，没有 \code{break} 会继续往下执行（称为“贯穿”，fall-through）；
  \item \code{default} 处理所有没有匹配到的情况，位置可以放在任意处，但习惯上放最后。
\end{itemize}

\begin{codebox}{switch-basic.c —— 根据数字输出星期}
#include <stdio.h>

int main(void)
{
    int day;

    printf("请输入 1-7：");
    scanf("%d", &day);

    switch (day) {
    case 1:
        printf("星期一\n");
        break;
    case 2:
        printf("星期二\n");
        break;
    case 3:
        printf("星期三\n");
        break;
    case 4:
        printf("星期四\n");
        break;
    case 5:
        printf("星期五\n");
        break;
    case 6:
        printf("星期六\n");
        break;
    case 7:
        printf("星期日\n");
        break;
    default:
        printf("输入不合法\n");
        break;
    }

    return 0;
}
\end{codebox}

\begin{outbox}{运行示例}
请输入 1-7：3
星期三
\end{outbox}

\subsection{贯穿（fall-through）}

如果某个 \code{case} 后面没有 \code{break}，程序会继续执行下一个 \code{case} 的语句，直到遇到 \code{break} 或 \code{switch} 结束。

\begin{figure}[htbp]
  \centering
  \begin{tikzpicture}[x=1cm,y=1cm,>=Stealth,node distance=6mm]
    \node[process,minimum width=30mm] (c1) at (0,0) {case 1};
    \node[process,minimum width=30mm] (c2) at (0,-1.2) {case 2};
    \node[process,minimum width=30mm] (c3) at (0,-2.4) {case 3};
    \node[process,minimum width=30mm] (c4) at (0,-3.6) {case 4};
    \draw[flow] (c1) -- node[lbl,right]{没有 break} (c2);
    \draw[flow] (c2) -- node[lbl,right]{没有 break} (c3);
    \draw[flow] (c3) -- node[lbl,right]{遇到 break} (c4);
    \node[lbl,right=8mm of c4,align=left] {跳出 switch};
  \end{tikzpicture}
  \caption{没有 \texttt{break} 时会“贯穿”到下一个 \texttt{case}}
  \label{fig:fallthrough}
\end{figure}

\begin{pitfall}[漏写 \texttt{break} 是最常见的 \texttt{switch} 错误]
\begin{codebox}{switch-bug.c —— 漏写 break}
#include <stdio.h>

int main(void)
{
    int n = 1;

    switch (n) {
    case 1:
        printf("one\n");
    case 2:
        printf("two\n");
    case 3:
        printf("three\n");
        break;
    default:
        printf("other\n");
        break;
    }
    return 0;
}
\end{codebox}
输出是三行：\code{one}、\code{two}、\code{three}，而不是只有 \code{one}。
\end{pitfall}

\subsection{故意利用贯穿}

有些场景下贯穿反而能简化代码，例如判断某个月有多少天：

\begin{codebox}{days.c —— 月份天数}
#include <stdio.h>

int main(void)
{
    int month;
    int year;
    int days;

    printf("请输入年份和月份：");
    scanf("%d%d", &year, &month);

    switch (month) {
    case 1: case 3: case 5: case 7:
    case 8: case 10: case 12:
        days = 31;
        break;
    case 4: case 6: case 9: case 11:
        days = 30;
        break;
    case 2:
        if ((year % 4 == 0 && year % 100 != 0) || (year % 400 == 0)) {
            days = 29;
        } else {
            days = 28;
        }
        break;
    default:
        printf("月份不合法\n");
        return 1;
    }

    printf("%d 年 %d 月有 %d 天\n", year, month, days);
    return 0;
}
\end{codebox}

\begin{outbox}{运行示例}
请输入年份和月份：2024 2
2024 年 2 月有 29 天
\end{outbox}

把多个 \code{case} 写在一起、共享同一段处理代码，是贯穿的正当用法。但要在代码里写注释说明“这是故意的”，避免后来的人误以为你漏了 \code{break}。

\subsection{例：简易计算器}

\begin{example}
输入两个数和四则运算符（\code{+}、\code{-}、\code{*}、\code{/}），输出运算结果。除法要判断除数是否为 0。
\end{example}
\begin{solution}
思路如下：

\begin{codebox}{calc.c —— 简易计算器}
#include <stdio.h>

int main(void)
{
    double a, b;
    char op;

    printf("请输入表达式（如 3.5 + 2）：");
    scanf("%lf %c %lf", &a, &op, &b);

    switch (op) {
    case '+':
        printf("%.2f\n", a + b);
        break;
    case '-':
        printf("%.2f\n", a - b);
        break;
    case '*':
        printf("%.2f\n", a * b);
        break;
    case '/':
        if (b == 0.0) {
            printf("除数不能为 0\n");
        } else {
            printf("%.2f\n", a / b);
        }
        break;
    default:
        printf("不支持的运算符：%c\n", op);
        break;
    }

    return 0;
}
\end{codebox}

\begin{outbox}{运行示例}
请输入表达式（如 3.5 + 2）：3.5 * 2
7.00
\end{outbox}

格式串 \texttt{"\%lf \%c \%lf"} 中，\texttt{\%c} 前面的空格很重要：它会跳过前面数字后面的空白字符（包括空格和换行），否则 \texttt{\%c} 可能读到一个空格。
\end{solution}

\section{分支的常见错误}

\begin{center}
\begin{tabular}{@{}p{0.36\textwidth}p{0.54\textwidth}@{}}
\toprule
错误写法 & 问题 \\
\midrule
\code{if (a = 5)} & 赋值，不是比较；条件永远为真 \\
\code{if (a == 5);} & 多了一个分号，\code{if} 后面变成了空语句 \\
\code{if (a > 0)} 后接多条语句但不加花括号 & 只有第一条受 \code{if} 控制 \\
\code{if (0 < x < 10)} & 连写比较，结果永远为真 \\
\texttt{if (a \& b)} & 按位与，不是逻辑与；应写 \texttt{\&\&} \\
\code{if (x == 0.1)} & 浮点数直接比较，可能不相等 \\
\code{switch (x)} 且 \code{x} 是 \code{double} & \code{switch} 只接受整数 \\
\code{case 1.5:} & \code{case} 必须是整数常量 \\
\code{default} 后忘记 \code{break} & 通常无害，但不规范 \\
\bottomrule
\end{tabular}
\end{center}

\begin{pitfall}[\texttt{if} 后面的分号]
\begin{codebox}{if-semicolon.c —— 致命的空语句}
if (a > 0);          /* 分号让 if 提前结束 */
{
    printf("a 是正数\n");   /* 这段代码永远执行 */
}
\end{codebox}
编译器通常不会报错，甚至不会警告。看到 \code{if (...)} 后面紧跟分号，立刻警觉。
\end{pitfall}

\begin{tipbox}[条件里的“范围”怎么写]
“\code{x} 在 0 到 10 之间”要写成 \texttt{x > 0 \&\& x < 10}；“\code{x} 不在 0 到 10 之间”写成 \code{x <= 0 || x >= 10}。注意逻辑或不是“与”的简单取反，写完后用边界值测试。
\end{tipbox}

\section{综合例题}

\subsection{例：闰年判断}

\begin{example}
输入年份，判断是否为闰年。
\end{example}
\begin{solution}
规则：能被 4 整除且不能被 100 整除，或者能被 400 整除。

\begin{codebox}{leap.c —— 闰年判断}
#include <stdio.h>

int main(void)
{
    int year;

    printf("请输入年份：");
    scanf("%d", &year);

    if ((year % 4 == 0 && year % 100 != 0) || (year % 400 == 0)) {
        printf("%d 年是闰年\n", year);
    } else {
        printf("%d 年不是闰年\n", year);
    }

    return 0;
}
\end{codebox}

\begin{outbox}{运行示例}
请输入年份：2000
2000 年是闰年
\end{outbox}

测试用例建议覆盖：1900（不是）、2000（是）、2024（是）、2023（不是）。
\end{solution}

\subsection{例：一元二次方程}

\begin{example}
输入 $a$、$b$、$c$，求方程 $ax^2 + bx + c = 0$ 的实根。要考虑 $a = 0$、判别式大于 0、等于 0、小于 0 四种情况。
\end{example}
\begin{solution}
思路如下：

\begin{codebox}{quadratic.c —— 一元二次方程}
#include <stdio.h>
#include <math.h>

int main(void)
{
    double a, b, c;
    double delta;
    double x1, x2;

    printf("请输入 a b c：");
    scanf("%lf%lf%lf", &a, &b, &c);

    if (a == 0.0) {
        if (b == 0.0) {
            printf("不是二次方程\n");
        } else {
            printf("x = %.2f\n", -c / b);
        }
    } else {
        delta = b * b - 4 * a * c;
        if (delta > 0.0) {
            x1 = (-b + sqrt(delta)) / (2 * a);
            x2 = (-b - sqrt(delta)) / (2 * a);
            printf("两个实根：x1 = %.2f, x2 = %.2f\n", x1, x2);
        } else if (delta == 0.0) {
            x1 = -b / (2 * a);
            printf("两个相等实根：x = %.2f\n", x1);
        } else {
            printf("没有实根\n");
        }
    }

    return 0;
}
\end{codebox}

\begin{outbox}{运行示例}
请输入 a b c：1 -3 2
两个实根：x1 = 2.00, x2 = 1.00
\end{outbox}

\begin{note}
在 Linux 上用 GCC 编译含 \file{math.h} 的程序时，需要在命令末尾加 \code{-lm} 链接数学库：\code{gcc quadratic.c -o quadratic -lm}。
\end{note}
\end{solution}

\subsection{例：三角形判断}

\begin{example}
输入三条边长，判断能否构成三角形；如果能，再判断是等边、等腰还是普通三角形。
\end{example}
\begin{solution}
思路如下：

\begin{codebox}{triangle.c —— 三角形判断}
#include <stdio.h>

int main(void)
{
    double a, b, c;

    printf("请输入三条边长：");
    scanf("%lf%lf%lf", &a, &b, &c);

    if (a + b > c && a + c > b && b + c > a) {
        if (a == b && b == c) {
            printf("等边三角形\n");
        } else if (a == b || a == c || b == c) {
            printf("等腰三角形\n");
        } else {
            printf("普通三角形\n");
        }
    } else {
        printf("不能构成三角形\n");
    }

    return 0;
}
\end{codebox}

\begin{outbox}{运行示例}
请输入三条边长：3 3 5
等腰三角形
\end{outbox}

\begin{pitfall}[浮点数判断相等的隐患]
这里用 \code{==} 判断边长相等，对整数输入没问题；如果边长是计算出来的小数，应改用 \code{fabs(a - b) < 1e-9}。
\end{pitfall}
\end{solution}

\subsection{例：出租车计费}

\begin{example}
某城市出租车计费规则：3 公里以内（含）收费 13 元；超过 3 公里的部分，每公里 2.5 元；超过 10 公里的部分，每公里加收 50\% 空驶费（即 3.75 元/公里）。输入行驶公里数，输出费用。
\end{example}
\begin{solution}
思路如下：

\begin{codebox}{taxi.c —— 分段计费}
#include <stdio.h>

int main(void)
{
    double km;
    double fee;

    printf("请输入公里数：");
    scanf("%lf", &km);

    if (km <= 3.0) {
        fee = 13.0;
    } else if (km <= 10.0) {
        fee = 13.0 + (km - 3.0) * 2.5;
    } else {
        fee = 13.0 + 7.0 * 2.5 + (km - 10.0) * 2.5 * 1.5;
    }

    printf("费用 = %.2f 元\n", fee);
    return 0;
}
\end{codebox}

\begin{outbox}{运行示例}
请输入公里数：12
费用 = 38.00
\end{outbox}

分段计费是 \code{else if} 的典型场景：每一段都只计算“超出部分”，千万不要重复计算前面的里程。
\end{solution}

\section{拔高：让分支更好读}

\subsection{卫语句：先处理异常，再处理正常}

当嵌套很深时，可以把“不满足条件就提前结束”的情况放在前面：

\begin{codebox}{guard.c —— 卫语句}
#include <stdio.h>

int main(void)
{
    int score;

    printf("请输入成绩：");
    scanf("%d", &score);

    if (score < 0 || score > 100) {
        printf("成绩不合法\n");
        return 1;
    }

    if (score >= 60) {
        printf("通过\n");
    } else {
        printf("未通过\n");
    }

    return 0;
}
\end{codebox}

先排除非法输入，后面的逻辑就不用再考虑它，代码更扁平。

\subsection{表驱动：用数组代替一长串 \texttt{if}}

如果分支是“输入一个编号，输出一个固定结果”，可以把结果放进数组，用下标直接取。第 6 章学完数组后，你会看到这种写法有多简洁。

\subsection{分支预测}

CPU 遇到 \code{if} 时会“猜”走哪条路，猜对了就快，猜错了要清空流水线。所以\textbf{把最可能发生的情况放在前面}，有时能提升性能。但对初学者来说，先保证正确和可读，性能是后话。

\section{本章小结}

\begin{center}
\begin{tabular}{@{}p{0.22\textwidth}p{0.68\textwidth}@{}}
\toprule
主题 & 要点 \\
\midrule
\code{if} & 条件非 0 则执行；建议总是加花括号 \\
\code{if-else} & 二选一；两个分支必有一个执行 \\
\code{else if} & 多路分支，从上到下命中即停；注意判断顺序 \\
嵌套 & 分支里套分支；\code{else} 与最近未配对的 \code{if} 匹配 \\
\code{switch} & 整数表达式 + 常量 \code{case}；\code{break} 跳出；可故意贯穿 \\
常见错误 & \code{=} 与 \code{==}、多余分号、漏花括号、连写比较 \\
综合应用 & 闰年、方程、三角形、计费、计算器 \\
\bottomrule
\end{tabular}
\end{center}

\begin{tipbox}[下一章预告]
分支让程序会“选”，但还不会“重复”。如果要计算 1 到 100 的和，你总不能写一百行加法。下一章学习循环，让程序学会“重复再重复”。
\end{tipbox}

\begin{problemset}
  \item 输入一个整数，判断它是正数、负数还是零。
  \item 输入一个年份，判断是否为闰年（分别用 \code{if} 和逻辑表达式两种写法）。
  \item 输入一个字符，判断它是大写字母、小写字母、数字还是其他字符。
  \item 输入三个整数，按从小到大的顺序输出。
  \item 输入学生成绩，输出等级（A/B/C/D/E），并处理非法输入。
  \item 用 \code{switch} 实现：输入月份，输出该月所属的季节（春/夏/秋/冬）。
  \item 输入一个一元二次方程的系数，输出实根（考虑各种情况）。
  \item 输入三条边，判断能否构成三角形，并判断类型。
  \item 解释 \code{if (a = 5)} 和 \code{if (a == 5)} 的区别。
  \item 思考题：为什么 \code{switch} 的 \code{case} 后面不能是变量？如果一定要根据变量分支，应该怎么写？
\end{problemset}

\chapter{循环：重复再重复}

\begin{introduction}
  \item 为什么需要循环：从“写一百行”到“写三行”
  \item \texttt{while}、\texttt{do-while}、\texttt{for} 三种循环
  \item C89 的循环变量声明规矩
  \item \texttt{break} 与 \texttt{continue}
  \item 嵌套循环与经典算法：累加、素数、斐波那契、辗转相除
  \item 循环的写法、边界与常见错误
\end{introduction}

\chaptermotto{「聪明人把重复交给机器，把思考留给自己。」}

如果要计算 $1+2+3+\cdots+100$，你可以写一百行 \code{sum = sum + i;}，但没人会这么做。循环就是为“重复做同一件事”而生的。

本章会学三种循环：\code{while}、\code{do-while}、\code{for}。它们能力相同，只是适合的场景不同。学完后，你会用循环解决求和、阶乘、素数、斐波那契、最大公约数等一大批问题。

\section{为什么需要循环}

先看一个笨办法：

\begin{codebox}{sum100-hard.c —— 笨办法：写一百行}
#include <stdio.h>

int main(void)
{
    int sum = 0;

    sum = sum + 1;
    sum = sum + 2;
    sum = sum + 3;
    /* ... 中间省略 96 行 ... */
    sum = sum + 100;

    printf("sum = %d\n", sum);
    return 0;
}
\end{codebox}

这段代码有 100 行几乎一样的语句。如果要求和到 10000，就得写一万行；如果要求到 1 亿，根本没法写。更糟糕的是，只要需求变一下（比如只加偶数），就要改一百处。

循环的思路是：把“重复的动作”写一次，让机器重复执行。

\begin{codebox}{sum100.c —— 用循环求和}
#include <stdio.h>

int main(void)
{
    int i;
    int sum = 0;

    for (i = 1; i <= 100; i++) {
        sum = sum + i;
    }

    printf("1 + 2 + ... + 100 = %d\n", sum);
    return 0;
}
\end{codebox}

\begin{outbox}{运行结果}
1 + 2 + ... + 100 = 5050
\end{outbox}

三行循环，替代了一百行。这就是循环的威力。

\begin{figure}[htbp]
  \centering
  \begin{tikzpicture}[x=1cm,y=1cm,>=Stealth,node distance=7mm]
    \node[terminator] (start) {开始};
    \node[process,below=of start] (init) {i = 1, sum = 0};
    \node[decision,below=9mm of init] (cond) {i $\le$ 100 ?};
    \node[process,below=9mm of cond] (body) {sum = sum + i};
    \node[process,below=7mm of body] (inc) {i = i + 1};
    \node[terminator,right=24mm of body] (end) {结束};

    \draw[flow] (start) -- (init);
    \draw[flow] (init) -- (cond);
    \draw[flow] (cond) -- node[lbl,left]{是} (body);
    \draw[flow] (cond) -- node[lbl,above]{否} (end);
    \draw[flow] (body) -- (inc);
    \draw[flow] (inc.east) -- ++(6mm,0) |- (cond.east);
  \end{tikzpicture}
  \caption{累加循环的流程图：条件成立就重复“加—自增—再判断”}
  \label{fig:sum-loop}
\end{figure}

\section{\texttt{while} 循环}

\subsection{基本语法}

\begin{center}
\fbox{\parbox{0.7\textwidth}{
\ttfamily while (条件) \{ \\
\phantom{xxxx}循环体; \\
\ttfamily \}
}}
\end{center}

执行逻辑：先判断条件，为真就执行循环体，执行完再判断条件……直到条件为假，循环结束。

\begin{codebox}{while.c —— while 倒计时}
#include <stdio.h>

int main(void)
{
    int i = 5;

    while (i > 0) {
        printf("%d\n", i);
        i--;            /* 千万别忘了改变循环变量 */
    }

    printf("发射！\n");
    return 0;
}
\end{codebox}

\begin{outbox}{运行结果}
5
4
3
2
1
发射！
\end{outbox}

\begin{pitfall}[死循环]
如果循环体里没有让条件最终变假的语句，循环会永远执行下去：
\begin{codebox}{infinite.c —— 死循环}
#include <stdio.h>

int main(void)
{
    int i = 1;

    while (i <= 5) {
        printf("%d\n", i);
        /* 忘了 i++，条件永远为真 */
    }
    return 0;
}
\end{codebox}
运行后终端会无限刷 1，只能按 \code{Ctrl+C} 强制结束。
\end{pitfall}

\subsection{用 \texttt{while} 实现累加}

\begin{codebox}{while-sum.c —— while 求和}
#include <stdio.h>

int main(void)
{
    int i = 1;
    int sum = 0;

    while (i <= 100) {
        sum = sum + i;
        i++;
    }

    printf("sum = %d\n", sum);
    return 0;
}
\end{codebox}

\begin{outbox}{运行结果}
sum = 5050
\end{outbox}

写 \code{while} 循环时，按“四步法”检查：

\begin{enumerate}
  \item \textbf{初始化}：循环变量从几开始？（\code{i = 1}）
  \item \textbf{条件}：什么时候继续？（\code{i <= 100}）
  \item \textbf{循环体}：每次重复做什么？（\code{sum += i}）
  \item \textbf{更新}：怎么让循环走向结束？（\code{i++}）
\end{enumerate}

少任何一步，循环都可能不工作或停不下来。

\subsection{循环次数不确定时用 \texttt{while}}

\code{while} 最适合“不知道要循环多少次，只知道什么时候停”的场景。例如不断读入数字，直到读到 0 为止：

\begin{codebox}{read-until-zero.c —— 读到 0 为止}
#include <stdio.h>

int main(void)
{
    int n;
    int sum = 0;

    printf("请输入整数（输入 0 结束）：");

    scanf("%d", &n);
    while (n != 0) {
        sum = sum + n;
        scanf("%d", &n);
    }

    printf("总和 = %d\n", sum);
    return 0;
}
\end{codebox}

\begin{outbox}{运行示例}
请输入整数（输入 0 结束）：3 5 7 0
总和 = 15
\end{outbox}

注意 \code{scanf} 出现了两次：循环前读一次，循环体里再读一次。这是“读取—判断—处理—再读取”的经典模式。

\section{\texttt{for} 循环}

\subsection{基本语法}

\begin{center}
\fbox{\parbox{0.86\textwidth}{
\ttfamily for (初始化; 条件; 更新) \{ \\
\phantom{xxxx}循环体; \\
\ttfamily \}
}}
\end{center}

\code{for} 把 \code{while} 的四步压缩到一行，最适合“次数确定”的循环。

\begin{codebox}{for.c —— for 求和}
#include <stdio.h>

int main(void)
{
    int i;
    int sum = 0;

    for (i = 1; i <= 100; i++) {
        sum = sum + i;
    }

    printf("sum = %d\n", sum);
    return 0;
}
\end{codebox}

\subsection{执行顺序}

\code{for} 的执行顺序是：

\begin{enumerate}
  \item 执行\textbf{初始化}（只执行一次）；
  \item 判断\textbf{条件}，为假则结束；
  \item 执行\textbf{循环体}；
  \item 执行\textbf{更新}；
  \item 回到第 2 步。
\end{enumerate}

\begin{figure}[htbp]
  \centering
  \begin{tikzpicture}[x=1cm,y=1cm,>=Stealth,node distance=6mm]
    \node[process] (init) at (0,3.2) {初始化\\\texttt{i = 1}};
    \node[decision] (cond) at (0,1.6) {条件\\\texttt{i <= 100} ?};
    \node[process] (body) at (0,0) {循环体\\\texttt{sum += i}};
    \node[process] (update) at (0,-1.5) {更新\\\texttt{i++}};
    \node[terminator] (end) at (5,1.6) {结束};

    \draw[flow] (init) -- (cond);
    \draw[flow] (cond) -- node[lbl,left]{真} (body);
    \draw[flow] (cond) -- node[lbl,above]{假} (end);
    \draw[flow] (body) -- (update);
    \draw[flow] (update.west) -- ++(-1.2,0) |- (cond.west);
  \end{tikzpicture}
  \caption{\texttt{for} 的执行顺序：初始化一次，然后“条件—循环体—更新”循环}
  \label{fig:for-order}
\end{figure}

\begin{c89box}[C89：循环变量必须提前声明]
C89 不允许在 \code{for} 的初始化部分声明变量。下面两段代码，左边是 C99 写法，右边才是本书要求的 C89 写法：
\begin{codebox}{wrong-for.c / right-for.c —— 两种写法的对比}
/* C99：不允许（本书不使用） */
for (int i = 0; i < n; i++) {
    ...
}

/* C89：正确写法 */
int i;
for (i = 0; i < n; i++) {
    ...
}
\end{codebox}
把 \code{int i;} 写在函数（或语句块）开头，这是本书所有代码的写法。
\end{c89box}

\subsection{\texttt{for} 的灵活写法}

\code{for} 的三个部分都可以为空或使用逗号表达式：

\begin{codebox}{for-flex.c —— for 的灵活用法}
#include <stdio.h>

int main(void)
{
    int i;
    int j;

    /* 从 10 倒数到 1 */
    for (i = 10; i >= 1; i--) {
        printf("%d ", i);
    }
    printf("\n");

    /* 每次加 2：打印 0 2 4 6 8 */
    for (i = 0; i < 10; i += 2) {
        printf("%d ", i);
    }
    printf("\n");

    /* 逗号表达式：同时控制两个变量 */
    for (i = 0, j = 10; i < j; i++, j--) {
        printf("(%d,%d) ", i, j);
    }
    printf("\n");

    return 0;
}
\end{codebox}

\begin{outbox}{运行结果}
10 9 8 7 6 5 4 3 2 1 
0 2 4 6 8 
(0,10) (1,9) (2,8) (3,7) (4,6) 
\end{outbox}

\begin{pitfall}[分号陷阱]
\code{for (i = 0; i < 10; i++);} 末尾的分号会让循环体变成\textbf{空语句}，循环空转十次，后面的代码只执行一次。这是很常见的错误。
\end{pitfall}

\section{\texttt{do-while} 循环}

\code{do-while} 的特点是：\textbf{先执行一次循环体，再判断条件}。所以循环体至少执行一次。

\begin{center}
\fbox{\parbox{0.78\textwidth}{
\ttfamily do \{ \\
\phantom{xxxx}循环体; \\
\ttfamily \} while (条件);
}}
\end{center}

\begin{codebox}{do-while.c —— 至少执行一次}
#include <stdio.h>

int main(void)
{
    int n;

    do {
        printf("请输入 1-10 之间的数：");
        scanf("%d", &n);
    } while (n < 1 || n > 10);

    printf("你输入了 %d\n", n);
    return 0;
}
\end{codebox}

\begin{outbox}{运行示例}
请输入 1-10 之间的数：0
请输入 1-10 之间的数：11
请输入 1-10 之间的数：5
你输入了 5
\end{outbox}

这种“至少问一次，不合法就再问”的场景，\code{do-while} 比 \code{while} 更自然。

\begin{codebox}{digits.c —— 数一数有几位}
#include <stdio.h>

int main(void)
{
    int n;
    int count = 0;

    printf("请输入一个非负整数：");
    scanf("%d", &n);

    do {
        count++;
        n = n / 10;
    } while (n != 0);

    printf("共有 %d 位\n", count);
    return 0;
}
\end{codebox}

\begin{outbox}{运行示例}
请输入一个非负整数：2026
共有 4 位
\end{outbox}

如果 \code{n} 是 0，\code{do-while} 也能正确输出 1 位；换成 \code{while (n != 0)} 则一次都不会执行，结果为 0。

\section{\texttt{break} 与 \texttt{continue}}

\subsection{\texttt{break}：提前跳出循环}

\code{break} 立即结束当前循环，跳到循环后面的语句。

\begin{codebox}{break.c —— 找到就停}
#include <stdio.h>

int main(void)
{
    int i;

    for (i = 1; i <= 100; i++) {
        if (i * i > 50) {
            printf("第一个平方大于 50 的数是 %d\n", i);
            break;
        }
    }

    return 0;
}
\end{codebox}

\begin{outbox}{运行结果}
第一个平方大于 50 的数是 8
\end{outbox}

\code{break} 也常用于“搜索”：找到目标就停止，不必把剩下的都检查一遍。

\subsection{\texttt{continue}：跳过本次，继续下一次}

\code{continue} 跳过循环体剩下的部分，直接进入下一次循环（对 \code{for} 来说，会先执行更新部分）。

\begin{codebox}{continue.c —— 只打印奇数}
#include <stdio.h>

int main(void)
{
    int i;

    for (i = 1; i <= 10; i++) {
        if (i % 2 == 0) {
            continue;       /* 偶数跳过 */
        }
        printf("%d ", i);
    }
    printf("\n");

    return 0;
}
\end{codebox}

\begin{outbox}{运行结果}
1 3 5 7 9 
\end{outbox}

\begin{pitfall}[\texttt{break} 只跳出最内层循环]
在嵌套循环里，\code{break} 只结束它所在的那一层，外层循环继续。如果要一次跳出多层，需要额外的标志变量或 \code{goto}（不推荐）。
\end{pitfall}

\begin{table}[htbp]
  \centering
  \caption{\texttt{break} 与 \texttt{continue} 的区别}
  \begin{tabular}{@{}lll@{}}
    \toprule
    语句 & 作用 & 之后执行什么 \\
    \midrule
    \code{break} & 立即结束当前循环 & 循环后面的语句 \\
    \code{continue} & 跳过本次剩余部分 & 下一次循环（\code{for} 先更新） \\
    \bottomrule
  \end{tabular}
\end{table}

\section{嵌套循环}

循环体里再写一个循环，叫\textbf{嵌套循环}。外层循环每执行一次，内层循环完整地执行一轮。

\begin{codebox}{nested-loop.c —— 嵌套循环}
#include <stdio.h>

int main(void)
{
    int i, j;

    for (i = 1; i <= 3; i++) {
        for (j = 1; j <= 4; j++) {
            printf("(%d,%d) ", i, j);
        }
        printf("\n");
    }

    return 0;
}
\end{codebox}

\begin{outbox}{运行结果}
(1,1) (1,2) (1,3) (1,4) 
(2,1) (2,2) (2,3) (2,4) 
(3,1) (3,2) (3,3) (3,4) 
\end{outbox}

外层 \code{i} 每取一个值，内层 \code{j} 就从 1 到 4 走一遍，所以共执行 $3 \times 4 = 12$ 次。

\subsection{例：九九乘法表}

\begin{example}
输出九九乘法表（下三角）。
\end{example}
\begin{solution}
思路如下：

\begin{codebox}{multiply.c —— 九九乘法表}
#include <stdio.h>

int main(void)
{
    int i, j;

    for (i = 1; i <= 9; i++) {
        for (j = 1; j <= i; j++) {
            printf("%d*%d=%-2d ", j, i, i * j);
        }
        printf("\n");
    }

    return 0;
}
\end{codebox}

\begin{outbox}{运行结果（节选）}
1*1=1  
1*2=2  2*2=4  
1*3=3  2*3=6  3*3=9  
...
1*9=9  2*9=18 3*9=27 ... 9*9=81
\end{outbox}

内层循环的上限是 \code{i}，所以每行输出 \code{i} 项，形成下三角。\texttt{\%-2d} 表示左对齐、占两格，让各列大致对齐。
\end{solution}

\subsection{例：星号三角形}

\begin{example}
输入行数 $n$，输出一个由星号组成的直角三角形。
\end{example}
\begin{solution}
思路如下：

\begin{codebox}{stars.c —— 星号三角形}
#include <stdio.h>

int main(void)
{
    int n;
    int i, j;

    printf("请输入行数：");
    scanf("%d", &n);

    for (i = 1; i <= n; i++) {
        for (j = 1; j <= i; j++) {
            printf("*");
        }
        printf("\n");
    }

    return 0;
}
\end{codebox}

\begin{outbox}{运行示例（n = 5）}
*
**
***
****
*****
\end{outbox}

如果要做“金字塔”（居中对齐），需要在星号前打印空格：第 \code{i} 行先打印 \code{n - i} 个空格，再打印 \code{2*i - 1} 个星号。这个练习留给你自己完成。
\end{solution}

\section{循环的经典算法}

\subsection{累加与累乘}

累加的模式是 \code{sum = sum + 项}，累乘的模式是 \code{product = product * 项}。注意累乘的初始值必须是 1，累加的初始值必须是 0。

\begin{codebox}{factorial.c —— 阶乘}
#include <stdio.h>

int main(void)
{
    int n;
    int i;
    int fact = 1;

    printf("请输入 n：");
    scanf("%d", &n);

    for (i = 1; i <= n; i++) {
        fact = fact * i;
    }

    printf("%d! = %d\n", n, fact);
    return 0;
}
\end{codebox}

\begin{outbox}{运行示例}
请输入 n：10
10! = 3628800
\end{outbox}

\begin{pitfall}[阶乘增长极快]
$13! = 6227020800$，已经超过 32 位 \code{int} 的范围。计算更大的阶乘需要 \code{long long}（C99）或自行实现大整数。不要假设 \code{int} 能装下所有结果。
\end{pitfall}

\subsection{最大值与最小值}

\begin{codebox}{max-min.c —— 求最大值和最小值}
#include <stdio.h>

int main(void)
{
    int n, i;
    int x;
    int max, min;

    printf("请输入数据个数：");
    scanf("%d", &n);

    printf("请输入 %d 个整数：", n);
    scanf("%d", &x);
    max = x;
    min = x;

    for (i = 2; i <= n; i++) {
        scanf("%d", &x);
        if (x > max) {
            max = x;
        }
        if (x < min) {
            min = x;
        }
    }

    printf("最大值 = %d，最小值 = %d\n", max, min);
    return 0;
}
\end{codebox}

\begin{outbox}{运行示例}
请输入数据个数：5
请输入 5 个整数：3 9 1 7 5
最大值 = 9，最小值 = 1
\end{outbox}

注意第一个数要先读进来作为 \code{max} 和 \code{min} 的初始值，而不是把初值设成 0——否则全是负数时最大值会错。

\subsection{判断素数}

素数是只能被 1 和自身整除的大于 1 的整数。最直接的判断方法是试除：

\begin{codebox}{prime.c —— 判断素数}
#include <stdio.h>

int main(void)
{
    int n;
    int i;
    int is_prime = 1;

    printf("请输入一个整数：");
    scanf("%d", &n);

    if (n <= 1) {
        is_prime = 0;
    } else {
        for (i = 2; i * i <= n; i++) {
            if (n % i == 0) {
                is_prime = 0;
                break;
            }
        }
    }

    if (is_prime) {
        printf("%d 是素数\n", n);
    } else {
        printf("%d 不是素数\n", n);
    }

    return 0;
}
\end{codebox}

\begin{outbox}{运行示例}
请输入一个整数：97
97 是素数
\end{outbox}

为什么循环条件是 \code{i * i <= n}？因为如果 $n$ 有大于 $\sqrt{n}$ 的因子，那么它一定还有一个小于 $\sqrt{n}$ 的因子。检查到 $\sqrt{n}$ 就够了。

\subsection{打印 100 以内的素数}

\begin{codebox}{primes100.c —— 100 以内的素数}
#include <stdio.h>

int main(void)
{
    int n, i, is_prime;

    for (n = 2; n <= 100; n++) {
        is_prime = 1;
        for (i = 2; i * i <= n; i++) {
            if (n % i == 0) {
                is_prime = 0;
                break;
            }
        }
        if (is_prime) {
            printf("%d ", n);
        }
    }
    printf("\n");

    return 0;
}
\end{codebox}

\begin{outbox}{运行结果}
2 3 5 7 11 13 17 19 23 29 31 37 41 43 47 53 59 61 67 71 73 79 83 89 97
\end{outbox}

这是嵌套循环的典型用法：外层枚举候选数，内层判断素数。

\subsection{斐波那契数列}

斐波那契数列：$F_1 = 1$，$F_2 = 1$，$F_n = F_{n-1} + F_{n-2}$。

\begin{codebox}{fib.c —— 前 20 项斐波那契数}
#include <stdio.h>

int main(void)
{
    int i;
    int a = 1;
    int b = 1;
    int c;

    printf("%d %d ", a, b);
    for (i = 3; i <= 20; i++) {
        c = a + b;
        printf("%d ", c);
        a = b;
        b = c;
    }
    printf("\n");

    return 0;
}
\end{codebox}

\begin{outbox}{运行结果}
1 1 2 3 5 8 13 21 34 55 89 144 233 377 610 987 1597 2584 4181 6765
\end{outbox}

这个“滚动变量”的技巧很常用：只保留最近两个值，不断往前推。

\subsection{辗转相除法求最大公约数}

\begin{codebox}{gcd.c —— 最大公约数}
#include <stdio.h>

int main(void)
{
    int a, b;
    int t;

    printf("请输入两个正整数：");
    scanf("%d%d", &a, &b);

    while (b != 0) {
        t = a % b;
        a = b;
        b = t;
    }

    printf("最大公约数 = %d\n", a);
    return 0;
}
\end{codebox}

\begin{outbox}{运行示例}
请输入两个正整数：48 18
最大公约数 = 6
\end{outbox}

辗转相除法（欧几里得算法）的思想是：$\gcd(a, b) = \gcd(b, a \bmod b)$，直到余数为 0，此时的除数就是最大公约数。

\subsection{水仙花数}

水仙花数是指一个三位数，其各位数字的立方和等于它本身，例如 $153 = 1^3 + 5^3 + 3^3$。

\begin{codebox}{narcissus.c —— 水仙花数}
#include <stdio.h>

int main(void)
{
    int n;
    int a, b, c;

    for (n = 100; n <= 999; n++) {
        a = n / 100;
        b = n / 10 % 10;
        c = n % 10;
        if (a * a * a + b * b * b + c * c * c == n) {
            printf("%d ", n);
        }
    }
    printf("\n");

    return 0;
}
\end{codebox}

\begin{outbox}{运行结果}
153 370 371 407
\end{outbox}

\subsection{反转整数}

\begin{codebox}{reverse-int.c —— 反转整数}
#include <stdio.h>

int main(void)
{
    int n;
    int rev = 0;

    printf("请输入一个整数：");
    scanf("%d", &n);

    while (n != 0) {
        rev = rev * 10 + n % 10;
        n = n / 10;
    }

    printf("反转后 = %d\n", rev);
    return 0;
}
\end{codebox}

\begin{outbox}{运行示例}
请输入一个整数：12345
反转后 = 54321
\end{outbox}

这个算法对负数也能工作，但反转后的值可能溢出（例如反转 2147483647），实际工程中要检查。

\section{怎样写好一个循环}

\subsection{四步检查法}

写循环前，先回答四个问题：

\begin{enumerate}
  \item \textbf{从哪里开始}？（初始化）
  \item \textbf{做到什么时候停}？（条件）
  \item \textbf{每次做什么}？（循环体）
  \item \textbf{怎么走向结束}？（更新）
\end{enumerate}

写完后，用三个值测试：\textbf{最小规模}（0 次或 1 次）、\textbf{典型值}、\textbf{边界值}。例如求和 1 到 100，要测试 \code{n = 0}、\code{n = 1}、\code{n = 100}。

\subsection{边界：\texttt{<=} 还是 \texttt{<}}

\begin{center}
\begin{tabular}{@{}lll@{}}
\toprule
写法 & 循环次数 & 说明 \\
\midrule
\code{for (i = 1; i <= n; i++)} & $n$ 次 & 包含 $n$ \\
\code{for (i = 0; i < n; i++)} & $n$ 次 & 不包含 $n$，适合数组下标 \\
\code{for (i = 1; i < n; i++)} & $n-1$ 次 & 少一次，常犯的 off-by-one 错误 \\
\bottomrule
\end{tabular}
\end{center}

数组下标从 0 开始时，\code{for (i = 0; i < n; i++)} 是最标准的写法（第 6 章会反复用）。

\subsection{循环不变量}

\textbf{循环不变量}是每次循环开始时都成立的性质。例如求和循环的不变量是：

\begin{center}
\kaishu 在每轮循环开始时，\code{sum} 等于 $1 + 2 + \cdots + (i-1)$。
\end{center}

验证一下：初始 \code{i = 1} 时，\code{sum = 0}，空和确实为 0；每轮把 \code{i} 加进 \code{sum}，然后 \code{i} 加 1，不变量继续成立；循环结束时 \code{i = 101}，\code{sum} 就是 $1 + \cdots + 100$。

养成想“不变量”的习惯，能帮你快速检查循环写得对不对。

\section{循环的常见错误}

\begin{center}
\begin{tabular}{@{}p{0.38\textwidth}p{0.52\textwidth}@{}}
\toprule
错误 & 后果 \\
\midrule
忘了更新循环变量 & 死循环 \\
\code{for (...);} 多写分号 & 循环体为空，后面的语句只执行一次 \\
\code{while (i <= n)} 里 \code{n} 被修改 & 循环次数不可预测 \\
用浮点数做循环变量 & 累积误差可能让循环多一次或少一次 \\
\code{unsigned} 变量写 \code{i >= 0} & 永远为真，死循环 \\
循环体内修改循环变量 & 次数混乱 \\
嵌套循环里 \code{break} & 只跳出内层 \\
\bottomrule
\end{tabular}
\end{center}

\begin{pitfall}[无符号变量的倒序循环]
\begin{codebox}{unsigned-loop.c —— 死循环陷阱}
#include <stdio.h>

int main(void)
{
    unsigned int i;

    for (i = 10; i >= 0; i--) {
        printf("%u ", i);   /* 0 之后再减 1 会变成 4294967295 */
    }
    return 0;
}
\end{codebox}
这个循环停不下来。正确写法是用 \code{int}，或者写成 \code{for (i = 10; i > 0; i--)} 并单独处理 0。
\end{pitfall}

\begin{pitfall}[浮点数循环变量]
\begin{codebox}{float-loop.c —— 浮点循环的隐患}
for (x = 0.0; x != 1.0; x += 0.1) {
    ...
}
\end{codebox}
0.1 在二进制里不能精确表示，累加十次后 \code{x} 可能不等于 1.0，循环停不下来。循环次数应该用整数控制。
\end{pitfall}

\section{综合例题}

\begin{example}
输入 $n$，计算 $1 + \frac{1}{2} + \frac{1}{3} + \cdots + \frac{1}{n}$，保留 4 位小数。
\end{example}
\begin{solution}
思路如下：

\begin{codebox}{harmonic.c —— 调和级数}
#include <stdio.h>

int main(void)
{
    int n, i;
    double sum = 0.0;

    printf("请输入 n：");
    scanf("%d", &n);

    for (i = 1; i <= n; i++) {
        sum = sum + 1.0 / i;
    }

    printf("sum = %.4f\n", sum);
    return 0;
}
\end{codebox}

注意 \code{1.0 / i} 而不是 \code{1 / i}。后者是整数除法，除了 \code{i = 1} 以外全是 0。
\end{solution}

\begin{example}
输入若干个整数（以 0 结束），统计正数、负数的个数，并计算平均值。
\end{example}
\begin{solution}
思路如下：

\begin{codebox}{stat.c —— 统计正负数}
#include <stdio.h>

int main(void)
{
    int n;
    int positive = 0;
    int negative = 0;
    int sum = 0;
    int count = 0;

    printf("请输入整数（0 结束）：");
    scanf("%d", &n);

    while (n != 0) {
        if (n > 0) {
            positive++;
        } else {
            negative++;
        }
        sum = sum + n;
        count++;
        scanf("%d", &n);
    }

    printf("正数 %d 个，负数 %d 个\n", positive, negative);
    if (count > 0) {
        printf("平均值 = %.2f\n", (double)sum / count);
    } else {
        printf("没有输入有效数据\n");
    }

    return 0;
}
\end{codebox}

这里用 \code{(double)sum / count} 做浮点除法；同时处理了“一个数都没输入”的情况，避免除以 0。
\end{solution}

\begin{example}
输出 1000 以内的所有完数。完数是指一个数恰好等于它的真因子之和，例如 $6 = 1 + 2 + 3$。
\end{example}
\begin{solution}
思路如下：

\begin{codebox}{perfect.c —— 完数}
#include <stdio.h>

int main(void)
{
    int n, i, sum;

    for (n = 2; n <= 1000; n++) {
        sum = 0;
        for (i = 1; i < n; i++) {
            if (n % i == 0) {
                sum = sum + i;
            }
        }
        if (sum == n) {
            printf("%d ", n);
        }
    }
    printf("\n");

    return 0;
}
\end{codebox}

\begin{outbox}{运行结果}
6 28 496
\end{outbox}
\end{solution}

\begin{example}
输入一个整数，输出它的二进制表示（从最高位到最低位）。
\end{example}
\begin{solution}
可以用“掩码”从高位往低位检查。以 32 位 \code{unsigned int} 为例：

\begin{codebox}{to-binary.c —— 输出二进制}
#include <stdio.h>

int main(void)
{
    unsigned int n;
    int i;

    printf("请输入一个非负整数：");
    scanf("%u", &n);

    for (i = 31; i >= 0; i--) {
        if (n & (1u << i)) {
            printf("1");
        } else {
            printf("0");
        }
        if (i % 8 == 0) {
            printf(" ");
        }
    }
    printf("\n");

    return 0;
}
\end{codebox}

\begin{outbox}{运行示例}
请输入一个非负整数：182
00000000 00000000 00000000 10110110 
\end{outbox}

\code{1u << i} 生成只有第 \code{i} 位为 1 的掩码，\texttt{n \& 掩码} 判断该位是否为 1。这是位运算和循环的经典结合。
\end{solution}

\section{拔高：循环的效率与递归}

\subsection{时间复杂度直觉}

循环执行次数决定了程序要花多少时间。求 1 到 $n$ 的和需要循环 $n$ 次，记作 $O(n)$；判断一个数是否为素数只需要试除到 $\sqrt{n}$，记作 $O(\sqrt{n})$。数据量增大时，$O(\sqrt{n})$ 比 $O(n)$ 快得多。

\begin{center}
\begin{tabular}{@{}lll@{}}
\toprule
$n$ & $O(n)$ 次操作 & $O(\sqrt{n})$ 次操作 \\
\midrule
100 & 100 & 10 \\
10000 & 10000 & 100 \\
1000000 & 1000000 & 1000 \\
\bottomrule
\end{tabular}
\end{center}

\subsection{递归：循环的另一种写法}

有些问题用“自己调用自己”描述更自然，例如阶乘：

\begin{codebox}{fact-rec.c —— 递归求阶乘}
#include <stdio.h>

int factorial(int n)
{
    if (n <= 1) {
        return 1;
    }
    return n * factorial(n - 1);
}

int main(void)
{
    printf("%d\n", factorial(5));
    return 0;
}
\end{codebox}

递归和循环可以互相转换。初学阶段先把循环写熟，递归等学完函数（后续章节）再深入。

\subsection{\texttt{goto}：不推荐，但要知道}

\code{goto} 可以跳到同一函数内的标号处，常用于“多层循环一次跳出”或“统一清理资源”。但它会破坏代码结构，现代工程中很少使用。知道有这个东西即可，不要主动用。

\section{本章小结}

\begin{center}
\begin{tabular}{@{}p{0.24\textwidth}p{0.66\textwidth}@{}}
\toprule
主题 & 要点 \\
\midrule
\code{while} & 先判断后执行；适合次数不确定的循环 \\
\code{for} & 初始化、条件、更新集中在一行；适合次数确定的循环 \\
\code{do-while} & 先执行一次再判断；适合“至少做一次”的场景 \\
C89 循环变量 & 必须在循环外声明 \\
\code{break} & 立即结束当前循环 \\
\code{continue} & 跳过本次，进入下一次 \\
嵌套循环 & 外层每走一步，内层走完整一轮 \\
写循环 & 初始化、条件、循环体、更新，四步缺一不可 \\
常见错误 & 死循环、多余分号、边界差一、无符号倒序 \\
\bottomrule
\end{tabular}
\end{center}

\begin{tipbox}[下一章预告]
有了循环，我们可以处理大量数据了。但数据放哪里？如果每读一个数就丢了，循环再厉害也没用。下一章学习\textbf{数组}——把一堆同类型的数据整齐地存在一起。
\end{tipbox}

\begin{problemset}
  \item 用 \code{for}、\code{while}、\code{do-while} 三种循环分别求 $1+2+\cdots+100$。
  \item 输入 $n$，求 $n!$，并测试 $n=10$、$n=13$ 的结果，解释差异。
  \item 输入 10 个整数，求最大值、最小值和平均值。
  \item 输出 1\textasciitilde100 中所有能被 3 或 5 整除的数，每行 5 个。
  \item 判断一个数是否为素数，输出 100 以内的所有素数。
  \item 输出前 30 项斐波那契数列。
  \item 输入两个正整数，用辗转相除法求最大公约数和最小公倍数。
  \item 输出所有三位水仙花数。
  \item 输入一个整数，输出它的位数、各位数字之和、逆序数。
  \item 思考题：为什么 \code{for (i = 0; i <= n; i++)} 在处理长度为 \code{n} 的数组时会越界？应该怎么改？
\end{problemset}

\chapter{数组：优雅的数据存取方式}

\begin{introduction}
  \item 为什么需要数组：一批同类型数据的家
  \item 一维数组的定义、初始化、下标与遍历
  \item 数组在内存中的连续排布
  \item 求最值、逆序、统计、查找等经典算法
  \item 二维数组与矩阵
  \item 数组的边界、退化与常见错误
\end{introduction}

\chaptermotto{「把一百个数据排好队，你只需要记住队首的位置。」}

前面的程序一次只能处理一两个数。可现实中我们经常要处理一批数据：全班 50 个成绩、一帧 128 个采样点、一张 $3 \times 4$ 的矩阵。如果每个数据都用一个变量，代码会变得无法维护。

数组就是为“一批同类型数据”准备的容器。它让数据排成连续的一队，用一个名字和下标来访问。这一章会把数组的“形”和“神”都讲清楚：既会用它，也知道它在内存里长什么样。

\section{为什么需要数组}

假设要读入 10 个成绩并求平均分。不用数组，你可能写成这样：

\begin{codebox}{scores-ugly.c —— 没有数组的笨办法}
#include <stdio.h>

int main(void)
{
    int s1, s2, s3, s4, s5, s6, s7, s8, s9, s10;
    int sum;

    printf("请输入 10 个成绩：");
    scanf("%d%d%d%d%d%d%d%d%d%d",
          &s1, &s2, &s3, &s4, &s5,
          &s6, &s7, &s8, &s9, &s10);

    sum = s1 + s2 + s3 + s4 + s5 + s6 + s7 + s8 + s9 + s10;
    printf("平均分 = %.2f\n", sum / 10.0);
    return 0;
}
\end{codebox}

这段代码能跑，但问题很明显：

\begin{itemize}
  \item 变量太多，读起来累；
  \item 想求最大值、排序、找某个分数，几乎没法写；
  \item 如果成绩变成 100 个，代码就没法维护了。
\end{itemize}

用数组，只需要一个名字 \code{scores} 和一个下标：

\begin{codebox}{scores-array.c —— 用数组读入 10 个成绩}
#include <stdio.h>

int main(void)
{
    int scores[10];
    int i;
    int sum = 0;

    printf("请输入 10 个成绩：");
    for (i = 0; i < 10; i++) {
        scanf("%d", &scores[i]);
    }

    for (i = 0; i < 10; i++) {
        sum = sum + scores[i];
    }

    printf("平均分 = %.2f\n", sum / 10.0);
    return 0;
}
\end{codebox}

\begin{outbox}{运行示例}
请输入 10 个成绩：90 85 77 92 88 60 95 73 81 100
平均分 = 84.10
\end{outbox}

\begin{center}
\fbox{\parbox{0.8\textwidth}{\centering\kaishu
数组 = 一段连续的内存 + 一个名字 + 一组从 0 开始的下标。}}
\end{center}

\section{一维数组}

\subsection{定义}

\begin{center}
\fbox{\parbox{0.62\textwidth}{
\ttfamily 类型 数组名[元素个数];
}}
\end{center}

例如：

\begin{codebox}{define-array.c —— 数组的定义}
int a[10];          /* 10 个 int，下标 0 到 9 */
double b[5];        /* 5 个 double */
char name[20];      /* 20 个 char，第 9 章会用作字符串 */
\end{codebox}

几个要点：

\begin{itemize}
  \item 元素个数必须是\textbf{编译期常量}，不能是变量（C89 不支持变长数组）；
  \item 下标从 \textbf{0} 开始，到 \textbf{元素个数 $-$ 1} 结束；
  \item 数组名是一个整体，代表这段连续内存的首地址（第 8 章会展开）。
\end{itemize}

\begin{c89box}[C89：数组长度必须是常量]
\begin{codebox}{vla.c —— 错误示范（C99 才允许）}
int n = 10;
int a[n];       /* C89 不允许：n 不是编译期常量 */
\end{codebox}
C89 里请写成 \code{int a[10];}。如果确实需要运行时决定长度，应该用 \code{malloc} 动态分配（第 8 章讲指针时介绍），或者定义一个足够大的数组。
\end{c89box}

\subsection{内存布局：连续排布}

数组最重要的性质是：\textbf{所有元素在内存中连续存放}。下图是一个 \code{int a[5]} 的内存示意（假设 \code{int} 占 4 字节，首地址为 0x1000）：

\begin{figure}[htbp]
  \centering
  \begin{tikzpicture}[x=1cm,y=1cm]
    \foreach \i/\v/\addr in {0/10/0x1000,1/20/0x1004,2/30/0x1008,3/40/0x100C,4/50/0x1010} {
      \node[membox,minimum width=1.6cm] (e\i) at (\i*1.8,0) {\v};
      \node[font=\scriptsize,color=black!60,below=1mm of e\i] {a[\i]};
      \node[font=\scriptsize,color=black!45,above=1mm of e\i] {\texttt{\addr}};
    }
    \draw[decorate,decoration={brace,amplitude=4pt},structurecolor!70,line width=0.7pt]
      (-0.8,-1.05) -- (7.4,-1.05);
    \node[font=\small,below=6mm] at (3.3,-1.05) {5 个元素，共 20 字节，连续存放};
  \end{tikzpicture}
  \caption{\texttt{int a[5]} 的内存布局：元素一个挨一个}
  \label{fig:array-memory}
\end{figure}

因为连续，所以第 \code{i} 个元素的地址可以由首地址算出来：

\[
\texttt{\&a[i]} = \texttt{\&a[0]} + i \times \texttt{sizeof(int)}.
\]

这个公式是后面理解指针的关键。现在你只需要记住：\textbf{下标就是“第几个”，地址是“在哪里”，两者可以互相换算}。

\subsection{初始化}

\begin{codebox}{init-array.c —— 数组的初始化}
int a[5] = {10, 20, 30, 40, 50};  /* 全部初始化 */
int b[5] = {1, 2};                /* 部分初始化，其余为 0 */
int c[5] = {0};                   /* 全部为 0 的常用写法 */
int d[]  = {1, 2, 3};             /* 省略长度，编译器数出 3 */
\end{codebox}

\begin{itemize}
  \item 部分初始化时，剩下的元素自动补 0；
  \item \texttt{\{0\}} 是“全部清零”的惯用写法；
  \item 省略长度时，编译器根据初值个数确定大小。
\end{itemize}

\begin{pitfall}[局部数组不会自动清零]
\code{int a[10];} 只是分配了空间，里面的值是随机的“垃圾值”。如果要用，必须先赋值或初始化。只有全局数组和静态数组才会自动清零。
\end{pitfall}

\subsection{下标访问与遍历}

\begin{codebox}{traverse.c —— 数组的遍历}
#include <stdio.h>

int main(void)
{
    int a[5] = {10, 20, 30, 40, 50};
    int i;

    for (i = 0; i < 5; i++) {
        printf("a[%d] = %d\n", i, a[i]);
    }

    return 0;
}
\end{codebox}

\begin{outbox}{运行结果}
a[0] = 10
a[1] = 20
a[2] = 30
a[3] = 40
a[4] = 50
\end{outbox}

遍历的标准写法是：

\begin{codebox}{traverse-template.c —— 遍历模板}
int i;
for (i = 0; i < 5; i++) {
    /* 使用 a[i] */
}
\end{codebox}

注意条件是 \code{i < 5}，不是 \code{i <= 5}。\code{a[5]} 已经越界了！

\subsection{数组大小与元素个数}

\begin{codebox}{array-size.c —— 数组占多少字节}
#include <stdio.h>

int main(void)
{
    int a[10];

    printf("sizeof(a)    = %d\n", (int)sizeof(a));
    printf("sizeof(a[0]) = %d\n", (int)sizeof(a[0]));
    printf("元素个数     = %d\n", (int)(sizeof(a) / sizeof(a[0])));

    return 0;
}
\end{codebox}

\begin{outbox}{运行结果（典型 32 位 int 环境）}
sizeof(a)    = 40
sizeof(a[0]) = 4
元素个数     = 10
\end{outbox}

\code{sizeof(a) / sizeof(a[0])} 是求数组元素个数的经典写法。它只在“数组名还是数组”的时候有效，一旦数组名退化成指针（传给函数后），这个写法就会失效——第 8 章会详细解释。

\section{数组的常见操作}

\subsection{例：求最大值及其下标}

\begin{example}
读入 10 个整数，输出最大值以及它在数组中的下标。
\end{example}
\begin{solution}
思路如下：

\begin{codebox}{max-index.c —— 最大值与下标}
#include <stdio.h>

int main(void)
{
    int a[10];
    int i;
    int max;
    int max_index = 0;

    printf("请输入 10 个整数：");
    for (i = 0; i < 10; i++) {
        scanf("%d", &a[i]);
    }

    max = a[0];
    for (i = 1; i < 10; i++) {
        if (a[i] > max) {
            max = a[i];
            max_index = i;
        }
    }

    printf("最大值 = %d，下标 = %d\n", max, max_index);
    return 0;
}
\end{codebox}

\begin{outbox}{运行示例}
请输入 10 个整数：3 9 1 7 9 5 2 8 6 4
最大值 = 9，下标 = 1
\end{outbox}

如果有多个相同的最大值，这个程序返回第一个。想返回最后一个，把 \code{>} 改成 \code{>=} 即可。
\end{solution}

\subsection{例：数组逆序}

\begin{example}
把数组中的元素原地逆序，例如 \texttt{\{1,2,3,4,5\}} 变成 \texttt{\{5,4,3,2,1\}}。
\end{example}
\begin{solution}
用两个下标，一个从头往后，一个从尾往前，交换后向中间靠拢：

\begin{figure}[htbp]
  \centering
  \begin{tikzpicture}[x=1cm,y=1cm]
    \foreach \i/\v in {0/1,1/2,2/3,3/4,4/5} {
      \node[membox,minimum width=1.4cm] (e\i) at (\i*1.6,0) {\v};
      \node[font=\scriptsize,color=black!55,below=1mm of e\i] {\i};
    }
    \draw[ptr] ([yshift=3mm]e0.north) to[bend left=25] ([yshift=3mm]e4.north);
    \draw[ptr] ([yshift=-3mm]e1.south) to[bend right=25] ([yshift=-3mm]e3.south);
    \node[font=\scriptsize,color=ptrline,above=10mm] at (3.2,0) {交换 a[0] 与 a[4]，a[1] 与 a[3]};
  \end{tikzpicture}
  \caption{双下标向中间靠拢，完成原地逆序}
  \label{fig:reverse-array}
\end{figure}

\begin{codebox}{reverse-array.c —— 原地逆序}
#include <stdio.h>

int main(void)
{
    int a[10];
    int i, j, temp;
    int n = 10;

    printf("请输入 10 个整数：");
    for (i = 0; i < n; i++) {
        scanf("%d", &a[i]);
    }

    for (i = 0, j = n - 1; i < j; i++, j--) {
        temp = a[i];
        a[i] = a[j];
        a[j] = temp;
    }

    printf("逆序后：");
    for (i = 0; i < n; i++) {
        printf("%d ", a[i]);
    }
    printf("\n");

    return 0;
}
\end{codebox}

\begin{outbox}{运行示例}
请输入 10 个整数：1 2 3 4 5 6 7 8 9 10
逆序后：10 9 8 7 6 5 4 3 2 1
\end{outbox}

循环条件 \code{i < j} 保证中间的元素不会被重复交换。
\end{solution}

\subsection{例：统计每个数字出现的次数}

\begin{example}
输入一个非负整数，统计它各个数字（0\textasciitilde9）出现的次数。
\end{example}
\begin{solution}
用数组 \code{count[10]} 当下标为 0\textasciitilde9 的“计数器”：

\begin{codebox}{digit-count.c —— 数字频次统计}
#include <stdio.h>

int main(void)
{
    int n;
    int digit;
    int count[10] = {0};
    int i;

    printf("请输入一个非负整数：");
    scanf("%d", &n);

    if (n == 0) {
        count[0] = 1;
    }
    while (n > 0) {
        digit = n % 10;
        count[digit]++;
        n = n / 10;
    }

    for (i = 0; i < 10; i++) {
        if (count[i] > 0) {
            printf("数字 %d 出现了 %d 次\n", i, count[i]);
        }
    }

    return 0;
}
\end{codebox}

\begin{outbox}{运行示例}
请输入一个非负整数：20260226
数字 0 出现了 2 次
数字 2 出现了 4 次
数字 6 出现了 2 次
\end{outbox}

用“值当下标”是数组最常用的技巧之一。统计、计数、查表都靠它。
\end{solution}

\subsection{例：线性查找}

\begin{example}
在数组中查找某个值，找到则输出下标，找不到输出提示。
\end{example}
\begin{solution}
思路如下：

\begin{codebox}{linear-search.c —— 线性查找}
#include <stdio.h>

int main(void)
{
    int a[10] = {3, 9, 1, 7, 5, 2, 8, 6, 4, 0};
    int key;
    int i;
    int found = -1;

    printf("请输入要查找的数：");
    scanf("%d", &key);

    for (i = 0; i < 10; i++) {
        if (a[i] == key) {
            found = i;
            break;
        }
    }

    if (found >= 0) {
        printf("找到了，下标是 %d\n", found);
    } else {
        printf("没有找到\n");
    }

    return 0;
}
\end{codebox}

线性查找从头到尾依次比较，最坏要比较 $n$ 次。
\end{solution}

\subsection{例：二分查找}

\begin{example}
在一个\textbf{已排序}的数组中查找某个值。
\end{example}
\begin{solution}
二分查找每次比较中间元素，根据大小关系把查找范围缩小一半：

\begin{figure}[htbp]
  \centering
  \begin{tikzpicture}[x=1cm,y=1cm]
    \foreach \i/\v in {0/1,1/3,2/5,3/7,4/9,5/11,6/13,7/15,8/17,9/19} {
      \node[membox,minimum width=1.1cm,minimum height=6mm] (e\i) at (\i*1.25,0) {\v};
      \node[font=\scriptsize,color=black!55,below=1mm of e\i] {\i};
    }
    \node[font=\scriptsize,color=structurecolor,above=3mm of e0] {low};
    \node[font=\scriptsize,color=structurecolor,above=3mm of e9] {high};
    \node[font=\scriptsize,color=ptrline,above=3mm of e4] {mid};
    \draw[ptr] ([yshift=8mm]e4.north) -- ([yshift=4mm]e4.north);
    \draw[decorate,decoration={brace,amplitude=3pt},structurecolor!60]
      ([yshift=-5mm]e0.south west) -- ([yshift=-5mm]e9.south east);
  \end{tikzpicture}
  \caption{二分查找：比较 \texttt{mid}，然后丢掉一半}
  \label{fig:binary-search}
\end{figure}

\begin{codebox}{binary-search.c —— 二分查找}
#include <stdio.h>

int main(void)
{
    int a[10] = {1, 3, 5, 7, 9, 11, 13, 15, 17, 19};
    int key;
    int low = 0;
    int high = 9;
    int mid;
    int found = -1;

    printf("请输入要查找的数：");
    scanf("%d", &key);

    while (low <= high) {
        mid = (low + high) / 2;
        if (a[mid] == key) {
            found = mid;
            break;
        } else if (a[mid] < key) {
            low = mid + 1;
        } else {
            high = mid - 1;
        }
    }

    if (found >= 0) {
        printf("找到了，下标是 %d\n", found);
    } else {
        printf("没有找到\n");
    }

    return 0;
}
\end{codebox}

二分查找的前提是数组\textbf{已经有序}。它的比较次数是 $O(\log n)$，比线性查找快得多：100 万个元素最多比较约 20 次。
\end{solution}

\section{二维数组}

\subsection{定义与初始化}

二维数组可以看成“数组的数组”，例如 \code{int a[3][4]} 表示 3 行 4 列，共 12 个元素。

\begin{codebox}{define-2d.c —— 二维数组的定义与初始化}
int a[3][4];                    /* 3 行 4 列，未初始化 */
int b[2][3] = {                 /* 按行初始化 */
    {1, 2, 3},
    {4, 5, 6}
};
int c[2][3] = {1, 2, 3, 4, 5, 6};   /* 也可以写成一维形式，按行填充 */
int d[3][4] = {{0}};                /* 全部清零 */
\end{codebox}

访问元素用两个下标：\code{a[i][j]}，其中 \code{i} 是行，\code{j} 是列，都从 0 开始。

\subsection{内存布局：按行连续}

二维数组在内存里仍然是\textbf{连续}的，按“先行后列”排列：

\begin{figure}[htbp]
  \centering
  \begin{tikzpicture}[x=1cm,y=1cm]
    \foreach \r/\lab in {0/第 0 行,1/第 1 行,2/第 2 行} {
      \node[font=\small,anchor=east] at (-1.1,-\r*1.0) {\lab};
      \foreach \c in {0,1,2,3} {
        \pgfmathtruncatemacro{\idx}{\r*4+\c}
        \node[membox,minimum width=1.35cm,minimum height=6.5mm]
          (e\r\c) at (\c*1.5,-\r*1.0) {\idx};
      }
    }
    \node[font=\scriptsize,color=black!55] at (2.25,0.8) {线性下标 $0,1,2,\ldots,11$ 按行递增};
  \end{tikzpicture}
  \caption{\texttt{int a[3][4]} 按行存放：第 0 行 4 个，接着第 1 行 4 个……}
  \label{fig:2d-memory}
\end{figure}

元素 \code{a[i][j]} 的地址是：

\[
\texttt{\&a[i][j]} = \texttt{\&a[0][0]} + (i \times 4 + j) \times \texttt{sizeof(int)}.
\]

所以“按行遍历”通常比“按列遍历”更快——因为前者访问的内存是连续的，更符合缓存的工作方式。

\subsection{二维数组的遍历}

\begin{codebox}{traverse-2d.c —— 遍历二维数组}
#include <stdio.h>

int main(void)
{
    int a[3][4] = {
        {1, 2, 3, 4},
        {5, 6, 7, 8},
        {9, 10, 11, 12}
    };
    int i, j;

    for (i = 0; i < 3; i++) {
        for (j = 0; j < 4; j++) {
            printf("%4d", a[i][j]);
        }
        printf("\n");
    }

    return 0;
}
\end{codebox}

\begin{outbox}{运行结果}
   1   2   3   4
   5   6   7   8
   9  10  11  12
\end{outbox}

\subsection{例：矩阵加法}

\begin{example}
输入两个 $3 \times 4$ 的矩阵，输出它们的和。
\end{example}
\begin{solution}
思路如下：

\begin{codebox}{matrix-add.c —— 矩阵加法}
#include <stdio.h>

int main(void)
{
    int a[3][4];
    int b[3][4];
    int c[3][4];
    int i, j;

    printf("请输入矩阵 A（3x4）：\n");
    for (i = 0; i < 3; i++) {
        for (j = 0; j < 4; j++) {
            scanf("%d", &a[i][j]);
        }
    }

    printf("请输入矩阵 B（3x4）：\n");
    for (i = 0; i < 3; i++) {
        for (j = 0; j < 4; j++) {
            scanf("%d", &b[i][j]);
        }
    }

    for (i = 0; i < 3; i++) {
        for (j = 0; j < 4; j++) {
            c[i][j] = a[i][j] + b[i][j];
        }
    }

    printf("A + B =\n");
    for (i = 0; i < 3; i++) {
        for (j = 0; j < 4; j++) {
            printf("%4d", c[i][j]);
        }
        printf("\n");
    }

    return 0;
}
\end{codebox}
\end{solution}

\subsection{例：矩阵转置}

\begin{example}
输入一个 $3 \times 3$ 矩阵，输出它的转置（行列互换）。
\end{example}
\begin{solution}
思路如下：

\begin{codebox}{transpose.c —— 矩阵转置}
#include <stdio.h>

int main(void)
{
    int a[3][3];
    int i, j;

    printf("请输入 3x3 矩阵：\n");
    for (i = 0; i < 3; i++) {
        for (j = 0; j < 3; j++) {
            scanf("%d", &a[i][j]);
        }
    }

    printf("转置后：\n");
    for (i = 0; i < 3; i++) {
        for (j = 0; j < 3; j++) {
            printf("%4d", a[j][i]);
        }
        printf("\n");
    }

    return 0;
}
\end{codebox}

注意输出时用的是 \code{a[j][i]}，而不是 \code{a[i][j]}。
\end{solution}

\section{数组的常见错误}

\begin{center}
\begin{tabular}{@{}p{0.38\textwidth}p{0.52\textwidth}@{}}
\toprule
错误 & 后果 \\
\midrule
下标越界（\code{a[10]} 用于 \code{int a[10]}） & 未定义行为，可能崩溃或悄悄破坏其他数据 \\
数组未初始化就读取 & 得到随机值 \\
用 \code{=} 拷贝数组 & 编译错误；数组不能整体赋值 \\
数组长度用变量 & C89 编译错误 \\
遍历条件写成 \code{i <= n} & 多访问一个元素，越界 \\
\code{sizeof(a)} 用在函数参数上 & 得到的是指针大小，不是数组大小 \\
\bottomrule
\end{tabular}
\end{center}

\begin{pitfall}[数组越界不会报错]
C 语言不会检查数组下标是否越界。\code{a[10]} 会去读写数组后面的内存，可能：
\begin{itemize}
  \item 什么都没发生（最危险，掩盖了 bug）；
  \item 覆盖了别的变量，导致诡异的结果；
  \item 触发段错误，程序崩溃。
\end{itemize}
写循环时一定要反复确认边界：\code{i < n}，不是 \code{i <= n}。
\end{pitfall}

\begin{pitfall}[数组不能整体赋值]
\begin{codebox}{array-assign.c —— 错误：数组整体赋值}
int a[5] = {1, 2, 3, 4, 5};
int b[5];

b = a;      /* 错误：数组名不能作为赋值目标 */
\end{codebox}
正确做法是用循环逐个复制：
\begin{codebox}{array-copy.c —— 逐个复制}
int i;
for (i = 0; i < 5; i++) {
    b[i] = a[i];
}
\end{codebox}
\end{pitfall}

\section{拔高：数组名、退化与栈空间}

\subsection{数组名不是普通的变量}

\code{sizeof(a)} 能得到整个数组的大小，但 \code{a} 在大多数表达式里会“退化”成指向首元素的指针。这就是为什么：

\begin{itemize}
  \item \code{sizeof(a)} 是 40（整个数组）；
  \item 把 \code{a} 传给函数后，函数里的 \code{sizeof(a)} 变成 8（指针大小，64 位环境）。
\end{itemize}

\begin{codebox}{array-decay.c —— 数组名退化}
#include <stdio.h>

void show(int arr[])
{
    printf("函数内 sizeof(arr) = %d\n", (int)sizeof(arr));
}

int main(void)
{
    int a[10];

    printf("main 内 sizeof(a) = %d\n", (int)sizeof(a));
    show(a);
    return 0;
}
\end{codebox}

\begin{outbox}{运行结果（64 位环境）}
main 内 sizeof(a) = 40
函数内 sizeof(arr) = 8
\end{outbox}

所以“在函数里用 \code{sizeof(arr) / sizeof(arr[0])} 求长度”是错的。正确做法是把长度作为参数一起传进去。

\subsection{数组在栈上}

局部数组分配在栈上，大小在编译时确定。栈空间通常只有几 MB，不要定义过大的局部数组：

\begin{codebox}{big-array.c —— 栈溢出的风险}
void f(void)
{
    int a[1000000];   /* 4 MB，可能直接栈溢出 */
    ...
}
\end{codebox}

需要大数组时，可以：
\begin{itemize}
  \item 定义为全局数组或静态数组（放在数据区，不在栈上）；
  \item 用 \code{malloc} 动态分配（第 8 章）。
\end{itemize}

\subsection{变长数组（VLA）}

C99 支持用变量定义数组长度（\code{int a[n];}），但 C89 不支持，而且 VLA 在栈上分配、大小不可控，不建议在嵌段式系统里使用。本书坚持 C89 写法。

\section{综合例题}

\begin{example}
读入 10 个整数，统计其中正数、负数和零的个数。
\end{example}
\begin{solution}
思路如下：

\begin{codebox}{count-sign.c —— 统计正负零}
#include <stdio.h>

int main(void)
{
    int a[10];
    int i;
    int positive = 0;
    int negative = 0;
    int zero = 0;

    printf("请输入 10 个整数：");
    for (i = 0; i < 10; i++) {
        scanf("%d", &a[i]);
    }

    for (i = 0; i < 10; i++) {
        if (a[i] > 0) {
            positive++;
        } else if (a[i] < 0) {
            negative++;
        } else {
            zero++;
        }
    }

    printf("正数 %d 个，负数 %d 个，零 %d 个\n",
           positive, negative, zero);
    return 0;
}
\end{codebox}
\end{solution}

\begin{example}
用数组求斐波那契数列的前 20 项。
\end{example}
\begin{solution}
思路如下：

\begin{codebox}{fib-array.c —— 用数组存斐波那契}
#include <stdio.h>

int main(void)
{
    int f[20];
    int i;

    f[0] = 1;
    f[1] = 1;
    for (i = 2; i < 20; i++) {
        f[i] = f[i - 1] + f[i - 2];
    }

    for (i = 0; i < 20; i++) {
        printf("%d ", f[i]);
    }
    printf("\n");

    return 0;
}
\end{codebox}

用数组的好处是：所有项都保留下来，可以随时回头访问 \code{f[k]}。
\end{solution}

\begin{example}
输入 10 个学生的成绩，输出高于平均分的成绩。
\end{example}
\begin{solution}
思路如下：

\begin{codebox}{above-average.c —— 高于平均分}
#include <stdio.h>

int main(void)
{
    int scores[10];
    int i;
    int sum = 0;
    double average;

    printf("请输入 10 个成绩：");
    for (i = 0; i < 10; i++) {
        scanf("%d", &scores[i]);
        sum = sum + scores[i];
    }

    average = sum / 10.0;
    printf("平均分 = %.2f\n", average);
    printf("高于平均分的成绩：");
    for (i = 0; i < 10; i++) {
        if (scores[i] > average) {
            printf("%d ", scores[i]);
        }
    }
    printf("\n");

    return 0;
}
\end{codebox}

注意 \code{scores[i] > average} 比较的是 \code{int} 和 \code{double}，\code{scores[i]} 会自动转成 \code{double}，没有问题。
\end{solution}

\section{本章小结}

\begin{center}
\begin{tabular}{@{}p{0.26\textwidth}p{0.64\textwidth}@{}}
\toprule
主题 & 要点 \\
\midrule
数组 & 一段连续内存，元素类型相同，下标从 0 开始 \\
定义 & \code{类型 名字[长度];}，长度必须是编译期常量 \\
初始化 & \texttt{\{...\}}；部分初始化补 0；\texttt{\{0\}} 清零 \\
遍历 & \code{for (i = 0; i < n; i++)}，注意边界 \\
内存 & 元素连续存放；下标可换算成地址偏移 \\
二维数组 & \code{a[i][j]}，按行连续存放 \\
常见算法 & 最值、逆序、统计、线性查找、二分查找、矩阵 \\
常见错误 & 越界、未初始化、整体赋值、长度用变量 \\
\bottomrule
\end{tabular}
\end{center}

\begin{tipbox}[下一章预告]
数组只能装同一种类型的数据。如果要描述一个“学生”，需要同时存姓名、学号、成绩，类型各不相同，数组就不够用了。下一章学习\textbf{结构体}——把不同类型的数据打包成一个整体。
\end{tipbox}

\begin{problemset}
  \item 定义 \code{int a[10]}，读入 10 个数并输出，要求每行 5 个。
  \item 读入 10 个数，求最大值、最小值、平均值。
  \item 读入 10 个数，原地逆序后输出。
  \item 读入 10 个数，统计其中偶数和奇数的个数。
  \item 读入一个整数，统计每个数字出现的次数（0\textasciitilde9）。
  \item 在有序数组中用二分查找查找某个数，输出比较次数。
  \item 输入两个 $3 \times 3$ 矩阵，输出它们的和与转置。
  \item 输入 10 个成绩，输出高于平均分的成绩及其下标。
  \item 用数组求斐波那契数列前 30 项（注意是否会溢出）。
  \item 思考题：为什么 \code{sizeof(a) / sizeof(a[0])} 在函数内部求数组长度会失效？应该怎么传长度？
\end{problemset}

\chapter{结构体：古希腊掌管分类的神}

\begin{introduction}
  \item 为什么需要结构体：把不同类型的数据打包成一个整体
  \item 结构体的定义、声明、初始化与成员访问
  \item \texttt{typedef} 与嵌套结构体
  \item 结构体数组与结构体作为函数参数
  \item 内存对齐：结构体到底占多少字节
  \item 位域与联合体（拔高）
\end{introduction}

\chaptermotto{「世界不是由同一种东西组成的，程序也不该只用一种类型描述世界。」}

数组很优雅，但它有一个前提：所有元素类型相同。可现实中的一个“学生”包含学号（整数）、姓名（字符串）、成绩（小数），类型各不相同。你当然可以用三个平行数组分别存，但数据一旦多了，维护起来就是灾难。

结构体（structure）就是为此而生：它把若干个不同类型的数据打包成一个整体，给这个整体起一个名字。学会结构体，你就从“处理数据”迈向了“描述对象”。

\section{为什么需要结构体}

假设要管理 3 个学生的信息。用平行数组的写法是这样的：

\begin{codebox}{parallel-arrays.c —— 用平行数组存学生信息}
#include <stdio.h>

int main(void)
{
    int ids[3] = {1001, 1002, 1003};
    char names[3][20] = {"Alice", "Bob", "Carol"};
    double scores[3] = {92.5, 78.0, 85.5};
    int i;

    for (i = 0; i < 3; i++) {
        printf("%d %s %.1f\n", ids[i], names[i], scores[i]);
    }

    return 0;
}
\end{codebox}

这个写法的问题在于：三个数组必须\textbf{一一对应}。如果排序时只交换了 \code{scores}，忘了同步交换 \code{ids} 和 \code{names}，数据就全乱了。新增一个字段（比如年龄）还要再开一个数组，处处要同步。

结构体把“一个学生的所有信息”绑在一起：

\begin{codebox}{struct-student.c —— 用结构体存学生信息}
#include <stdio.h>

struct Student {
    int id;
    char name[20];
    double score;
};

int main(void)
{
    struct Student s;

    s.id = 1001;
    s.score = 92.5;
    /* s.name 是字符数组，赋值方式第 9 章讲，这里先用初始化 */

    printf("%d %.1f\n", s.id, s.score);
    return 0;
}
\end{codebox}

\begin{figure}[htbp]
  \centering
  \begin{tikzpicture}[x=1cm,y=1cm]
    % 平行数组
    \node[font=\small\bfseries] at (0,3.0) {平行数组};
    \node[membox,minimum width=2.4cm] (id0) at (0,2.2) {1001};
    \node[membox,minimum width=2.4cm,right=0mm of id0] (id1) {1002};
    \node[membox,minimum width=2.4cm,right=0mm of id1] (id2) {1003};
    \node[font=\scriptsize,color=black!55,left=1mm of id0] {ids};

    \node[membox,minimum width=2.4cm] (nm0) at (0,1.3) {Alice};
    \node[membox,minimum width=2.4cm,right=0mm of nm0] (nm1) {Bob};
    \node[membox,minimum width=2.4cm,right=0mm of nm1] (nm2) {Carol};
    \node[font=\scriptsize,color=black!55,left=1mm of nm0] {names};

    \node[membox,minimum width=2.4cm] (sc0) at (0,0.4) {92.5};
    \node[membox,minimum width=2.4cm,right=0mm of sc0] (sc1) {78.0};
    \node[membox,minimum width=2.4cm,right=0mm of sc1] (sc2) {85.5};
    \node[font=\scriptsize,color=black!55,left=1mm of sc0] {scores};

    % 结构体数组
    \node[font=\small\bfseries] at (9.2,3.0) {结构体数组};
    \node[draw=structurecolor!70,fill=structurecolor!6,rounded corners=2mm,
          minimum width=3.2cm,minimum height=1.1cm,align=left,font=\scriptsize] (st0) at (9.2,2.2)
      {id: 1001\\name: Alice\\score: 92.5};
    \node[draw=structurecolor!70,fill=structurecolor!6,rounded corners=2mm,
          minimum width=3.2cm,minimum height=1.1cm,align=left,font=\scriptsize,below=2mm of st0] (st1)
      {id: 1002\\name: Bob\\score: 78.0};
    \node[draw=structurecolor!70,fill=structurecolor!6,rounded corners=2mm,
          minimum width=3.2cm,minimum height=1.1cm,align=left,font=\scriptsize,below=2mm of st1] (st2)
      {id: 1003\\name: Carol\\score: 85.5};
    \node[font=\scriptsize,color=black!55,right=2mm of st1] {每人一“包”};
  \end{tikzpicture}
  \caption{平行数组 vs.\ 结构体数组：后者把一个人的信息绑在一起}
  \label{fig:struct-vs-parallel}
\end{figure}

\section{结构体的定义与使用}

\subsection{定义}

\begin{center}
\fbox{\parbox{0.82\textwidth}{
\ttfamily struct 结构体名 \{ \\
\phantom{xxxx}类型 成员1; \\
\phantom{xxxx}类型 成员2; \\
\phantom{xxxx}... \\
\ttfamily \}; \quad /* 注意这个分号 */
}}
\end{center}

\begin{codebox}{struct-def.c —— 定义一个学生结构体}
struct Student {
    int id;             /* 学号 */
    char name[20];      /* 姓名 */
    double score;       /* 成绩 */
};
\end{codebox}

\begin{pitfall}[结构体定义末尾的分号不能少]
\texttt{struct Student \{ ... \};} 最后的分号是定义的一部分，漏掉会导致一连串奇怪的编译错误。这个分号很容易忘，因为函数定义后面不写分号。
\end{pitfall}

\code{struct Student} 是一个\textbf{类型名}，就像 \code{int} 一样。注意：\code{Student} 单独使用不合法，必须写 \code{struct Student}，除非用 \code{typedef} 起别名（见后文）。

\subsection{声明变量与访问成员}

\begin{codebox}{struct-use.c —— 声明、初始化、访问成员}
#include <stdio.h>

struct Student {
    int id;
    char name[20];
    double score;
};

int main(void)
{
    struct Student s1;
    struct Student s2 = {1002, "Bob", 78.0};

    s1.id = 1001;
    s1.score = 92.5;

    printf("s1: %d, %.1f\n", s1.id, s1.score);
    printf("s2: %d, %s, %.1f\n", s2.id, s2.name, s2.score);

    return 0;
}
\end{codebox}

\begin{outbox}{运行结果}
s1: 1001, 92.5
s2: 1002, Bob, 78.0
\end{outbox}

访问成员用点号 \code{.}，读作“的”：\code{s1.id} 就是“s1 的 id”。

\subsection{整体赋值}

和数组不同，\textbf{同类型的结构体变量可以整体赋值}，编译器会把所有成员逐字节复制过去：

\begin{codebox}{struct-assign.c —— 结构体整体赋值}
#include <stdio.h>

struct Point {
    int x;
    int y;
};

int main(void)
{
    struct Point p1 = {3, 4};
    struct Point p2;

    p2 = p1;        /* 合法：把 p1 的所有成员复制给 p2 */

    printf("p2 = (%d, %d)\n", p2.x, p2.y);
    return 0;
}
\end{codebox}

\begin{outbox}{运行结果}
p2 = (3, 4)
\end{outbox}

\subsection{\texttt{typedef}：给类型起别名}

每次写 \code{struct Student} 有点啰嗦。可以用 \code{typedef} 起一个别名：

\begin{codebox}{typedef-struct.c —— typedef 的两种写法}
/* 写法一：先定义，再起别名 */
struct Student {
    int id;
    char name[20];
    double score;
};
typedef struct Student Student;

/* 写法二：定义的同时起别名（更常用） */
typedef struct {
    int x;
    int y;
} Point;
\end{codebox}

之后就可以直接写 \code{Student s;}、\code{Point p;}，和 \code{int a;} 一样自然。

\begin{tipbox}[本书的风格]
本书在需要频繁使用某个结构体时，会用 \code{typedef} 起一个首字母大写的别名；只在定义处出现一次的结构体，可以直接用 \code{struct} 写法。
\end{tipbox}

\subsection{嵌套结构体}

结构体的成员也可以是另一个结构体：

\begin{codebox}{nested-struct.c —— 嵌套结构体}
#include <stdio.h>

typedef struct {
    int year;
    int month;
    int day;
} Date;

typedef struct {
    int id;
    char name[20];
    Date birthday;      /* 成员是另一个结构体 */
    double score;
} Student;

int main(void)
{
    Student s = {1001, "Alice", {2005, 6, 18}, 92.5};

    printf("%s 出生于 %d 年 %d 月 %d 日\n",
           s.name, s.birthday.year, s.birthday.month, s.birthday.day);

    return 0;
}
\end{codebox}

\begin{outbox}{运行结果}
Alice 出生于 2005 年 6 月 18 日
\end{outbox}

访问嵌套成员就一层层点下去：\code{s.birthday.year}。

\section{结构体在内存里长什么样}

结构体的成员按声明顺序依次存放，但\textbf{不一定紧挨着}——编译器会在成员之间插入“填充字节”（padding），让每个成员的地址满足对齐要求。

\begin{figure}[htbp]
  \centering
  \begin{tikzpicture}[x=1cm,y=1cm]
    % 标尺
    \foreach \i in {0,...,15} {
      \node[font=\scriptsize,color=black!45] at (\i*0.9+0.45,1.3) {\i};
    }
    % 成员
    \draw[structurecolor!70,fill=ptrfill] (0,0) rectangle (0.9,-1.2);
    \node[font=\scriptsize] at (0.45,-0.6) {a};
    \draw[black!15,fill=black!6] (0.9,0) rectangle (3.6,-1.2);
    \node[font=\scriptsize,color=black!45] at (2.25,-0.6) {填充};
    \draw[structurecolor!70,fill=structurecolor!12] (3.6,0) rectangle (7.2,-1.2);
    \node[font=\scriptsize] at (5.4,-0.6) {b};
    \draw[structurecolor!70,fill=okfill] (7.2,0) rectangle (14.4,-1.2);
    \node[font=\scriptsize] at (10.8,-0.6) {c};

    \node[font=\scriptsize,color=black!55,below] at (0.45,-1.2) {偏移 0};
    \node[font=\scriptsize,color=black!55,below] at (5.4,-1.2) {偏移 4};
    \node[font=\scriptsize,color=black!55,below] at (10.8,-1.2) {偏移 8};
    \node[font=\scriptsize,color=black!55,below] at (14.4,-1.2) {总大小 16};
  \end{tikzpicture}
  \caption{\texttt{struct \{ char a; int b; double c; \}} 的内存布局（典型 8 字节对齐）}
  \label{fig:struct-padding}
\end{figure}

为什么需要对齐？因为 CPU 读取内存时，按“字”对齐的访问速度最快，某些架构甚至不允许非对齐访问。编译器为了性能，会在成员之间补空。

\begin{codebox}{sizeof-struct.c —— 结构体的大小}
#include <stdio.h>

struct A {
    char a;
    int b;
    double c;
};

struct B {
    char a;
    char b;
    int c;
};

int main(void)
{
    printf("sizeof(struct A) = %d\n", (int)sizeof(struct A));
    printf("sizeof(struct B) = %d\n", (int)sizeof(struct B));
    return 0;
}
\end{codebox}

\begin{outbox}{运行结果（典型 64 位环境）}
sizeof(struct A) = 16
sizeof(struct B) = 8
\end{outbox}

\code{struct A} 的成员一共 $1 + 4 + 8 = 13$ 字节，但因为对齐，实际占 16 字节。\code{struct B} 的两个 \code{char} 挤在一起，后面补 2 字节再放 \code{int}，共 8 字节。

\begin{tipbox}[省内存的小技巧]
把成员按“从大到小”排列，通常能减少填充。例如 \texttt{struct \{double c; int b; char a;\}} 在很多平台上只占 16 字节，而 \texttt{\{char a; int b; double c;\}} 也是 16——但当成员更多时，差别会明显。
\end{tipbox}

\begin{pitfall}[不要把结构体的内存布局当成协议]
虽然你可以用 \code{memcpy} 把结构体直接写进文件或发到网络，但填充字节的内容是不确定的，不同编译器/平台的对齐规则也不同。跨平台传输数据时，应该逐字段序列化，而不是直接发结构体。
\end{pitfall}

\section{结构体数组}

结构体数组就是“每个元素都是一个结构体”的数组：

\begin{codebox}{struct-array.c —— 结构体数组}
#include <stdio.h>

typedef struct {
    int id;
    char name[20];
    double score;
} Student;

int main(void)
{
    Student cls[3] = {
        {1001, "Alice", 92.5},
        {1002, "Bob",   78.0},
        {1003, "Carol", 85.5}
    };
    int i;
    double sum = 0.0;

    for (i = 0; i < 3; i++) {
        printf("%d %-8s %.1f\n", cls[i].id, cls[i].name, cls[i].score);
        sum = sum + cls[i].score;
    }

    printf("平均分 = %.2f\n", sum / 3.0);
    return 0;
}
\end{codebox}

\begin{outbox}{运行结果}
1001 Alice    92.5
1002 Bob      78.0
1003 Carol    85.5
平均分 = 85.33
\end{outbox}

\begin{note}
\texttt{\%-8s} 表示字符串左对齐、占 8 格，用来让姓名列对齐。这个格式符第 9 章会详细讲。
\end{note}

\section{结构体与函数}

结构体可以像普通变量一样传给函数、从函数返回。这里先讲最基本的用法，函数的完整语法见附录~\ref{app:function}。

\subsection{按值传递：函数拿到一份副本}

\begin{codebox}{struct-param.c —— 结构体作为参数}
#include <stdio.h>

typedef struct {
    int id;
    char name[20];
    double score;
} Student;

void print_student(Student s)
{
    printf("%d %s %.1f\n", s.id, s.name, s.score);
}

int main(void)
{
    Student s = {1001, "Alice", 92.5};

    print_student(s);       /* 传递的是一份副本 */
    return 0;
}
\end{codebox}

按值传递时，函数里修改 \code{s} 不会影响 \code{main} 里的原变量。这很安全，但如果结构体很大，复制会浪费时间。

\subsection{按地址传递：函数可以修改原变量}

如果想让函数修改原结构体，或者结构体很大想避免复制，就传地址：

\begin{codebox}{struct-pointer.c —— 传地址修改结构体}
#include <stdio.h>

typedef struct {
    int id;
    double score;
} Student;

void set_score(Student *p, double new_score)
{
    p->score = new_score;       /* 等价于 (*p).score */
}

int main(void)
{
    Student s = {1001, 92.5};

    set_score(&s, 95.0);        /* 传地址 */
    printf("%.1f\n", s.score);
    return 0;
}
\end{codebox}

\begin{outbox}{运行结果}
95.0
\end{outbox}

\code{p->score} 是“指针访问成员”的写法，等价于 \code{(*p).score}。下一章学指针时会详细解释，现在先记住这个符号。

\section{综合例题}

\subsection{例：学生成绩统计}

\begin{example}
有 5 个学生，每个学生有学号、姓名、成绩。输出平均分、最高分学生，以及高于平均分的人数。
\end{example}
\begin{solution}
思路如下：

\begin{codebox}{student-stats.c —— 学生成绩统计}
#include <stdio.h>

typedef struct {
    int id;
    char name[20];
    double score;
} Student;

int main(void)
{
    Student cls[5] = {
        {1001, "Alice", 92.5},
        {1002, "Bob",   78.0},
        {1003, "Carol", 85.5},
        {1004, "Dave",  66.0},
        {1005, "Eve",   95.0}
    };
    int i;
    double sum = 0.0;
    double average;
    int max_index = 0;
    int above = 0;

    for (i = 0; i < 5; i++) {
        sum = sum + cls[i].score;
    }
    average = sum / 5.0;

    for (i = 1; i < 5; i++) {
        if (cls[i].score > cls[max_index].score) {
            max_index = i;
        }
    }

    for (i = 0; i < 5; i++) {
        if (cls[i].score > average) {
            above++;
        }
    }

    printf("平均分 = %.2f\n", average);
    printf("最高分：%s (%.1f)\n",
           cls[max_index].name, cls[max_index].score);
    printf("高于平均分的人数：%d\n", above);

    return 0;
}
\end{codebox}

\begin{outbox}{运行结果}
平均分 = 83.40
最高分：Eve (95.0)
高于平均分的人数：2
\end{outbox}
\end{solution}

\subsection{例：复数运算}

\begin{example}
定义复数结构体，输入两个复数，输出它们的和与积。
\end{example}
\begin{solution}
复数 $a + bi$ 与 $c + di$ 相加为 $(a+c) + (b+d)i$，相乘为 $(ac-bd) + (ad+bc)i$。

\begin{codebox}{complex.c —— 复数运算}
#include <stdio.h>

typedef struct {
    double re;      /* 实部 */
    double im;      /* 虚部 */
} Complex;

int main(void)
{
    Complex x, y, sum, product;

    printf("请输入第一个复数（实部 虚部）：");
    scanf("%lf%lf", &x.re, &x.im);
    printf("请输入第二个复数（实部 虚部）：");
    scanf("%lf%lf", &y.re, &y.im);

    sum.re = x.re + y.re;
    sum.im = x.im + y.im;

    product.re = x.re * y.re - x.im * y.im;
    product.im = x.re * y.im + x.im * y.re;

    printf("和 = %.2f + %.2fi\n", sum.re, sum.im);
    printf("积 = %.2f + %.2fi\n", product.re, product.im);

    return 0;
}
\end{codebox}

\begin{outbox}{运行示例}
请输入第一个复数（实部 虚部）：1 2
请输入第二个复数（实部 虚部）：3 4
和 = 4.00 + 6.00i
积 = -5.00 + 10.00i
\end{outbox}
\end{solution}

\subsection{例：日期与下一天}

\begin{example}
定义日期结构体，输入一个日期，输出它的下一天。要考虑闰年和月份天数。
\end{example}
\begin{solution}
思路如下：

\begin{codebox}{next-day.c —— 下一天}
#include <stdio.h>

typedef struct {
    int year;
    int month;
    int day;
} Date;

int is_leap(int year)
{
    return (year % 4 == 0 && year % 100 != 0) || (year % 400 == 0);
}

int main(void)
{
    Date d;
    int days[13] = {0, 31, 28, 31, 30, 31, 30, 31, 31, 30, 31, 30, 31};

    printf("请输入日期（年 月 日）：");
    scanf("%d%d%d", &d.year, &d.month, &d.day);

    if (is_leap(d.year)) {
        days[2] = 29;
    }

    d.day++;
    if (d.day > days[d.month]) {
        d.day = 1;
        d.month++;
        if (d.month > 12) {
            d.month = 1;
            d.year++;
        }
    }

    printf("下一天：%d-%02d-%02d\n", d.year, d.month, d.day);
    return 0;
}
\end{codebox}

\begin{outbox}{运行示例}
请输入日期（年 月 日）：2024 2 28
下一天：2024-02-29
\end{outbox}

这里用数组 \code{days} 存每月天数，下标 1\textasciitilde12 对应 1\textasciitilde12 月，\code{days[0]} 占位不用。
\end{solution}

\subsection{例：电机参数}

\begin{example}
定义一个电机参数结构体，包含编号、PID 三个参数和输出限幅；输入一组参数并打印。
\end{example}
\begin{solution}
思路如下：

\begin{codebox}{motor.c —— 电机参数结构体}
#include <stdio.h>

typedef struct {
    int id;
    double kp;
    double ki;
    double kd;
    double limit;
} MotorParam;

int main(void)
{
    MotorParam m;

    printf("请输入电机编号与 PID 参数：");
    scanf("%d%lf%lf%lf", &m.id, &m.kp, &m.ki, &m.kd);

    printf("请输入输出限幅：");
    scanf("%lf", &m.limit);

    printf("电机 %d: kp=%.3f, ki=%.3f, kd=%.3f, limit=%.2f\n",
           m.id, m.kp, m.ki, m.kd, m.limit);

    return 0;
}
\end{codebox}

结构体让一组相关参数“打包”在一起，后面写 PID 函数时，传一个 \code{MotorParam} 就够了，不用传五个参数。
\end{solution}

\section{常见错误}

\begin{center}
\begin{tabular}{@{}p{0.38\textwidth}p{0.52\textwidth}@{}}
\toprule
错误 & 问题 \\
\midrule
结构体定义后漏分号 & 后续代码出现莫名其妙的错误 \\
\code{Student s;} 但没用 \code{typedef} & 应写 \code{struct Student s;} \\
\code{s == t} & 结构体不能直接用 \code{==} 比较 \\
\code{s.name = "Bob"} & 数组不能整体赋值，要用 \code{strcpy} \\
\texttt{scanf("\%s", \&s.name)} & \code{s.name} 本身就是地址，不要加 \texttt{\&} \\
忘记给数值成员加 \texttt{\&} & \code{scanf} 写不进去 \\
\bottomrule
\end{tabular}
\end{center}

\begin{pitfall}[结构体比较要逐成员]
\code{if (s1 == s2)} 在 C 里不合法。正确做法是逐成员比较，或者用 \code{memcmp}——但 \code{memcmp} 会把填充字节也比进去，不可靠。最稳妥的是逐成员比较：
\begin{codebox}{struct-compare.c —— 逐成员比较}
int equal(Student a, Student b)
{
    return a.id == b.id &&
           a.score == b.score;   /* 字符串比较需用 strcmp */
}
\end{codebox}
\end{pitfall}

\begin{pitfall}[字符串成员不能直接赋值]
\code{s.name = "Bob";} 是错的，因为 \code{name} 是数组，数组名不能做赋值目标。正确写法是 \code{strcpy(s.name, "Bob");}，第 9 章会讲。
\end{pitfall}

\section{拔高：位域与联合体}

\subsection{位域}

有时结构体成员只需要几个二进制位（例如标志位），用 \code{int} 会浪费空间。C 语言提供了\textbf{位域}：

\begin{codebox}{bitfield.c —— 位域}
#include <stdio.h>

struct Flags {
    unsigned int enable : 1;    /* 占 1 位 */
    unsigned int mode   : 2;    /* 占 2 位，取值 0-3 */
    unsigned int speed  : 5;    /* 占 5 位，取值 0-31 */
};

int main(void)
{
    struct Flags f;

    f.enable = 1;
    f.mode = 3;
    f.speed = 20;

    printf("enable=%u mode=%u speed=%u\n", f.enable, f.mode, f.speed);
    printf("sizeof(struct Flags) = %d\n", (int)sizeof(struct Flags));
    return 0;
}
\end{codebox}

位域在嵌入式寄存器映射中很常见，但它的具体布局（谁在高位、跨字节怎么排）由编译器决定，跨平台代码要谨慎。

\subsection{联合体}

\code{union} 让多个成员\textbf{共用同一块内存}，同一时刻只有一个成员有意义：

\begin{codebox}{union.c —— 联合体}
#include <stdio.h>

union Data {
    int i;
    float f;
    unsigned char bytes[4];
};

int main(void)
{
    union Data d;

    d.i = 0x12345678;
    printf("i = %x\n", d.i);
    printf("bytes = %02x %02x %02x %02x\n",
           d.bytes[0], d.bytes[1], d.bytes[2], d.bytes[3]);
    printf("sizeof(union Data) = %d\n", (int)sizeof(union Data));

    return 0;
}
\end{codebox}

联合体常用于：
\begin{itemize}
  \item 查看一个数据的字节表示（例如上面的小端序演示）；
  \item 节省内存：同一块空间在不同时刻存不同类型；
  \item 协议解析：把一帧数据按不同格式解释。
\end{itemize}

\begin{note}
上面代码里 \code{d.i = 0x12345678} 后读 \code{d.bytes}，属于“用另一个成员读取”，在 C 标准里是允许的（类型双关），但不同电脑的存放顺序可能不同。这里只用于演示，不要依赖它做可移植代码。
\end{note}

\section{本章小结}

\begin{center}
\begin{tabular}{@{}p{0.24\textwidth}p{0.66\textwidth}@{}}
\toprule
主题 & 要点 \\
\midrule
结构体 & 把不同类型的成员打包成一个整体 \\
定义 & \texttt{struct 名字 \{ ... \};}，末尾分号不能少 \\
访问 & 用 \code{.} 访问成员；整体赋值是逐字节复制 \\
\code{typedef} & 给类型起别名，少写 \code{struct} \\
嵌套 & 成员可以是另一个结构体，逐层点访问 \\
数组 & 结构体数组适合管理一批对象 \\
函数 & 按值传递复制一份；传地址用 \code{->} 修改原值 \\
内存 & 成员按声明顺序排列，可能有填充字节 \\
拔高 & 位域省空间；联合体共用内存 \\
\bottomrule
\end{tabular}
\end{center}

\begin{tipbox}[下一章预告]
结构体让我们能“描述对象”，但还有一个问题没解决：如何让函数直接修改调用者的变量？如何让一个变量“指向”另一个变量？如何高效地传递大块数据？答案就是\textbf{指针}——C 语言最锋利、也最容易伤到自己的工具。
\end{tipbox}

\begin{problemset}
  \item 定义一个 \code{Point} 结构体（\code{x}、\code{y}），输入两个点，输出它们的中点。
  \item 定义一个 \code{Rect} 结构体（左上角、右下角），计算面积和周长。
  \item 定义 \code{Student} 结构体，输入 5 个学生信息，输出平均分和最高分学生。
  \item 用结构体数组实现：输入 5 个学生的成绩，按输入顺序输出，并统计及格人数。
  \item 定义 \code{Date} 结构体，判断两个日期哪个更早。
  \item 定义 \code{Complex} 结构体，实现复数的加法、减法、乘法。
  \item 定义一个包含 \code{char}、\code{int}、\code{double} 的结构体，用 \code{sizeof} 观察它的大小，并解释为什么不是各成员大小之和。
  \item 写一个函数，接收 \code{Student} 的地址，把成绩加 5 分（不超过 100）。
  \item 用位域定义一个 8 位的状态字，包含 1 位使能、2 位模式、5 位速度，并打印各字段。
  \item 思考题：为什么结构体可以整体赋值，而数组不可以？
\end{problemset}

\chapter{指针：会者不难}

\begin{introduction}
  \item 地址、指针、取地址与解引用
  \item 指针与变量、数组、结构体的关系
  \item 指针作为函数参数：让函数修改调用者的数据
  \item 指针运算与数组遍历
  \item 动态内存：\texttt{malloc} 与 \texttt{free}
  \item 二级指针、函数指针与常见陷阱
\end{introduction}

\chaptermotto{「指针不是魔法，它只是“存地址的变量”。」}

很多同学一听到“指针”就紧张。其实指针的概念非常朴素：\textbf{一个变量，里面存的是另一个变量的地址}。你已经在第 1 章见过地址，在第 2 章用 \code{scanf} 时用过 \texttt{\&}，在第 6 章知道数组元素连续存放。指针只是把这些知识串起来。

这一章会慢一点讲：先把“地址—指针—解引用”这三件事用图讲清楚，再谈指针和数组、函数、结构体的配合，最后讲动态内存和常见陷阱。请务必亲手画几遍内存图。

\section{地址与指针}

\subsection{变量住在哪里}

第 1 章说过：每个变量都占一块内存，每块内存有地址。例如：

\begin{codebox}{address.c —— 打印变量的地址}
#include <stdio.h>

int main(void)
{
    int a = 42;

    printf("a 的值   = %d\n", a);
    printf("a 的地址 = %p\n", (void *)&a);

    return 0;
}
\end{codebox}

\begin{outbox}{运行结果（地址每次可能不同）}
a 的值   = 42
a 的地址 = 000000000061FE1C
\end{outbox}

\texttt{\&a} 表示“取 \code{a} 的地址”。你可以把地址理解成酒店房间号：房间号不是房间本身，但拿着房间号就能找到房间。

\subsection{指针变量}

\textbf{指针}就是一个专门用来存地址的变量。定义方式：

\begin{center}
\fbox{\parbox{0.6\textwidth}{
\ttfamily 类型 *指针名;
}}
\end{center}

例如 \code{int *p;} 表示“\code{p} 是一个指针，它指向一个 \code{int} 类型的数据”。

\begin{codebox}{pointer-basic.c —— 指针的基本用法}
#include <stdio.h>

int main(void)
{
    int a = 42;
    int *p;         /* p 是指向 int 的指针 */

    p = &a;         /* 把 a 的地址存进 p */

    printf("a 的地址 = %p\n", (void *)&a);
    printf("p 的值   = %p\n", (void *)p);
    printf("p 指向的值 = %d\n", *p);

    return 0;
}
\end{codebox}

\begin{outbox}{运行结果（地址每次可能不同）}
a 的地址 = 000000000061FE1C
p 的值   = 000000000061FE1C
p 指向的值 = 42
\end{outbox}

这里出现了两个新符号：

\begin{itemize}
  \item \texttt{\&}：取地址。\texttt{\&a} 得到 \code{a} 的地址；
  \item \code{*}：解引用（间接访问）。\code{*p} 表示“\code{p} 指向的那个变量”。
\end{itemize}

\begin{figure}[htbp]
  \centering
  \begin{tikzpicture}[x=1cm,y=1cm,>=Stealth]
    % a
    \node[draw=memline,fill=memfill,minimum width=2.2cm,minimum height=1.2cm,align=center] (a) at (0,0) {\texttt{a}\\42};
    \node[font=\scriptsize,color=black!55,above=1mm of a] {地址 0x1000};
    % p
    \node[draw=ptrline,fill=ptrfill,minimum width=2.2cm,minimum height=1.2cm,align=center] (p) at (5.2,0) {\texttt{p}\\0x1000};
    \node[font=\scriptsize,color=black!55,above=1mm of p] {地址 0x2000};
    \draw[ptr,line width=1pt] (p.west) -- node[lbl,above]{存的是 a 的地址} (a.east);
    \node[font=\small,color=ptrline,below=4mm of a] {变量 a};
    \node[font=\small,color=ptrline,below=4mm of p] {指针 p};
  \end{tikzpicture}
  \caption{指针 \texttt{p} 里存的是变量 \texttt{a} 的地址；通过 \texttt{*p} 可以访问 \texttt{a}}
  \label{fig:pointer-basic}
\end{figure}

\begin{tipbox}[一句话记住]
\code{p} 是地址，\code{*p} 是地址里的东西。就像 \code{p} 是房间号，\code{*p} 是房间里的行李。
\end{tipbox}

\subsection{通过指针修改变量}

有了指针，就可以“隔着地址”修改原变量：

\begin{codebox}{pointer-modify.c —— 通过指针修改}
#include <stdio.h>

int main(void)
{
    int a = 42;
    int *p = &a;

    *p = 100;       /* 等价于 a = 100 */
    printf("a = %d\n", a);
    printf("*p = %d\n", *p);

    return 0;
}
\end{codebox}

\begin{outbox}{运行结果}
a = 100
*p = 100
\end{outbox}

这看起来多此一举，但它是很多高级技巧的基础：函数可以通过指针修改调用者的变量，数据结构可以用指针连接起来。

\subsection{空指针与野指针}

\begin{itemize}
  \item \textbf{空指针}：不指向任何有效地址，用 \code{NULL} 表示；
  \item \textbf{野指针}：指向一个不确定/已失效地址的指针，极其危险。
\end{itemize}

\begin{codebox}{null-pointer.c —— 空指针}
#include <stdio.h>
#include <stdlib.h>

int main(void)
{
    int *p = NULL;

    if (p != NULL) {
        printf("%d\n", *p);
    } else {
        printf("p 是空指针，不能解引用\n");
    }

    return 0;
}
\end{codebox}

\begin{pitfall}[野指针]
\begin{codebox}{wild-pointer.c —— 野指针}
int *p;         /* 未初始化，p 里是垃圾值 */
*p = 10;        /* 往一个随机地址写数据：崩溃或更糟 */
\end{codebox}
指针在使用前必须初始化：要么指向一个有效变量，要么赋值为 \code{NULL}。
\end{pitfall}

\section{指针与函数}

\subsection{让函数修改调用者的变量}

C 语言的函数参数是“按值传递”的：函数拿到的是实参的副本。所以下面这个交换函数是\textbf{错的}：

\begin{codebox}{swap-wrong.c —— 交换失败}
#include <stdio.h>

void swap(int x, int y)
{
    int t;
    t = x;
    x = y;
    y = t;
}

int main(void)
{
    int a = 3;
    int b = 5;

    swap(a, b);
    printf("a = %d, b = %d\n", a, b);   /* 还是 3 和 5 */
    return 0;
}
\end{codebox}

\begin{outbox}{运行结果}
a = 3, b = 5
\end{outbox}

因为 \code{swap} 里的 \code{x}、\code{y} 只是 \code{a}、\code{b} 的副本，交换副本不影响原件。要修改原件，就得把地址传进去：

\begin{codebox}{swap-pointer.c —— 用指针交换}
#include <stdio.h>

void swap(int *px, int *py)
{
    int t;

    t = *px;
    *px = *py;
    *py = t;
}

int main(void)
{
    int a = 3;
    int b = 5;

    swap(&a, &b);       /* 传地址 */
    printf("a = %d, b = %d\n", a, b);
    return 0;
}
\end{codebox}

\begin{outbox}{运行结果}
a = 5, b = 3
\end{outbox}

\begin{figure}[htbp]
  \centering
  \begin{tikzpicture}[x=1cm,y=1cm,>=Stealth]
    \node[draw=memline,fill=memfill,minimum width=2cm,minimum height=1cm,align=center] (a) at (0,0) {\texttt{a}\\3};
    \node[draw=memline,fill=memfill,minimum width=2cm,minimum height=1cm,align=center] (b) at (6,0) {\texttt{b}\\5};
    \node[draw=ptrline,fill=ptrfill,minimum width=2cm,minimum height=1cm,align=center] (px) at (0,-2.2) {\texttt{px}};
    \node[draw=ptrline,fill=ptrfill,minimum width=2cm,minimum height=1cm,align=center] (py) at (6,-2.2) {\texttt{py}};
    \draw[ptr] (px.north) -- (a.south);
    \draw[ptr] (py.north) -- (b.south);
    \node[font=\small,below=3mm of px] {指向 a};
    \node[font=\small,below=3mm of py] {指向 b};
  \end{tikzpicture}
  \caption{传地址后，函数通过指针直接修改 \texttt{a} 和 \texttt{b}}
  \label{fig:swap-pointer}
\end{figure}

\subsection{指针参数与数组参数}

数组作为函数参数时，会“退化”成指向首元素的指针。所以下面两种写法等价：

\begin{codebox}{array-param.c —— 数组参数就是指针}
void print_array(int a[], int n);   /* 写法一 */
void print_array(int *a, int n);    /* 写法二，完全等价 */
\end{codebox}

在函数内部，\code{sizeof(a)} 得到的是指针大小，不是数组大小。所以长度必须额外传进来：

\begin{codebox}{print-array.c —— 打印数组}
#include <stdio.h>

void print_array(int *a, int n)
{
    int i;
    for (i = 0; i < n; i++) {
        printf("%d ", a[i]);
    }
    printf("\n");
}

int main(void)
{
    int arr[5] = {1, 2, 3, 4, 5};

    print_array(arr, 5);        /* 数组名就是首元素地址 */
    return 0;
}
\end{codebox}

\begin{outbox}{运行结果}
1 2 3 4 5
\end{outbox}

\section{指针与数组}

\subsection{数组名的含义}

数组名在大多数表达式里会转换成“指向首元素的指针”，即：

\[
\texttt{a} \quad \text{等价于} \quad \texttt{\&a[0]}.
\]

所以下面四种写法访问的是同一个元素：

\begin{center}
\begin{tabular}{@{}ll@{}}
\toprule
写法 & 含义 \\
\midrule
\code{a[i]} & 下标访问 \\
\code{*(a + i)} & 先算地址，再解引用 \\
\code{*(p + i)} & \code{p} 指向 \code{a[0]} 时 \\
\code{p[i]} & 指针也可以下标访问 \\
\bottomrule
\end{tabular}
\end{center}

\begin{codebox}{pointer-array.c —— 指针遍历数组}
#include <stdio.h>

int main(void)
{
    int a[5] = {10, 20, 30, 40, 50};
    int *p = a;         /* p 指向 a[0] */
    int i;

    for (i = 0; i < 5; i++) {
        printf("a[%d] = %d, *(p+%d) = %d\n",
               i, a[i], i, *(p + i));
    }

    return 0;
}
\end{codebox}

\begin{outbox}{运行结果}
a[0] = 10, *(p+0) = 10
a[1] = 20, *(p+1) = 20
a[2] = 30, *(p+2) = 30
a[3] = 40, *(p+3) = 40
a[4] = 50, *(p+4) = 50
\end{outbox}

\begin{figure}[htbp]
  \centering
  \begin{tikzpicture}[x=1cm,y=1cm,>=Stealth]
    \foreach \i/\v in {0/10,1/20,2/30,3/40,4/50} {
      \node[membox,minimum width=1.5cm] (a\i) at (\i*1.7,0) {\v};
      \node[font=\scriptsize,color=black!55,below=1mm of a\i] {\texttt{a[\i]}};
    }
    \node[ptrbox,minimum width=1.4cm] (p) at (-2.4,-1.6) {\texttt{p}};
    \draw[ptr] (p.north) -- (a0.south);
    \node[ptrbox,minimum width=1.4cm] (p1) at (-0.7,-1.6) {\texttt{p+1}};
    \draw[ptr] (p1.north) -- (a1.south);
    \node[ptrbox,minimum width=1.4cm] (p2) at (1.0,-1.6) {\texttt{p+2}};
    \draw[ptr] (p2.north) -- (a2.south);
    \node[font=\small,color=ptrline] at (3.4,-1.6) {指针加 1，指向下一个元素};
  \end{tikzpicture}
  \caption{指针算术：\texttt{p+i} 指向第 \texttt{i} 个元素}
  \label{fig:pointer-arithmetic}
\end{figure}

\subsection{指针算术}

指针加 1，不是地址加 1 字节，而是加“一个元素的大小”：

\begin{itemize}
  \item \code{int *p}：\code{p + 1} 的地址增加 \code{sizeof(int)}（通常 4 字节）；
  \item \code{double *p}：\code{p + 1} 增加 8 字节；
  \item \code{char *p}：\code{p + 1} 增加 1 字节。
\end{itemize}

两个同类型指针相减，得到它们之间相隔的元素个数：

\begin{codebox}{pointer-diff.c —— 指针相减}
#include <stdio.h>

int main(void)
{
    int a[5] = {10, 20, 30, 40, 50};
    int *p = &a[1];
    int *q = &a[4];

    printf("q - p = %d\n", (int)(q - p));   /* 3 个元素 */

    return 0;
}
\end{codebox}

\begin{outbox}{运行结果}
q - p = 3
\end{outbox}

\subsection{数组名、首元素地址与整个数组的地址}

\code{a}、\texttt{\&a[0]}、\texttt{\&a} 的值通常相同，但类型不同：

\begin{itemize}
  \item \code{a} 和 \texttt{\&a[0]}：类型是 \code{int *}；
  \item \texttt{\&a}：类型是“指向整个数组的指针”，\texttt{\&a + 1} 会跳过整个数组。
\end{itemize}

\begin{codebox}{array-address.c —— 三种写法的区别}
#include <stdio.h>

int main(void)
{
    int a[5] = {1, 2, 3, 4, 5};

    printf("a     = %p\n", (void *)a);
    printf("&a[0] = %p\n", (void *)&a[0]);
    printf("&a    = %p\n", (void *)&a);
    printf("a + 1   = %p\n", (void *)(a + 1));
    printf("&a + 1  = %p\n", (void *)(&a + 1));

    return 0;
}
\end{codebox}

\begin{outbox}{运行结果（地址示意）}
a     = 000000000061FE00
&a[0] = 000000000061FE00
&a    = 000000000061FE00
a + 1   = 000000000061FE04
&a + 1  = 000000000061FE14
\end{outbox}

\code{a + 1} 前进了 4 字节（一个 \code{int}），\texttt{\&a + 1} 前进了 20 字节（整个数组）。初学阶段不必深究，知道有区别即可。

\section{指针与结构体}

指向结构体的指针用 \code{->} 访问成员：

\begin{codebox}{struct-pointer2.c —— 结构体指针}
#include <stdio.h>

typedef struct {
    int id;
    double score;
} Student;

int main(void)
{
    Student s = {1001, 92.5};
    Student *p = &s;

    printf("%d\n", p->id);          /* 等价于 (*p).id */
    p->score = 95.0;
    printf("%.1f\n", s.score);

    return 0;
}
\end{codebox}

\begin{outbox}{运行结果}
1001
95.0
\end{outbox}

结构体指针在链表、树、函数参数里非常常用。传结构体指针给函数，既避免复制，又能修改原结构体。

\section{二级指针与函数指针（拔高）}

\subsection{二级指针}

指针本身也有地址，所以可以有“指向指针的指针”：

\begin{codebox}{double-pointer.c —— 二级指针}
#include <stdio.h>

int main(void)
{
    int a = 5;
    int *p = &a;
    int **pp = &p;

    **pp = 10;          /* 通过二级指针修改 a */
    printf("a = %d\n", a);

    return 0;
}
\end{codebox}

\begin{outbox}{运行结果}
a = 10
\end{outbox}

二级指针常用于：函数需要修改“指针本身”（例如在函数里 \code{malloc} 并让调用者的指针指向新内存），或者二维数组、字符串数组。

\subsection{函数指针}

函数也有地址，可以定义指向函数的指针：

\begin{codebox}{function-pointer.c —— 函数指针}
#include <stdio.h>

int add(int a, int b)
{
    return a + b;
}

int main(void)
{
    int (*fp)(int, int);    /* fp 是指向“两个 int 参数、返回 int”的函数的指针 */

    fp = add;
    printf("%d\n", fp(2, 3));   /* 通过函数指针调用 */
    return 0;
}
\end{codebox}

\begin{outbox}{运行结果}
5
\end{outbox}

函数指针在回调、状态机、驱动表里很常见。初学阶段知道语法即可，不必急着用。

\section{动态内存}

\subsection{为什么需要动态内存}

数组的长度必须在编译时确定，如果运行时才知道要多少数据怎么办？用 \code{malloc} 在堆上动态申请内存。

\begin{figure}[htbp]
  \centering
  \begin{tikzpicture}[x=1cm,y=1cm]
    \fill[structurecolor!10] (0,0) rectangle (6,0.9);
    \fill[structurecolor!18] (0,1.1) rectangle (6,1.8);
    \fill[structurecolor!8] (0,2.0) rectangle (6,4.0);
    \fill[structurecolor!18] (0,4.2) rectangle (6,5.2);
    \node[font=\small\bfseries] at (3,0.45) {代码区};
    \node[font=\small\bfseries] at (3,1.45) {数据区};
    \node[font=\small\bfseries] at (3,3.0) {堆（malloc 申请）};
    \node[font=\small\bfseries] at (3,4.7) {栈（局部变量）};
    \node[font=\scriptsize,color=black!55,anchor=west] at (6.2,0.45) {机器指令};
    \node[font=\scriptsize,color=black!55,anchor=west] at (6.2,1.45) {全局/静态变量};
    \node[font=\scriptsize,color=black!55,anchor=west] at (6.2,3.0) {向上生长，手动释放};
    \node[font=\scriptsize,color=black!55,anchor=west] at (6.2,4.7) {向下生长，自动释放};
  \end{tikzpicture}
  \caption{堆与栈：\texttt{malloc} 在堆上分配，\texttt{free} 释放}
  \label{fig:heap-stack}
\end{figure}

\subsection{\texttt{malloc} 与 \texttt{free}}

\begin{codebox}{malloc.c —— 动态数组}
#include <stdio.h>
#include <stdlib.h>

int main(void)
{
    int n;
    int i;
    int *a;

    printf("请输入元素个数：");
    scanf("%d", &n);

    a = (int *)malloc(n * sizeof(int));
    if (a == NULL) {
        printf("内存分配失败\n");
        return 1;
    }

    printf("请输入 %d 个整数：", n);
    for (i = 0; i < n; i++) {
        scanf("%d", &a[i]);
    }

    for (i = 0; i < n; i++) {
        printf("%d ", a[i]);
    }
    printf("\n");

    free(a);        /* 用完必须释放 */
    a = NULL;       /* 防止野指针 */

    return 0;
}
\end{codebox}

要点：

\begin{itemize}
  \item \code{malloc(n * sizeof(int))} 申请 $n$ 个 \code{int} 的空间，返回首地址（\code{void *}）；
  \item 强制转换成 \code{(int *)} 是为了类型匹配；
  \item 申请可能失败，返回 \code{NULL}，必须检查；
  \item 用完调用 \code{free(a)} 释放；
  \item 释放后把指针置 \code{NULL}，避免“悬空指针”。
\end{itemize}

\begin{pitfall}[动态内存的四个大坑]
\begin{enumerate}
  \item \textbf{忘记 free}：内存泄漏，程序越跑越占内存；
  \item \textbf{重复 free}：未定义行为，可能崩溃；
  \item \textbf{free 后继续使用}：悬空指针，数据已被回收；
  \item \textbf{越界写}：\code{a[n]} 越界，破坏堆结构。
\end{enumerate}
\end{pitfall}

\subsection{\texttt{calloc} 与 \texttt{realloc}}

\begin{codebox}{calloc-realloc.c —— calloc 与 realloc}
/* calloc：申请 n 个元素，并自动清零 */
int *a = (int *)calloc(n, sizeof(int));

/* realloc：把已申请的空间调整为 new_n 个元素 */
int *b = (int *)realloc(a, new_n * sizeof(int));
if (b == NULL) {
    /* 调整失败，原来的 a 仍然有效 */
}
\end{codebox}

\section{指针的常见错误}

\begin{center}
\begin{tabular}{@{}p{0.4\textwidth}p{0.5\textwidth}@{}}
\toprule
错误 & 后果 \\
\midrule
使用未初始化的指针 & 野指针，崩溃或破坏数据 \\
\code{*p = 10;} 但 \code{p} 未指向有效内存 & 同上 \\
数组越界后继续用指针 & 未定义行为 \\
\code{free} 后继续使用 & 悬空指针 \\
重复 \code{free} & 崩溃 \\
忘记 \code{free} & 内存泄漏 \\
\code{p = 10;} & 类型错误：不能把整数赋给指针 \\
返回局部变量的地址 & 函数结束后地址失效 \\
\bottomrule
\end{tabular}
\end{center}

\begin{pitfall}[不要返回局部变量的地址]
\begin{codebox}{return-local.c —— 错误：返回局部变量地址}
int *bad(void)
{
    int x = 10;
    return &x;      /* x 在函数返回后就失效了 */
}
\end{codebox}
函数返回后，局部变量所在的栈空间会被回收，返回的地址指向一块不再有效的地方。
\end{pitfall}

\section{指针与 \texttt{const}}

\code{const} 和指针结合时，位置不同含义不同：

\begin{center}
\begin{tabular}{@{}ll@{}}
\toprule
写法 & 含义 \\
\midrule
\code{const int *p} & 指向的内容不能改，指针本身可以改 \\
\code{int * const p} & 指针本身不能改，指向的内容可以改 \\
\code{const int * const p} & 两者都不能改 \\
\bottomrule
\end{tabular}
\end{center}

\begin{codebox}{const-pointer.c —— const 指针}
void print(const int *p)
{
    /* *p = 10; 错误：不能通过 p 修改内容 */
    printf("%d\n", *p);
}
\end{codebox}

函数参数里用 \code{const int *} 表示“我只读，不改”，既清晰又安全。

\section{综合例题}

\begin{example}
用指针实现数组逆序。
\end{example}
\begin{solution}
思路如下：

\begin{codebox}{reverse-pointer.c —— 指针逆序}
#include <stdio.h>

int main(void)
{
    int a[10] = {1, 2, 3, 4, 5, 6, 7, 8, 9, 10};
    int *left = a;
    int *right = a + 9;
    int temp;

    while (left < right) {
        temp = *left;
        *left = *right;
        *right = temp;

        left++;
        right--;
    }

    {
        int i;
        for (i = 0; i < 10; i++) {
            printf("%d ", a[i]);
        }
        printf("\n");
    }

    return 0;
}
\end{codebox}

\begin{outbox}{运行结果}
10 9 8 7 6 5 4 3 2 1
\end{outbox}

在 C89 里不能在 \code{while} 后面临时声明 \code{i}，所以代码里用了一个花括号块，把 \code{i} 的声明放在块开头。
\end{solution}

\begin{example}
写一个函数，用指针求数组的最大值。
\end{example}
\begin{solution}
思路如下：

\begin{codebox}{max-pointer.c —— 指针求最大值}
#include <stdio.h>

int max_array(const int *a, int n)
{
    int i;
    int max = a[0];

    for (i = 1; i < n; i++) {
        if (a[i] > max) {
            max = a[i];
        }
    }
    return max;
}

int main(void)
{
    int a[5] = {3, 9, 1, 7, 5};

    printf("最大值 = %d\n", max_array(a, 5));
    return 0;
}
\end{codebox}
\end{solution}

\begin{example}
用 \texttt{malloc} 申请一个数组，读入 $n$ 个数，求平均值后释放。
\end{example}
\begin{solution}
思路如下：

\begin{codebox}{dynamic-average.c —— 动态数组求平均}
#include <stdio.h>
#include <stdlib.h>

int main(void)
{
    int n;
    int i;
    int *a;
    double sum = 0.0;

    printf("请输入个数：");
    scanf("%d", &n);

    a = (int *)malloc(n * sizeof(int));
    if (a == NULL) {
        printf("分配失败\n");
        return 1;
    }

    printf("请输入 %d 个数：", n);
    for (i = 0; i < n; i++) {
        scanf("%d", &a[i]);
        sum = sum + a[i];
    }

    printf("平均值 = %.2f\n", sum / n);

    free(a);
    a = NULL;
    return 0;
}
\end{codebox}
\end{solution}

\section{本章小结}

\begin{center}
\begin{tabular}{@{}p{0.24\textwidth}p{0.66\textwidth}@{}}
\toprule
主题 & 要点 \\
\midrule
指针 & 存地址的变量；\texttt{\&} 取地址，\code{*} 解引用 \\
类型 & 指针有类型，\code{int *} 和 \code{double *} 不同 \\
空指针 & \code{NULL}，不能解引用 \\
函数 & 传指针才能修改调用者的变量；数组参数退化为指针 \\
数组 & \code{a[i]} 等价于 \code{*(a+i)}；指针加 1 走一个元素 \\
结构体 & \code{p->member} 等价于 \code{(*p).member} \\
动态内存 & \code{malloc}/\code{free}；检查 \code{NULL}；释放后置空 \\
常见错误 & 野指针、悬空指针、内存泄漏、越界、返回局部地址 \\
\bottomrule
\end{tabular}
\end{center}

\begin{tipbox}[下一章预告]
指针和字符数组结合，就是 C 语言的字符串。字符串是电控里处理命令、日志、协议时最常用的数据形式。下一章我们会把字符、字符数组、字符串常量、指针之间的关系彻底理清。
\end{tipbox}

\begin{problemset}
  \item 定义变量 \code{int a = 10}，用指针输出 \code{a} 的值和地址，并通过指针把 \code{a} 改成 20。
  \item 写函数 \code{swap(int *a, int *b)} 交换两个整数，并在 \code{main} 中验证。
  \item 用指针遍历数组，输出所有元素。
  \item 用指针实现数组逆序。
  \item 用指针求数组的最大值、最小值和平均值。
  \item 写函数 \code{void add_score(Student *s, double x)}，给学生成绩加 \code{x} 分（不超过 100）。
  \item 用 \code{malloc} 申请 $n$ 个整数，读入并排序（可以用冒泡），然后释放。
  \item 解释 \code{a}、\texttt{\&a[0]}、\texttt{\&a} 的区别。
  \item 解释 \code{const int *p} 和 \code{int * const p} 的区别。
  \item 思考题：为什么函数参数用数组形式 \code{int a[]} 时，函数内部无法用 \code{sizeof} 求长度？
\end{problemset}

\chapter{字符串}

\begin{introduction}
  \item 字符、字符数组与字符串常量
  \item 字符串的结束标志 \texttt{\textbackslash0}
  \item \texttt{scanf}、\texttt{fgets}、\texttt{puts} 的输入输出
  \item \texttt{<string.h>} 常用函数：求长、复制、连接、比较、查找
  \item 用指针和数组处理字符串
  \item 中文字符编码与常见错误
\end{introduction}

\chaptermotto{「字符串是字符排成的队伍，\texttt{\textbackslash0} 是队伍的终点线。」}

C 语言没有专门的“字符串类型”。字符串本质上是一个\textbf{以 \texttt{\textbackslash0} 结尾的字符数组}，或者一个指向这种数组的指针。理解这一点，字符串的所有“怪现象”都能解释清楚。

电控程序里到处都是字符串：串口命令、日志、文件名、协议报文、调试信息。这一章会把你从“会打印 hello”带到“能自己写字符串处理函数”。

\section{字符与字符串}

\subsection{字符类型回顾}

\code{char} 存的是一个字符的编码（ASCII 码），本质是小整数：

\begin{codebox}{char-basic.c —— 字符与编码}
#include <stdio.h>

int main(void)
{
    char c = 'A';

    printf("字符：%c\n", c);
    printf("编码：%d\n", c);
    return 0;
}
\end{codebox}

\begin{outbox}{运行结果}
字符：A
编码：65
\end{outbox}

常用字符的 ASCII 码：

\begin{center}
\begin{tabular}{@{}llll@{}}
\toprule
字符 & 编码 & 字符 & 编码 \\
\midrule
\code{'0'} \textasciitilde \code{'9'} & 48 \textasciitilde 57 & \code{'A'} \textasciitilde \code{'Z'} & 65 \textasciitilde 90 \\
\code{'a'} \textasciitilde \code{'z'} & 97 \textasciitilde 122 & 空格 & 32 \\
\texttt{'\textbackslash n'} & 10 & \texttt{'\textbackslash0'} & 0 \\
\bottomrule
\end{tabular}
\end{center}

记住三组关系：

\begin{itemize}
  \item 数字字符 \code{'0'} 的编码是 48，所以 \code{'7' - '0'} 等于整数 7；
  \item 大写字母 \code{'A'} 是 65，小写字母 \code{'a'} 是 97，相差 32；
  \item 字符 \code{'0'} 和整数 0 完全不是一回事。
\end{itemize}

\subsection{字符串就是字符数组}

\begin{codebox}{string-array.c —— 字符串的内存}
#include <stdio.h>

int main(void)
{
    char s[6] = "hello";

    printf("%s\n", s);
    printf("sizeof(s) = %d\n", (int)sizeof(s));

    return 0;
}
\end{codebox}

\begin{outbox}{运行结果}
hello
sizeof(s) = 6
\end{outbox}

字符串 \code{"hello"} 有 5 个字符，但数组要 6 个字节，因为最后还有一个 \textbf{\texttt{\textbackslash0}}（空字符，编码 0）作为结束标志：

\begin{figure}[htbp]
  \centering
  \begin{tikzpicture}[x=1cm,y=1cm]
    \foreach \i/\c in {0/h,1/e,2/l,3/l,4/o,5/{$\backslash$0}} {
      \node[membox,minimum width=1.3cm] (c\i) at (\i*1.5,0) {\ttfamily\c};
      \node[font=\scriptsize,color=black!55,below=1mm of c\i] {\i};
    }
    \draw[decorate,decoration={brace,amplitude=4pt},structurecolor!70]
      (-0.65,0.7) -- (7.25,0.7) node[midway,above=4pt,font=\scriptsize]{6 个字节};
    \node[font=\small,color=ptrline,below=8mm] at (3.75,0) {\texttt{\textbackslash0} 告诉所有字符串函数：到这里为止};
  \end{tikzpicture}
  \caption{\texttt{char s[6] = "hello"} 的内存布局：5 个字符 + 1 个结束符}
  \label{fig:string-memory}
\end{figure}

\begin{pitfall}[没有 \texttt{\textbackslash0} 就不是字符串]
\begin{codebox}{no-null.c —— 缺少结束符}
char s[5] = {'h', 'e', 'l', 'l', 'o'};   /* 没有 '\0' */
printf("%s\n", s);                       /* 会一直往后读，直到碰巧遇到 0 */
\end{codebox}
打印时 \code{printf} 不知道什么时候停，会读到数组后面的内存，输出乱码甚至崩溃。\textbf{字符串必须以 \texttt{\textbackslash0} 结尾}。
\end{pitfall}

\subsection{字符串常量}

双引号括起来的是字符串常量，例如 \code{"hello"}。编译器会在末尾自动加 \texttt{\textbackslash0}：

\begin{codebox}{string-literal.c —— 字符串常量}
#include <stdio.h>

int main(void)
{
    printf("%s\n", "hello");        /* 直接使用字符串常量 */
    printf("sizeof(\"hello\") = %d\n", (int)sizeof("hello"));

    return 0;
}
\end{codebox}

\begin{outbox}{运行结果}
hello
sizeof("hello") = 6
\end{outbox}

注意 \code{sizeof("hello")} 是 6，不是 5。

\section{字符数组与字符指针}

这是字符串最容易混淆的地方。两种写法看起来都能用，但含义不同：

\begin{codebox}{array-vs-pointer.c —— 两种“字符串”}
char s1[] = "hello";    /* 字符数组：在栈上复制一份，可修改 */
char *s2 = "hello";     /* 字符指针：指向只读的字符串常量 */
\end{codebox}

\begin{figure}[htbp]
  \centering
  \begin{tikzpicture}[x=1cm,y=1cm]
    % 数组
    \node[font=\small\bfseries] at (0,2.6) {\texttt{char s1[] = "hello"}};
    \foreach \i/\c in {0/h,1/e,2/l,3/l,4/o,5/{$\backslash$0}} {
      \node[membox,minimum width=1.1cm] at (\i*1.3,1.8) {\ttfamily\c};
    }
    \node[font=\scriptsize,color=black!55] at (2.6,0.9) {数组：6 个字节都在栈上，可以改};

    % 指针
    \node[font=\small\bfseries] at (0,0.2) {\texttt{char *s2 = "hello"}};
    \node[ptrbox,minimum width=1.4cm] (p) at (0,-0.6) {\texttt{s2}};
    \node[draw=black!30,fill=black!7,rounded corners=1mm,minimum width=6.5cm,minimum height=0.8cm]
      (lit) at (5.2,-0.6) {\ttfamily h e l l o \textbackslash0};
    \draw[ptr] (p.east) -- (lit.west);
    \node[font=\scriptsize,color=black!55] at (4.5,-1.4) {字符串常量在只读区，不能改};
  \end{tikzpicture}
  \caption{字符数组会在栈上复制一份；字符指针只指向只读的字符串常量}
  \label{fig:string-array-pointer}
\end{figure}

\begin{center}
\begin{tabular}{@{}lll@{}}
\toprule
比较 & \code{char s[] = "hello";} & \code{char *p = "hello";} \\
\midrule
存储 & 数组，6 字节在栈上 & 指针，指向只读常量区 \\
能否修改 & 可以 \code{s[0] = 'H';} & \textbf{不可以}，未定义行为 \\
\code{sizeof} & 6（整个数组） & 8（指针大小，64 位） \\
重新赋值 & 不可以 \code{s = ...} & 可以 \code{p = ...} \\
\bottomrule
\end{tabular}
\end{center}

\begin{codebox}{modify-string.c —— 修改字符串}
#include <stdio.h>

int main(void)
{
    char s[] = "hello";

    s[0] = 'H';
    printf("%s\n", s);      /* Hello */

    return 0;
}
\end{codebox}

\begin{pitfall}[不要修改字符串常量]
\begin{codebox}{modify-literal.c —— 错误：修改只读常量}
char *p = "hello";
p[0] = 'H';     /* 未定义行为，很多平台上直接崩溃 */
\end{codebox}
如果需要修改，用字符数组；如果只读，用 \code{const char *p = "hello";} 让编译器帮你检查。
\end{pitfall}

\section{字符串的输入与输出}

\subsection{输出：\texttt{printf} 与 \texttt{puts}}

\begin{codebox}{string-output.c —— 输出字符串}
#include <stdio.h>

int main(void)
{
    char s[] = "hello";

    printf("%s\n", s);      /* printf 输出，不自动换行 */
    puts(s);                /* puts 输出并自动换行 */
    return 0;
}
\end{codebox}

\begin{outbox}{运行结果}
hello
hello
\end{outbox}

\subsection{输入：\texttt{scanf("\%s")} 的局限}

\begin{codebox}{string-input.c —— 读入一个单词}
#include <stdio.h>

int main(void)
{
    char name[20];

    printf("请输入姓名：");
    scanf("%19s", name);    /* 19 是最大宽度，防止溢出 */

    printf("你好，%s\n", name);
    return 0;
}
\end{codebox}

\texttt{scanf("\%s", name)} 会跳过前导空白，然后读到下一个空白（空格、制表符、换行）为止。所以它只能读“一个单词”，不能读带空格的句子。

\begin{pitfall}[\texttt{scanf("\%s")} 的缓冲区溢出]
\begin{codebox}{scanf-overflow.c —— 危险写法}
char name[20];
scanf("%s", name);      /* 如果输入超过 19 个字符，就会越界写 */
\end{codebox}
务必写成 \texttt{scanf("\%19s", name)}，或者改用 \code{fgets}。
\end{pitfall}

\subsection{读入一整行：\texttt{fgets}}

\begin{codebox}{fgets.c —— 读入一行}
#include <stdio.h>
#include <string.h>

int main(void)
{
    char line[100];
    int i;

    printf("请输入一行文字：");
    fgets(line, sizeof(line), stdin);

    /* fgets 会把换行符也读进来，去掉它 */
    for (i = 0; line[i] != '\0'; i++) {
        if (line[i] == '\n') {
            line[i] = '\0';
            break;
        }
    }

    printf("你输入的是：%s\n", line);
    return 0;
}
\end{codebox}

\begin{outbox}{运行示例}
请输入一行文字：hello world
你输入的是：hello world
\end{outbox}

\code{fgets} 最多读 \code{size - 1} 个字符，并自动补 \texttt{\textbackslash0}，比 \code{scanf} 安全。它的缺点是会把换行符也存进来，需要手动去掉。

\begin{tipbox}[输入方式怎么选]
\begin{itemize}
  \item 读一个单词（无空格）：\texttt{scanf("\%19s", s)};
  \item 读一整行：\code{fgets(s, sizeof(s), stdin)}，再去掉换行；
  \item 读单个字符：\code{getchar()}。
\end{itemize}
\end{tipbox}

\section{常用字符串函数}

C 标准库 \file{string.h} 提供了一组字符串函数。使用前要 \code{#include <string.h>}。

\begin{center}
\begin{tabular}{@{}lll@{}}
\toprule
函数 & 作用 & 注意 \\
\midrule
\code{strlen(s)} & 求长度（不含 \texttt{\textbackslash0}） & 返回值是 \code{size_t} \\
\code{strcpy(dst, src)} & 复制字符串 & 目标空间必须够大 \\
\code{strncpy(dst, src, n)} & 最多复制 n 个字符 & 可能不补 \texttt{\textbackslash0} \\
\code{strcat(dst, src)} & 连接字符串 & 目标空间必须够大 \\
\code{strncat(dst, src, n)} & 最多连接 n 个字符 & 会自动补 \texttt{\textbackslash0} \\
\code{strcmp(a, b)} & 比较字符串 & 返回 0 表示相等 \\
\code{strncmp(a, b, n)} & 比较前 n 个字符 & \\
\code{strchr(s, c)} & 查找字符 & 返回指针或 \code{NULL} \\
\code{strstr(s, sub)} & 查找子串 & 返回指针或 \code{NULL} \\
\bottomrule
\end{tabular}
\end{center}

\subsection{\texttt{strlen}：求长度}

\begin{codebox}{strlen.c —— strlen 与 sizeof}
#include <stdio.h>
#include <string.h>

int main(void)
{
    char s[] = "hello";

    printf("strlen(s) = %d\n", (int)strlen(s));
    printf("sizeof(s) = %d\n", (int)sizeof(s));
    return 0;
}
\end{codebox}

\begin{outbox}{运行结果}
strlen(s) = 5
sizeof(s) = 6
\end{outbox}

\code{strlen} 数到 \texttt{\textbackslash0} 为止，不包含它；\code{sizeof} 是数组占的总字节数，包含 \texttt{\textbackslash0}。

\subsection{\texttt{strcpy} 与 \texttt{strcat}}

\begin{codebox}{strcpy.c —— 复制与连接}
#include <stdio.h>
#include <string.h>

int main(void)
{
    char dst[50];

    strcpy(dst, "hello");
    strcat(dst, ", world");
    printf("%s\n", dst);

    return 0;
}
\end{codebox}

\begin{outbox}{运行结果}
hello, world
\end{outbox}

\begin{pitfall}[\texttt{strcpy} 不检查目标空间]
\begin{codebox}{strcpy-overflow.c —— 溢出}
char dst[5];
strcpy(dst, "hello world");   /* 目标只有 5 字节，严重越界 */
\end{codebox}
使用前必须确认 \code{dst} 足够大：至少 \code{strlen(src) + 1} 字节。
\end{pitfall}

\subsection{\texttt{strcmp}：比较字符串}

\begin{codebox}{strcmp.c —— 比较字符串}
#include <stdio.h>
#include <string.h>

int main(void)
{
    char a[] = "apple";
    char b[] = "banana";

    if (strcmp(a, b) == 0) {
        printf("相等\n");
    } else if (strcmp(a, b) < 0) {
        printf("a 在 b 前面\n");
    } else {
        printf("a 在 b 后面\n");
    }

    return 0;
}
\end{codebox}

\begin{outbox}{运行结果}
a 在 b 前面
\end{outbox}

\begin{pitfall}[字符串不能用 \texttt{==} 比较]
\code{if (a == b)} 比较的是两个数组的地址，不是内容。比较内容必须用 \code{strcmp(a, b) == 0}。
\end{pitfall}

\subsection{查找：\texttt{strchr} 与 \texttt{strstr}}

\begin{codebox}{strchr.c —— 查找字符与子串}
#include <stdio.h>
#include <string.h>

int main(void)
{
    char s[] = "hello world";
    char *p;

    p = strchr(s, 'w');
    if (p != NULL) {
        printf("找到 w，位置：%d\n", (int)(p - s));
        printf("从 w 开始：%s\n", p);
    }

    p = strstr(s, "world");
    if (p != NULL) {
        printf("找到子串，位置：%d\n", (int)(p - s));
    }

    return 0;
}
\end{codebox}

\begin{outbox}{运行结果}
找到 w，位置：6
从 w 开始：world
找到子串，位置：6
\end{outbox}

\subsection{自己实现 \texttt{strlen} 和 \texttt{strcpy}}

理解原理的最好方式是亲手写一遍：

\begin{codebox}{my-string.c —— 自己实现字符串函数}
#include <stdio.h>

int my_strlen(const char *s)
{
    int len = 0;

    while (s[len] != '\0') {
        len++;
    }
    return len;
}

void my_strcpy(char *dst, const char *src)
{
    int i = 0;

    while (src[i] != '\0') {
        dst[i] = src[i];
        i++;
    }
    dst[i] = '\0';      /* 别忘了结束符 */
}

int main(void)
{
    char buf[20];

    my_strcpy(buf, "hello");
    printf("%s, len = %d\n", buf, my_strlen(buf));
    return 0;
}
\end{codebox}

用指针写会更简洁：

\begin{codebox}{my-string-pointer.c —— 指针版}
int my_strlen(const char *s)
{
    const char *p = s;

    while (*p != '\0') {
        p++;
    }
    return (int)(p - s);
}

void my_strcpy(char *dst, const char *src)
{
    while ((*dst = *src) != '\0') {
        dst++;
        src++;
    }
}
\end{codebox}

\section{字符串的处理}

\subsection{统计字符}

\begin{codebox}{count-chars.c —— 统计字母、数字、其他字符}
#include <stdio.h>

int main(void)
{
    char s[100];
    int i;
    int letters = 0;
    int digits = 0;
    int others = 0;

    printf("请输入一行文字：");
    fgets(s, sizeof(s), stdin);

    for (i = 0; s[i] != '\0'; i++) {
        if ((s[i] >= 'a' && s[i] <= 'z') || (s[i] >= 'A' && s[i] <= 'Z')) {
            letters++;
        } else if (s[i] >= '0' && s[i] <= '9') {
            digits++;
        } else if (s[i] != '\n') {
            others++;
        }
    }

    printf("字母 %d 个，数字 %d 个，其他 %d 个\n",
           letters, digits, others);
    return 0;
}
\end{codebox}

\begin{outbox}{运行示例}
请输入一行文字：abc123 !!
字母 3 个，数字 3 个，其他 3 个
\end{outbox}

\subsection{大小写转换}

\begin{codebox}{case-convert.c —— 大小写转换}
#include <stdio.h>

int main(void)
{
    char s[100];
    int i;

    printf("请输入一行文字：");
    fgets(s, sizeof(s), stdin);

    for (i = 0; s[i] != '\0'; i++) {
        if (s[i] >= 'a' && s[i] <= 'z') {
            s[i] = s[i] - 32;       /* 小写转大写 */
        } else if (s[i] >= 'A' && s[i] <= 'Z') {
            s[i] = s[i] + 32;       /* 大写转小写 */
        }
    }

    printf("转换后：%s\n", s);
    return 0;
}
\end{codebox}

也可以用 \file{ctype.h} 里的 \code{toupper}、\code{tolower}、\code{isdigit}、\code{isalpha}，它们更可移植。

\subsection{判断回文}

\begin{codebox}{palindrome.c —— 判断回文}
#include <stdio.h>
#include <string.h>

int main(void)
{
    char s[100];
    int i;
    int j;
    int is_pal = 1;

    printf("请输入一个字符串：");
    scanf("%99s", s);

    i = 0;
    j = (int)strlen(s) - 1;

    while (i < j) {
        if (s[i] != s[j]) {
            is_pal = 0;
            break;
        }
        i++;
        j--;
    }

    if (is_pal) {
        printf("是回文\n");
    } else {
        printf("不是回文\n");
    }

    return 0;
}
\end{codebox}

\begin{outbox}{运行示例}
请输入一个字符串：level
是回文
\end{outbox}

\subsection{统计单词个数}

\begin{codebox}{word-count.c —— 统计单词个数}
#include <stdio.h>

int main(void)
{
    char s[200];
    int i;
    int count = 0;
    int in_word = 0;

    printf("请输入一行文字：");
    fgets(s, sizeof(s), stdin);

    for (i = 0; s[i] != '\0'; i++) {
        if (s[i] == ' ' || s[i] == '\n' || s[i] == '\t') {
            in_word = 0;
        } else {
            if (!in_word) {
                count++;
                in_word = 1;
            }
        }
    }

    printf("单词个数：%d\n", count);
    return 0;
}
\end{codebox}

思路：用一个标志 \code{in_word} 记录“当前是否在单词里”。从空白进入非空白时，单词数加 1。

\section{字符串与指针}

\subsection{用指针遍历字符串}

\begin{codebox}{string-pointer.c —— 指针遍历}
#include <stdio.h>

int main(void)
{
    char s[] = "hello";
    char *p;

    for (p = s; *p != '\0'; p++) {
        printf("%c ", *p);
    }
    printf("\n");

    return 0;
}
\end{codebox}

\begin{outbox}{运行结果}
h e l l o
\end{outbox}

\code{p++} 每次指向下一个字符，\texttt{*p != '\textbackslash0'} 作为结束条件。

\subsection{字符串数组}

如果要存多个字符串，有两种常见方式：

\begin{codebox}{string-list.c —— 多个字符串}
/* 方式一：二维字符数组，每行一个字符串，可修改 */
char names[3][20] = {"Alice", "Bob", "Carol"};

/* 方式二：字符指针数组，每个指针指向一个字符串常量，只读 */
char *colors[3] = {"red", "green", "blue"};
\end{codebox}

方式一每个字符串固定占 20 字节，浪费空间但可修改；方式二只占指针和常量本身，但不能修改内容。根据需求选择。

\section{中文与字符编码}

一个英文字符在 ASCII 里占 1 字节，但中文不是。以 UTF-8 为例，一个汉字通常占 3 字节：

\begin{codebox}{utf8.c —— 中文的长度}
#include <stdio.h>
#include <string.h>

int main(void)
{
    char s[] = "你好";

    printf("strlen(s) = %d\n", (int)strlen(s));
    printf("sizeof(s) = %d\n", (int)sizeof(s));
    return 0;
}
\end{codebox}

\begin{outbox}{运行结果（UTF-8 环境）}
strlen(s) = 6
sizeof(s) = 7
\end{outbox}

“你好”只有 2 个汉字，\code{strlen} 却是 6——因为 UTF-8 下每个汉字占 3 字节。\code{strlen} 数的是\textbf{字节数}，不是字符数。

\begin{itemize}
  \item \textbf{GBK}：一个汉字通常 2 字节；
  \item \textbf{UTF-8}：一个汉字通常 3 字节；
  \item 文件编码、终端编码、程序里的编码不一致，就会出现乱码。
\end{itemize}

\begin{tipbox}[处理中文的原则]
\begin{itemize}
  \item 输入输出用 \code{char} 数组存 UTF-8 字节，一般能原样显示；
  \item 不要用 \code{strlen} 判断“几个汉字”；
  \item 截断字符串时不能从中间截断一个 UTF-8 字符，否则会乱码；
  \item 真正需要按“字符”处理时，用宽字符 \code{wchar_t}（平台相关，初学了解即可）。
\end{itemize}
\end{tipbox}

\section{常见错误}

\begin{center}
\begin{tabular}{@{}p{0.38\textwidth}p{0.52\textwidth}@{}}
\toprule
错误 & 后果 \\
\midrule
字符数组没有 \texttt{\textbackslash0} & 打印越界、乱码 \\
\code{char *p = "abc"; p[0] = 'A';} & 修改只读常量，崩溃 \\
\code{if (s1 == s2)} & 比较地址，不是内容 \\
\code{s = "abc";}（数组） & 数组不能整体赋值，用 \code{strcpy} \\
\texttt{scanf("\%s", s)} 不限宽 & 缓冲区溢出 \\
\code{strcpy} 目标太小 & 溢出 \\
\code{strncpy} 后忘记补 \texttt{\textbackslash0} & 字符串没有结束符 \\
把 \code{strlen} 当字符数（中文） & 结果比预期大 \\
\bottomrule
\end{tabular}
\end{center}

\begin{pitfall}[\texttt{strncpy} 不保证补零]
\begin{codebox}{strncpy.c —— strncpy 的坑}
char dst[5];
strncpy(dst, "hello", 5);   /* dst 里没有 '\0'！ */
dst[4] = '\0';              /* 需要手动补 */
\end{codebox}
\code{strncpy} 只在源字符串长度小于 \code{n} 时补 \texttt{\textbackslash0}；如果源字符串更长，它不会补。
\end{pitfall}

\section{综合例题}

\begin{example}
输入一个字符串，删除其中所有的空格，输出结果。
\end{example}
\begin{solution}
用两个下标：\code{i} 读，\code{j} 写。遇到空格就跳过，否则写到 \code{j} 位置。

\begin{codebox}{remove-space.c —— 删除空格}
#include <stdio.h>

int main(void)
{
    char s[200];
    int i;
    int j = 0;

    printf("请输入一行文字：");
    fgets(s, sizeof(s), stdin);

    for (i = 0; s[i] != '\0'; i++) {
        if (s[i] != ' ') {
            s[j] = s[i];
            j++;
        }
    }
    s[j] = '\0';

    printf("删除空格后：%s\n", s);
    return 0;
}
\end{codebox}
\end{solution}

\begin{example}
输入一个字符串，统计每个字母出现的次数（只统计小写字母）。
\end{example}
\begin{solution}
思路如下：

\begin{codebox}{letter-count.c —— 字母频次}
#include <stdio.h>

int main(void)
{
    char s[200];
    int count[26] = {0};
    int i;

    printf("请输入一行文字：");
    fgets(s, sizeof(s), stdin);

    for (i = 0; s[i] != '\0'; i++) {
        if (s[i] >= 'a' && s[i] <= 'z') {
            count[s[i] - 'a']++;
        }
    }

    for (i = 0; i < 26; i++) {
        if (count[i] > 0) {
            printf("%c: %d\n", 'a' + i, count[i]);
        }
    }

    return 0;
}
\end{codebox}
\end{solution}

\begin{example}
输入两个字符串，判断第二个是不是第一个的子串（不用 \texttt{strstr}）。
\end{example}
\begin{solution}
思路如下：

\begin{codebox}{substring.c —— 判断子串}
#include <stdio.h>

int main(void)
{
    char s[200];
    char sub[100];
    int i, j, k;

    printf("请输入主串：");
    scanf("%199s", s);
    printf("请输入子串：");
    scanf("%99s", sub);

    for (i = 0; s[i] != '\0'; i++) {
        j = 0;
        k = i;
        while (sub[j] != '\0' && s[k] == sub[j]) {
            j++;
            k++;
        }
        if (sub[j] == '\0') {
            printf("是子串，起始位置 %d\n", i);
            return 0;
        }
    }

    printf("不是子串\n");
    return 0;
}
\end{codebox}
\end{solution}

\section{本章小结}

\begin{center}
\begin{tabular}{@{}p{0.24\textwidth}p{0.66\textwidth}@{}}
\toprule
主题 & 要点 \\
\midrule
字符串 & 以 \texttt{\textbackslash0} 结尾的字符数组 \\
字符数组 & 可修改；\code{sizeof} 是数组总大小 \\
字符指针 & 指向字符串常量；不能修改内容 \\
输入 & \texttt{scanf("\%19s")} 防溢出；\code{fgets} 读整行 \\
输出 & \texttt{printf("\%s")} 不换行；\code{puts} 自动换行 \\
常用函数 & \code{strlen}、\code{strcpy}、\code{strcat}、\code{strcmp}、\code{strchr}、\code{strstr} \\
比较 & 用 \code{strcmp}，不能用 \code{==} \\
中文 & \code{strlen} 数字节，不数字符 \\
\bottomrule
\end{tabular}
\end{center}

\begin{tipbox}[下一章预告]
有了数组和字符串，我们就可以处理一批数据了。但数据往往是乱序的，查找、统计、输出都需要先排序。下一章学习几种经典排序算法，并比较它们的效率。
\end{tipbox}

\begin{problemset}
  \item 输入一个字符串，输出它的长度（分别用 \code{strlen} 和自己写的函数）。
  \item 输入两个字符串，把第二个连接到第一个后面（用 \code{strcat} 和自己写的函数）。
  \item 输入两个字符串，比较它们的大小（用 \code{strcmp}）。
  \item 输入一个字符串，统计其中字母、数字、空格和其他字符的个数。
  \item 输入一个字符串，把所有小写字母转成大写。
  \item 判断一个字符串是否为回文。
  \item 输入一行文字，统计单词个数。
  \item 输入一个字符串，删除所有空格。
  \item 输入一个字符串，输出其中每个字母出现的次数。
  \item 思考题：为什么 \code{char s[] = "hello";} 可以修改，而 \code{char *p = "hello";} 不能修改？它们的内存分别在哪里？
\end{problemset}

\chapter{排序算法}

\begin{introduction}
  \item 为什么要排序
  \item 冒泡排序：两两比较，大的往后走
  \item 选择排序：每次挑最小的
  \item 插入排序：像整理扑克牌
  \item 三种排序怎么选
\end{introduction}

\chaptermotto{「把乱七八糟的数据排好队，后面的事情就简单了。」}

考试成绩要排名，找最大值要先比较，查资料要按顺序找。生活里到处都需要“排好队”。

这一章只学三种最基础、最好懂的排序：冒泡排序、选择排序、插入排序。我们不讲复杂的公式，只用手把过程走一遍，再把过程翻译成 C 语言。

\section{为什么要排序}

先看一组没有排序的分数：

\begin{center}
\begin{tabular}{@{}cccccc@{}}
\toprule
90 & 65 & 78 & 100 & 82 & 55 \\
\bottomrule
\end{tabular}
\end{center}

如果老师要问“最高分是谁”，你得从头看一遍。如果数据有 1000 个，就更麻烦了。

排好序以后：

\begin{center}
\begin{tabular}{@{}cccccc@{}}
\toprule
55 & 65 & 78 & 82 & 90 & 100 \\
\bottomrule
\end{tabular}
\end{center}

最高分在最后，最低分在最前面，找起来一眼就能看到。后面学查找、统计、排名，也会方便很多。

\textbf{排序}就是把一组数据按照从小到大（或从大到小）的顺序重新排列。

\section{冒泡排序}

\subsection{先用手做一遍}

冒泡排序的规则很简单：\textbf{从左到右，两两比较，前面的比后面大就交换。}

先拿 5 个数试一试：

\begin{center}
\begin{tabular}{@{}ccccc@{}}
\toprule
5 & 2 & 9 & 1 & 7 \\
\bottomrule
\end{tabular}
\end{center}

第一轮：

\begin{itemize}
  \item 比较 5 和 2：5 比 2 大，交换，变成 \texttt{2 5 9 1 7}；
  \item 比较 5 和 9：5 比 9 小，不换，还是 \texttt{2 5 9 1 7}；
  \item 比较 9 和 1：9 比 1 大，交换，变成 \texttt{2 5 1 9 7}；
  \item 比较 9 和 7：9 比 7 大，交换，变成 \texttt{2 5 1 7 9}。
\end{itemize}

第一轮结束，最大的 9 已经“冒”到了最后面。就像水里的气泡，大的会往上冒。

第二轮再从左边开始，但最后一个 9 已经排好，不用再比。这样一轮一轮下去，直到所有数都排好。

\subsection{写成 C 语言}

\begin{codebox}{bubble.c —— 冒泡排序}
#include <stdio.h>

int main(void)
{
    int a[8] = {5, 2, 9, 1, 7, 3, 8, 6};
    int n = 8;
    int i, j, temp;

    /* 排序前 */
    for (i = 0; i < n; i++) {
        printf("%d ", a[i]);
    }
    printf("\n");

    /* 冒泡排序 */
    for (i = 0; i < n - 1; i++) {
        for (j = 0; j < n - 1 - i; j++) {
            if (a[j] > a[j + 1]) {
                temp = a[j];
                a[j] = a[j + 1];
                a[j + 1] = temp;
            }
        }
    }

    /* 排序后 */
    for (i = 0; i < n; i++) {
        printf("%d ", a[i]);
    }
    printf("\n");

    return 0;
}
\end{codebox}

\begin{outbox}{运行结果}
5 2 9 1 7 3 8 6
1 2 3 5 6 7 8 9
\end{outbox}

\subsection{这段代码在做什么}

\begin{itemize}
  \item 外层循环 \texttt{i} 表示“第几轮”；
  \item 内层循环 \texttt{j} 表示“这一轮比较到第几个”；
  \item \texttt{n - 1 - i} 的意思是：每轮结束后，后面已经有 \texttt{i} 个数排好了，不用再比较；
  \item 如果前面的数比后面的大，就借助 \texttt{temp} 交换它们。
\end{itemize}

\begin{tipbox}[交换两个数]
交换要用临时变量，就像交换两杯水需要第三个空杯子：
\begin{codebox}{swap-three.c}
temp = a[j];
a[j] = a[j + 1];
a[j + 1] = temp;
\end{codebox}
\end{tipbox}

\begin{note}
冒泡排序每轮都把当前最大的数送到后面，所以叫“冒泡”。数据少的时候很好用，数据多了会比较慢。初学阶段先把它写熟。
\end{note}

\section{选择排序}

\subsection{先用手做一遍}

选择排序的规则是：\textbf{每一轮找出最小的数，放到最前面。}

还是这 5 个数：

\begin{center}
\begin{tabular}{@{}ccccc@{}}
\toprule
5 & 2 & 9 & 1 & 7 \\
\bottomrule
\end{tabular}
\end{center}

\begin{enumerate}
  \item 从全部数里找最小的：是 1；
  \item 把 1 和第一个数 5 交换：\texttt{1 2 9 5 7}；
  \item 从剩下的 \texttt{2 9 5 7} 里找最小的：是 2，已经在第二位，不用换；
  \item 从剩下的 \texttt{9 5 7} 里找最小的：是 5，和 9 交换：\texttt{1 2 5 9 7}；
  \item 从剩下的 \texttt{9 7} 里找最小的：是 7，和 9 交换：\texttt{1 2 5 7 9}。
\end{enumerate}

排好了。

\subsection{写成 C 语言}

\begin{codebox}{selection.c —— 选择排序}
#include <stdio.h>

int main(void)
{
    int a[8] = {5, 2, 9, 1, 7, 3, 8, 6};
    int n = 8;
    int i, j, min_index, temp;

    for (i = 0; i < n; i++) {
        printf("%d ", a[i]);
    }
    printf("\n");

    for (i = 0; i < n - 1; i++) {
        min_index = i;
        for (j = i + 1; j < n; j++) {
            if (a[j] < a[min_index]) {
                min_index = j;
            }
        }
        if (min_index != i) {
            temp = a[i];
            a[i] = a[min_index];
            a[min_index] = temp;
        }
    }

    for (i = 0; i < n; i++) {
        printf("%d ", a[i]);
    }
    printf("\n");

    return 0;
}
\end{codebox}

\begin{outbox}{运行结果}
5 2 9 1 7 3 8 6
1 2 3 5 6 7 8 9
\end{outbox}

\subsection{这段代码在做什么}

\begin{itemize}
  \item \texttt{i} 表示“现在要填第几个位置”；
  \item 先假设第 \texttt{i} 个位置上的数最小，记下它的下标 \texttt{min\_index}；
  \item 用 \texttt{j} 从后面找一圈，发现更小的就更新 \texttt{min\_index}；
  \item 找完以后，如果最小数不在第 \texttt{i} 个位置，就交换过去。
\end{itemize}

选择排序的好处是\textbf{交换次数少}，每一轮最多交换一次。

\section{插入排序}

\subsection{先用手做一遍}

插入排序就像整理扑克牌：

\begin{itemize}
  \item 左手拿已经排好的牌；
  \item 右手每次摸一张新牌；
  \item 从右往左比，把新牌插到合适的位置。
\end{itemize}

例如 \texttt{5 2 9 1 7}：

\begin{enumerate}
  \item 先看第一个数 5，它自己就是有序的；
  \item 摸到 2，比 5 小，插到 5 前面：\texttt{2 5}；
  \item 摸到 9，比 5 大，放在后面：\texttt{2 5 9}；
  \item 摸到 1，比 9、5、2 都小，插到最前面：\texttt{1 2 5 9}；
  \item 摸到 7，插到 5 和 9 之间：\texttt{1 2 5 7 9}。
\end{enumerate}

\subsection{写成 C 语言}

\begin{codebox}{insertion.c —— 插入排序}
#include <stdio.h>

int main(void)
{
    int a[8] = {5, 2, 9, 1, 7, 3, 8, 6};
    int n = 8;
    int i, j, key;

    for (i = 0; i < n; i++) {
        printf("%d ", a[i]);
    }
    printf("\n");

    for (i = 1; i < n; i++) {
        key = a[i];             /* 摸到的新牌 */
        j = i - 1;

        /* 从右往左找位置，比 key 大的都往后挪 */
        while (j >= 0 && a[j] > key) {
            a[j + 1] = a[j];
            j--;
        }
        a[j + 1] = key;         /* 把新牌放进去 */
    }

    for (i = 0; i < n; i++) {
        printf("%d ", a[i]);
    }
    printf("\n");

    return 0;
}
\end{codebox}

\begin{outbox}{运行结果}
5 2 9 1 7 3 8 6
1 2 3 5 6 7 8 9
\end{outbox}

\subsection{这段代码在做什么}

\begin{itemize}
  \item 从第二个数开始，把当前数当作“新牌” \texttt{key}；
  \item 从它前面一个数开始，从右往左看；
  \item 如果前面的数比 \texttt{key} 大，就把前面的数往后挪一格；
  \item 挪到合适的位置后，把 \texttt{key} 放进去。
\end{itemize}

插入排序在数据\textbf{原本就比较有序}的时候特别快，因为不需要挪很多次。

\section{三种排序比一比}

\begin{center}
\begin{tabular}{@{}p{0.16\textwidth}p{0.24\textwidth}p{0.24\textwidth}p{0.24\textwidth}@{}}
\toprule
排序方法 & 怎么做的 & 优点 & 适合的情况 \\
\midrule
冒泡排序 & 相邻两个比较，大的往后换 & 最好理解 & 数据很少，练手 \\
选择排序 & 每轮找最小的放到前面 & 交换次数少 & 交换比较麻烦的时候 \\
插入排序 & 像抓牌，插入已排好的部分 & 数据基本有序时很快 & 数据不多、接近有序 \\
\bottomrule
\end{tabular}
\end{center}

\begin{tipbox}[初学怎么选]
初学阶段，先把\textbf{冒泡排序}写熟。它最直观，出错也容易查。等以后数据变多，再学更快的排序方法。
\end{tipbox}

\begin{note}
这三种排序在数据很多时都比较慢。真正的大数据排序，计算机会用更聪明的算法。那是后面进阶的内容，现在不用着急。
\end{note}

\section{小应用：给成绩排序}

下面把 5 个成绩从高到低排序，用的是冒泡排序。注意比较符号从 \texttt{>} 变成了 \texttt{<}：

\begin{codebox}{sort-scores.c —— 成绩从高到低}
#include <stdio.h>

int main(void)
{
    int score[5] = {78, 92, 65, 88, 70};
    int n = 5;
    int i, j, temp;

    for (i = 0; i < n - 1; i++) {
        for (j = 0; j < n - 1 - i; j++) {
            if (score[j] < score[j + 1]) {
                temp = score[j];
                score[j] = score[j + 1];
                score[j + 1] = temp;
            }
        }
    }

    printf("从高到低：");
    for (i = 0; i < n; i++) {
        printf("%d ", score[i]);
    }
    printf("\n");

    return 0;
}
\end{codebox}

\begin{outbox}{运行结果}
从高到低：92 88 78 70 65
\end{outbox}

只要把比较的“大于”改成“小于”，就能从升序变成降序。

\section{本章小结}

\begin{center}
\begin{tabular}{@{}p{0.24\textwidth}p{0.64\textwidth}@{}}
\toprule
主题 & 要点 \\
\midrule
排序 & 把数据按从小到大或从大到小的顺序排好 \\
冒泡排序 & 相邻比较，大的往后换；每轮把最大值送到末尾 \\
选择排序 & 每轮找最小值，放到前面 \\
插入排序 & 把当前数插入左边已排好的部分 \\
交换 & 需要临时变量 \texttt{temp} \\
升序/降序 & 改一下比较符号即可 \\
\bottomrule
\end{tabular}
\end{center}

\begin{tipbox}[下一章预告]
最后我们学宏定义：给一段代码起个名字，让写程序更方便。学完这一章，这本书的正文就结束啦。
\end{tipbox}

\begin{problemset}
  \item 用冒泡排序把 \texttt{5 3 8 1 2} 从小到大排序，写出每一轮的结果。
  \item 用选择排序把 \texttt{9 4 7 2 6} 从小到大排序，写出每一轮的结果。
  \item 用插入排序把 \texttt{6 2 8 4 1} 从小到大排序，写出每一步。
  \item 把本章的冒泡排序程序改成从大到小排序。
  \item 输入 10 个整数，用冒泡排序从小到大输出。
  \item 输入 10 个整数，用选择排序从大到小输出。
  \item 输入 10 个学生的成绩，用任意一种排序从高到低输出。
  \item 想一想：如果数据本来就已经排好序了，哪种排序最省事？为什么？
  \item 思考题：冒泡排序为什么要比较 \texttt{n - 1 - i} 次，而不是每次都比较 \texttt{n - 1} 次？
\end{problemset}

\chapter{宏定义}

\begin{introduction}
  \item 预处理器在编译之前做了什么
  \item 对象式宏与函数式宏
  \item 括号、副作用与 \texttt{do-while(0)}
  \item \texttt{\#} 与 \texttt{\#\#}：字符串化与连接
  \item 条件编译与头文件保护
  \item 预定义宏、调试宏与常见陷阱
\end{introduction}

\chaptermotto{「宏是编译之前的文本替换：威力很大，也很容易伤到自己。」}

宏定义（macro）由\textbf{预处理器}处理，发生在真正的编译之前。它本质上就是“文本替换”：你写一个名字，预处理器把它替换成一段文本。

宏可以用来定义常量、写“函数式”代码、做条件编译、开关调试信息。它没有类型检查，也不遵守作用域规则，所以用得好是利器，用不好就是灾难。这一章会把宏的能力和边界都讲清楚。

\section{预处理器}

\subsection{编译之前的步骤}

第 1 章讲过，C 程序从源码到可执行文件要经过预处理、编译、汇编、链接。预处理器负责处理所有以 \code{#} 开头的行：

\begin{figure}[htbp]
  \centering
  \begin{tikzpicture}[x=1cm,y=1cm,>=Stealth,node distance=8mm]
    \node[process] (src) {源文件\\\texttt{main.c}};
    \node[process,right=of src] (pp) {预处理器};
    \node[process,right=of pp] (cc) {编译器};
    \node[process,right=of cc] (obj) {目标文件};
    \draw[flow] (src) -- (pp);
    \draw[flow] (pp) -- node[lbl,above]{\texttt{main.i}} (cc);
    \draw[flow] (cc) -- (obj);
    \node[font=\scriptsize,color=black!55,below=4mm of pp] {展开 \texttt{\#include}、\texttt{\#define}、条件编译};
  \end{tikzpicture}
  \caption{预处理器在编译之前工作，只做文本层面的处理}
  \label{fig:preprocess}
\end{figure}

预处理器的输出仍然是 C 代码，只是所有宏都被替换了、头文件被展开了、条件编译的分支被裁剪了。

\subsection{预处理指令}

\begin{center}
\begin{tabular}{@{}ll@{}}
\toprule
指令 & 作用 \\
\midrule
\texttt{\#include} & 包含头文件 \\
\texttt{\#define} & 定义宏 \\
\texttt{\#undef} & 取消宏定义 \\
\texttt{\#if} / \texttt{\#ifdef} / \texttt{\#ifndef} & 条件编译 \\
\texttt{\#else} / \texttt{\#elif} / \texttt{\#endif} & 条件编译分支 \\
\texttt{\#error} & 产生编译错误 \\
\texttt{\#pragma} & 编译器相关指令 \\
\bottomrule
\end{tabular}
\end{center}

\section{对象式宏}

\subsection{定义常量}

\begin{codebox}{define-const.c —— 定义常量}
#include <stdio.h>

#define PI 3.1415926
#define MAX_SIZE 100
#define VERSION "1.0"

int main(void)
{
    printf("PI = %f\n", PI);
    printf("MAX_SIZE = %d\n", MAX_SIZE);
    printf("VERSION = %s\n", VERSION);

    return 0;
}
\end{codebox}

\begin{outbox}{运行结果}
PI = 3.141593
MAX_SIZE = 100
VERSION = 1.0
\end{outbox}

预处理器会把代码里所有的 \code{PI} 替换成 \code{3.1415926}，就像你在编辑器里用了“全部替换”。

\begin{pitfall}[宏定义不要加分号]
\begin{codebox}{define-semicolon.c —— 错误：多了分号}
#define MAX_SIZE 100;      /* 错误 */
int a[MAX_SIZE];            /* 展开成 int a[100;]; */
\end{codebox}
宏不是语句，末尾不加分号。否则替换后会多出一个分号，引发莫名其妙的错误。
\end{pitfall}

\subsection{宏与常量的区别}

\begin{codebox}{macro-vs-const.c —— 宏、const、enum}
#define MAX 100                 /* 宏：无类型，预处理替换 */
const int max_value = 100;      /* const 变量：有类型，编译器检查 */
enum { MAX_ENUM = 100 };        /* 枚举常量：有类型，可作数组长度 */
\end{codebox}

\begin{center}
\begin{tabular}{@{}lll@{}}
\toprule
方式 & 优点 & 缺点 \\
\midrule
\texttt{\#define} & 无类型限制、可用于预处理条件 & 无类型检查、难调试、可能重复求值 \\
\code{const} & 有类型、可调试 & C89 中不能作数组长度 \\
\code{enum} & 有类型、可作数组长度 & 只能是整数 \\
\bottomrule
\end{tabular}
\end{center}

\begin{tipbox}[常量怎么选]
整数常量优先用 \code{enum} 或 \code{const}；只在需要预处理条件（\texttt{\#if}）或跨平台配置时用 \texttt{\#define}。
\end{tipbox}

\section{函数式宏}

\subsection{基本写法}

\begin{codebox}{square-macro.c —— 函数式宏}
#include <stdio.h>

#define SQUARE(x) ((x) * (x))

int main(void)
{
    int a = 5;

    printf("%d\n", SQUARE(a));
    printf("%d\n", SQUARE(a + 1));

    return 0;
}
\end{codebox}

\begin{outbox}{运行结果}
25
36
\end{outbox}

\code{SQUARE(a + 1)} 展开成 \code{((a + 1) * (a + 1))}，结果是 36。

\subsection{括号！括号！括号！}

\begin{codebox}{macro-paren.c —— 不加括号的灾难}
#define SQUARE_BAD(x) x * x
#define SQUARE_GOOD(x) ((x) * (x))

SQUARE_BAD(2 + 3)    /* 展开成 2 + 3 * 2 + 3 = 11，而不是 25 */
SQUARE_GOOD(2 + 3)   /* 展开成 ((2 + 3) * (2 + 3)) = 25 */
\end{codebox}

\begin{pitfall}[宏的参数和整体都要加括号]
\begin{codebox}{macro-paren2.c —— 另一个经典错误}
#define DOUBLE(x) (x) + (x)
int y = DOUBLE(3) * 2;   /* 展开成 (3) + (3) * 2 = 9，而不是 12 */
\end{codebox}
正确写法：\texttt{\#define DOUBLE(x) ((x) + (x))}。记住：\textbf{每个参数加括号，整个表达式再加一层括号}。
\end{pitfall}

\subsection{副作用：参数被求值多次}

\begin{codebox}{macro-side-effect.c —— 宏的副作用}
#define SQUARE(x) ((x) * (x))

int i = 5;
int y = SQUARE(i++);    /* 展开成 ((i++) * (i++))，i 被加了两次 */
\end{codebox}

这属于未定义行为：\code{i} 在同一个表达式中被修改了两次。宏的参数不是“传值”，而是\textbf{文本替换}，所以会重复求值。

\begin{pitfall}[带副作用的参数不要传给宏]
\code{i++}、\code{func()}、\code{*p++} 这类表达式传给函数式宏时，可能被求值多次。要么先用临时变量接住，要么改用函数。
\end{pitfall}

\subsection{用 \texttt{do-while(0)} 包装多语句宏}

\begin{codebox}{do-while-macro.c —— 多语句宏}
#define SWAP(a, b) do { \
    int temp = (a);     \
    (a) = (b);          \
    (b) = temp;         \
} while (0)
\end{codebox}

为什么要包成 \texttt{do \{ ... \} while (0)}？因为这样宏展开后仍然是一条语句，可以安全地放在 \code{if-else} 里：

\begin{codebox}{swap-macro-use.c —— 使用多语句宏}
if (x > y)
    SWAP(x, y);
else
    printf("不用交换\n");
\end{codebox}

如果直接写成 \texttt{\{ ... \}}，展开后 \code{if} 和 \code{else} 的配对就会出问题。

\subsection{常用的函数式宏}

\begin{codebox}{common-macros.c —— 常见宏}
#define MAX(a, b) ((a) > (b) ? (a) : (b))
#define MIN(a, b) ((a) < (b) ? (a) : (b))
#define ABS(x)    ((x) < 0 ? -(x) : (x))
#define ARRAY_SIZE(a) (sizeof(a) / sizeof((a)[0]))
\end{codebox}

\code{ARRAY_SIZE} 只能用于真正的数组，不能用于指针参数。

\section{\texttt{\#} 与 \texttt{\#\#}}

\subsection{字符串化 \texttt{\#}}

在宏参数前加 \texttt{\#}，会把参数变成字符串：

\begin{codebox}{stringize.c —— 字符串化}
#include <stdio.h>

#define STR(x) #x

int main(void)
{
    printf("%s\n", STR(hello));
    printf("%s\n", STR(1 + 2));
    printf("%s\n", STR(MAX_SIZE));

    return 0;
}
\end{codebox}

\begin{outbox}{运行结果}
hello
1 + 2
MAX_SIZE
\end{outbox}

注意：\code{STR(MAX_SIZE)} 得到的是字符串 \code{"MAX_SIZE"}，不会先把宏展开。要展开后再字符串化，需要两层宏（进阶技巧）。

\subsection{连接 \texttt{\#\#}}

\texttt{\#\#} 把两个记号拼成一个：

\begin{codebox}{token-paste.c —— 记号连接}
#include <stdio.h>

#define MAKE_VAR(n) var_##n

int main(void)
{
    int MAKE_VAR(1) = 10;       /* 展开成 int var_1 = 10; */
    int MAKE_VAR(2) = 20;

    printf("%d %d\n", var_1, var_2);
    return 0;
}
\end{codebox}

\begin{outbox}{运行结果}
10 20
\end{outbox}

\texttt{\#\#} 常用于自动生成变量名、函数名，或者 X-Macro 技巧。

\section{条件编译}

\subsection{\texttt{\#ifdef} / \texttt{\#ifndef}}

\begin{codebox}{ifdef.c —— 条件编译}
#include <stdio.h>

#define DEBUG 1

int main(void)
{
#ifdef DEBUG
    printf("调试模式：程序开始运行\n");
#endif

    printf("程序主体\n");

    return 0;
}
\end{codebox}

常见用途：

\begin{itemize}
  \item 用 \code{DEBUG} 宏控制调试输出；
  \item 用 \texttt{\#ifdef \_WIN32} 区分 Windows 和 Linux；
  \item 用 \texttt{\#ifndef} 做头文件保护。
\end{itemize}

\subsection{头文件保护}

头文件可能被多次包含，导致重复定义。标准做法是用条件编译保护：

\begin{codebox}{student.h —— 头文件保护}
#ifndef STUDENT_H
#define STUDENT_H

struct Student {
    int id;
    char name[20];
    double score;
};

#endif
\end{codebox}

第一次包含时，\code{STUDENT_H} 未定义，进入内容并定义它；第二次包含时，\code{STUDENT_H} 已定义，直接跳过。

\begin{tipbox}[头文件保护的名字]
宏名通常用文件名大写 + 下划线，例如 \code{STUDENT_H}、\code{MOTOR_CTRL_H}。确保唯一，避免冲突。
\end{tipbox}

\subsection{\texttt{\#if} 与 \texttt{defined}}

\begin{codebox}{if-defined.c —— 更灵活的判断}
#if defined(DEBUG) && DEBUG > 0
    printf("调试输出\n");
#endif

#if VERSION >= 2
    printf("新版本\n");
#else
    printf("旧版本\n");
#endif
\end{codebox}

\texttt{\#if} 后面必须是一个整数常量表达式，不能使用变量。

\section{预定义宏}

预处理器还提供了一些预定义宏，调试时很有用：

\begin{center}
\begin{tabular}{@{}ll@{}}
\toprule
宏 & 含义 \\
\midrule
\code{__FILE__} & 当前源文件名（字符串） \\
\code{__LINE__} & 当前行号（整数） \\
\code{__DATE__} & 编译日期（字符串） \\
\code{__TIME__} & 编译时间（字符串） \\
\code{__STDC__} & 是否遵循标准 C \\
\bottomrule
\end{tabular}
\end{center}

\begin{codebox}{predefined.c —— 预定义宏}
#include <stdio.h>

int main(void)
{
    printf("文件：%s\n", __FILE__);
    printf("行号：%d\n", __LINE__);
    printf("编译日期：%s\n", __DATE__);
    printf("编译时间：%s\n", __TIME__);

    return 0;
}
\end{codebox}

\begin{outbox}{运行结果（示例）}
文件：predefined.c
行号：7
编译日期：Oct 21 2026
编译时间：13:30:00
\end{outbox}

\section{调试宏}

把 \code{__FILE__}、\code{__LINE__} 和 \texttt{\#\#} 结合起来，可以做一个简单的调试输出宏：

\begin{codebox}{debug-macro.c —— 调试宏}
#include <stdio.h>

#ifdef DEBUG
#define DBG_PRINT(fmt, ...) \
    printf("[%s:%d] " fmt "\n", __FILE__, __LINE__, __VA_ARGS__)
#else
#define DBG_PRINT(fmt, ...) ((void)0)
#endif
\end{codebox}

\begin{note}
上面用了 \code{...} 和 \code{__VA_ARGS__}，这是 C99 的可变参数宏，C89 编译器不一定支持。C89 下可以用固定参数或直接写 \code{printf}。本书正文代码不使用可变参数宏，这里只作了解。
\end{note}

C89 下更简单的做法：

\begin{codebox}{debug-c89.c —— C89 调试宏}
#define DBG(msg) printf("[%s:%d] %s\n", __FILE__, __LINE__, msg)
\end{codebox}

\section{宏的常见陷阱}

\begin{center}
\begin{tabular}{@{}p{0.38\textwidth}p{0.52\textwidth}@{}}
\toprule
错误 & 后果 \\
\midrule
\texttt{\#define MAX 100;} & 多出分号 \\
\texttt{\#define SQUARE(x) x*x} & 优先级错误 \\
参数有副作用 & 被求值多次 \\
宏名与变量名冲突 & 意外替换 \\
宏里写注释 & 注释可能影响替换 \\
宏没有类型检查 & 错误类型也能通过编译 \\
递归宏 & 展开不终止或行为异常 \\
\bottomrule
\end{tabular}
\end{center}

\begin{pitfall}[宏名不要用短小的常见词]
\texttt{\#define max 100} 会把代码里所有叫 \code{max} 的变量、函数都替换掉。宏名用大写，且尽量具体，例如 \code{MAX_BUFFER_SIZE}。
\end{pitfall}

\section{拔高：X-Macro 与代码生成}

有时候需要“一份数据，多处使用”。例如定义一组错误码，既要枚举值，又要转成字符串。X-Macro 可以做到：

\begin{codebox}{x-macro.c —— X-Macro 示意}
#define ERROR_LIST \
    X(OK, 0)       \
    X(TIMEOUT, 1)  \
    X(OVERFLOW, 2)

#define X(name, code) ERR_##name = code,
enum ErrorCode {
    ERROR_LIST
};
#undef X

#define X(name, code) case ERR_##name: return #name;
const char *error_name(int code)
{
    switch (code) {
        ERROR_LIST
    }
    return "UNKNOWN";
}
#undef X
\end{codebox}

这种写法在大型工程中用于生成“枚举 + 字符串 + 查表”等重复代码。初学阶段知道有这种技巧即可。

\section{本章小结}

\begin{center}
\begin{tabular}{@{}p{0.24\textwidth}p{0.66\textwidth}@{}}
\toprule
主题 & 要点 \\
\midrule
预处理 & 编译前的文本替换 \\
对象式宏 & 定义常量；不加分号；无类型检查 \\
函数式宏 & 参数整体加括号；注意副作用 \\
多语句宏 & 用 \texttt{do \{ ... \} while (0)} 包装 \\
\texttt{\#} / \texttt{\#\#} & 字符串化 / 记号连接 \\
条件编译 & \texttt{\#ifdef}、\texttt{\#ifndef}、\texttt{\#if} \\
头文件保护 & \texttt{\#ifndef} + \texttt{\#define} + \texttt{\#endif} \\
预定义宏 & \code{__FILE__}、\code{__LINE__} 等 \\
原则 & 能不用宏就不用；用就大写、加括号、写清楚 \\
\bottomrule
\end{tabular}
\end{center}

\begin{tipbox}[写在最后]
到这里，本书的正文就结束了。你从“计算机怎么存数据”出发，走过了顺序、分支、循环、数组、结构体、指针、字符串、排序和宏。接下来最重要的不是继续学新语法，而是\textbf{写代码}：找一些小项目（计算器、学生管理、串口协议解析、简单 PID），把学过的知识串起来。遇到问题回来看对应章节，进步会很快。
\end{tipbox}

\begin{problemset}
  \item 用宏定义 \code{PI}、\code{E} 和 \code{GOLDEN}，输出它们的值。
  \item 写一个 \code{SQUARE} 宏，并测试 \code{SQUARE(2+3)}。
  \item 写一个 \code{MAX3(a,b,c)} 宏，求三个数的最大值。
  \item 写一个 \code{SWAP} 多语句宏，用 \code{do-while(0)} 包装。
  \item 用 \texttt{\#} 写一个 \code{PRINT_VAR(x)} 宏，输出变量名和值。
  \item 用 \texttt{\#\#} 生成 \code{var_1}、\code{var_2} 两个变量。
  \item 用条件编译实现：定义 \code{DEBUG} 时输出调试信息，否则不输出。
  \item 写一个头文件保护，防止重复包含。
  \item 解释 \texttt{\#define SQUARE(x) x*x} 的问题，并改正。
  \item 思考题：为什么宏参数 \code{i++} 会导致问题？应该怎么避免？
\end{problemset}


\appendix
\chapter{C 语言速查表}
\label{app:cheatsheet}

本附录把全书最常用的语法、类型、格式符和标准库函数集中在一起，方便写代码时查阅。

\section{关键字（C89）}

C89 共有 32 个关键字，不能用作变量名或函数名：

\begin{center}
\small
\begin{tabular}{@{}llll@{}}
\toprule
\code{auto} & \code{break} & \code{case} & \code{char} \\
\code{const} & \code{continue} & \code{default} & \code{do} \\
\code{double} & \code{else} & \code{enum} & \code{extern} \\
\code{float} & \code{for} & \code{goto} & \code{if} \\
\code{int} & \code{long} & \code{register} & \code{return} \\
\code{short} & \code{signed} & \code{sizeof} & \code{static} \\
\code{struct} & \code{switch} & \code{typedef} & \code{union} \\
\code{unsigned} & \code{void} & \code{volatile} & \code{while} \\
\bottomrule
\end{tabular}
\end{center}

C99 新增的 \code{inline}、\code{restrict}、\code{_Bool} 等不属于本书范围。

\section{基本类型}

\begin{center}
\begin{tabular}{@{}llll@{}}
\toprule
类型 & 典型大小 & 范围/精度 & 格式符 \\
\midrule
\code{char} & 1 & $-128 \sim 127$ & \texttt{\%c} 或 \texttt{\%d} \\
\code{unsigned char} & 1 & $0 \sim 255$ & \texttt{\%u} \\
\code{short} & 2 & $-32768 \sim 32767$ & \texttt{\%d} \\
\code{unsigned short} & 2 & $0 \sim 65535$ & \texttt{\%u} \\
\code{int} & 4 & 约 $\pm 21$ 亿 & \texttt{\%d} \\
\code{unsigned int} & 4 & $0 \sim 4294967295$ & \texttt{\%u} \\
\code{long} & 4/8 & 平台相关 & \texttt{\%ld} \\
\code{float} & 4 & 约 6\textasciitilde7 位有效数字 & \texttt{\%f} \\
\code{double} & 8 & 约 15\textasciitilde16 位有效数字 & \texttt{\%f} / \texttt{\%lf} \\
\code{long double} & 8/16 & 平台相关 & \texttt{\%Lf} \\
\bottomrule
\end{tabular}
\end{center}

\section{常用格式符}

\begin{center}
\begin{tabular}{@{}lll@{}}
\toprule
格式符 & 含义 & 例子 \\
\midrule
\texttt{\%d} & 有符号十进制整数 & \texttt{printf("\%d", -42);} \\
\texttt{\%u} & 无符号十进制整数 & \texttt{printf("\%u", 42u);} \\
\texttt{\%x} / \texttt{\%X} & 十六进制（小/大写） & \texttt{\%x} 输出 \code{ff} \\
\texttt{\%o} & 八进制 & \\
\texttt{\%c} & 字符 & \texttt{printf("\%c", 'A');} \\
\texttt{\%s} & 字符串 & \texttt{printf("\%s", "hi");} \\
\texttt{\%f} & 浮点，默认 6 位小数 & \texttt{printf("\%f", 3.14);} \\
\texttt{\%.2f} & 保留 2 位小数 & \code{3.14} \\
\texttt{\%8.2f} & 总宽 8、2 位小数 & 右对齐 \\
\texttt{\%-8d} & 宽度 8、左对齐 & \\
\texttt{\%05d} & 宽度 5、补 0 & \code{00042} \\
\texttt{\%p} & 指针地址 & \texttt{(void *)\&a} \\
\texttt{\%\%} & 一个百分号 & \\
\bottomrule
\end{tabular}
\end{center}

\section{转义字符}

\begin{center}
\begin{tabular}{@{}llll@{}}
\toprule
转义 & 含义 & 转义 & 含义 \\
\midrule
\texttt{\textbackslash n} & 换行 & \texttt{\textbackslash t} & 制表符 \\
\texttt{\textbackslash r} & 回车 & \texttt{\textbackslash b} & 退格 \\
\texttt{\textbackslash\textbackslash} & 反斜杠 & \texttt{\textbackslash "} & 双引号 \\
\texttt{\textbackslash '} & 单引号 & \texttt{\textbackslash 0} & 空字符 \\
\bottomrule
\end{tabular}
\end{center}

\section{运算符优先级（高到低）}

\begin{center}
\small
\begin{tabular}{@{}clll@{}}
\toprule
优先级 & 运算符 & 含义 & 结合性 \\
\midrule
1 & \code{()} \code{[]} \code{->} \code{.} & 调用、下标、成员 & 左 \\
2 & \code{!} \code{\~} \code{++} \code{--} \code{+} \code{-} \code{*} \texttt{\&} \code{(类型)} \code{sizeof} & 单目 & 右 \\
3 & \code{*} \code{/} \texttt{\%} & 乘除余 & 左 \\
4 & \code{+} \code{-} & 加减 & 左 \\
5 & \code{<<} \code{>>} & 移位 & 左 \\
6 & \code{<} \code{<=} \code{>} \code{>=} & 关系 & 左 \\
7 & \code{==} \code{!=} & 相等 & 左 \\
8 & \texttt{\&} & 按位与 & 左 \\
9 & \code{\^} & 按位异或 & 左 \\
10 & \code{|} & 按位或 & 左 \\
11 & \texttt{\&\&} & 逻辑与 & 左 \\
12 & \code{||} & 逻辑或 & 左 \\
13 & \code{?:} & 条件 & 右 \\
14 & 赋值类 & \code{= += -= ...} & 右 \\
15 & \code{,} & 逗号 & 左 \\
\bottomrule
\end{tabular}
\end{center}

\section{常用标准库函数}

\subsection{\file{stdio.h}}

\begin{center}
\begin{tabular}{@{}ll@{}}
\toprule
函数 & 用途 \\
\midrule
\code{printf(fmt, ...)} & 格式化输出 \\
\code{scanf(fmt, ...)} & 格式化输入（变量要加 \texttt{\&}） \\
\code{getchar()} & 读一个字符 \\
\code{putchar(c)} & 输出一个字符 \\
\code{gets(s)} & 读一行（危险，不要用） \\
\code{fgets(s, n, stdin)} & 安全地读一行 \\
\code{puts(s)} & 输出字符串并换行 \\
\code{fopen/fclose} & 打开/关闭文件 \\
\bottomrule
\end{tabular}
\end{center}

\subsection{\file{string.h}}

\begin{center}
\begin{tabular}{@{}ll@{}}
\toprule
函数 & 用途 \\
\midrule
\code{strlen(s)} & 字符串长度（不含 \texttt{\textbackslash0}） \\
\code{strcpy(dst, src)} & 复制字符串 \\
\code{strncpy(dst, src, n)} & 最多复制 n 个字符 \\
\code{strcat(dst, src)} & 连接字符串 \\
\code{strcmp(a, b)} & 比较字符串，相等返回 0 \\
\code{strchr(s, c)} & 查找字符 \\
\code{strstr(s, sub)} & 查找子串 \\
\code{memcpy(dst, src, n)} & 复制内存块 \\
\code{memset(p, c, n)} & 填充内存块 \\
\bottomrule
\end{tabular}
\end{center}

\subsection{\file{stdlib.h}}

\begin{center}
\begin{tabular}{@{}ll@{}}
\toprule
函数 & 用途 \\
\midrule
\code{malloc(size)} & 申请内存 \\
\code{calloc(n, size)} & 申请并清零 \\
\code{realloc(p, size)} & 调整内存大小 \\
\code{free(p)} & 释放内存 \\
\code{abs(n)} & 整数绝对值 \\
\code{atoi(s)} & 字符串转整数 \\
\code{qsort(...)} & 快速排序 \\
\code{rand()/srand()} & 随机数 \\
\bottomrule
\end{tabular}
\end{center}

\subsection{\file{math.h} 与 \file{ctype.h}}

\begin{center}
\begin{tabular}{@{}ll@{}}
\toprule
函数 & 用途 \\
\midrule
\code{sqrt(x)} & 平方根 \\
\code{fabs(x)} & 浮点绝对值 \\
\code{pow(x, y)} & 幂 \\
\code{sin/cos/tan} & 三角函数（弧度） \\
\code{isalpha(c)} & 是否字母 \\
\code{isdigit(c)} & 是否数字 \\
\code{toupper(c)/tolower(c)} & 大小写转换 \\
\bottomrule
\end{tabular}
\end{center}

\section{常用代码片段}

\begin{codebox}{snippets.c —— 常用片段}
/* 交换两个整数 */
temp = a; a = b; b = temp;

/* 求数组最大值 */
max = a[0];
for (i = 1; i < n; i++) {
    if (a[i] > max) max = a[i];
}

/* 遍历数组 */
for (i = 0; i < n; i++) {
    printf("%d ", a[i]);
}

/* 字符串长度（自己实现） */
int len = 0;
while (s[len] != '\0') len++;

/* 判断闰年 */
if ((year % 4 == 0 && year % 100 != 0) || year % 400 == 0) {
    /* 闰年 */
}

/* 动态数组 */
p = (int *)malloc(n * sizeof(int));
if (p == NULL) { /* 处理失败 */ }
free(p);
p = NULL;
\end{codebox}

\chapter{编译、调试与常见错误}
\label{app:toolchain}

本附录介绍如何安装编译器、编译运行 C 程序、读懂报错，以及用调试器定位问题。不同系统细节不同，但核心命令是相通的。

\section{选择工具链}

\begin{center}
\begin{tabular}{@{}lll@{}}
\toprule
系统 & 推荐编译器 & 说明 \\
\midrule
Windows & MinGW-w64（GCC） & 免费，命令与 Linux 接近 \\
Linux & GCC / Clang & 通常已预装 \\
macOS & Clang & 安装 Xcode Command Line Tools \\
\bottomrule
\end{tabular}
\end{center}

本书推荐 \textbf{MinGW-w64 + VS Code}：编辑器轻量，GCC 的报错信息清晰，和后面的嵌入式开发环境也接近。

\section{Windows：安装 MinGW-w64}

最简单的方式是通过 MSYS2：

\begin{enumerate}
  \item 从 \url{https://www.msys2.org/} 下载并安装 MSYS2；
  \item 打开 MSYS2 UCRT64 终端，执行：
\begin{outbox}{MSYS2 命令}
pacman -S mingw-w64-ucrt-x86_64-gcc
\end{outbox}
  \item 把 \texttt{C:\textbackslash msys64\textbackslash ucrt64\textbackslash bin} 加入系统环境变量 \code{PATH}；
  \item 打开新的命令行，执行 \code{gcc --version}，能显示版本号就成功了。
\end{enumerate}

\section{第一个程序的编译与运行}

新建 \file{hello.c}：

\begin{codebox}{hello.c}
#include <stdio.h>

int main(void)
{
    printf("Hello, world!\n");
    return 0;
}
\end{codebox}

在文件所在目录打开终端，执行：

\begin{outbox}{编译与运行}
gcc hello.c -o hello -std=c89 -Wall -Wextra -pedantic
hello.exe        （Windows）
./hello          （Linux / macOS）
\end{outbox}

\begin{center}
\begin{tabular}{@{}ll@{}}
\toprule
参数 & 作用 \\
\midrule
\code{-o hello} & 指定输出文件名 \\
\code{-std=c89} & 按 C89 标准编译（本书要求） \\
\code{-Wall} & 打开常用警告 \\
\code{-Wextra} & 打开更多警告 \\
\code{-pedantic} & 严格按标准检查 \\
\code{-g} & 生成调试信息（配合 gdb） \\
\code{-lm} & 链接数学库（用了 \file{math.h} 时加在命令末尾） \\
\bottomrule
\end{tabular}
\end{center}

\begin{tipbox}[把警告当错误看]
初学阶段建议加 \code{-Wall -Wextra}。很多“能跑但结果不对”的 bug，编译器其实已经用警告提醒你了。更严格可以加 \code{-Werror}，把警告直接变成错误。
\end{tipbox}

\section{读懂错误信息}

GCC 的报错格式通常是：

\begin{outbox}{错误信息结构}
文件名:行号:列号: error: 错误描述
    源代码行
    ^
\end{outbox}

例如：

\begin{outbox}{示例}
hello.c:5:5: error: expected ';' before 'return'
     return 0;
     ^~~~~~
\end{outbox}

意思是第 5 行 \code{return} 前面少了一个分号。\textbf{先看第一条错误}，后面的错误常常是连锁反应。

\section{常见编译错误对照}

\begin{center}
\begin{tabular}{@{}p{0.42\textwidth}p{0.5\textwidth}@{}}
\toprule
错误信息 & 常见原因 \\
\midrule
\code{expected ';' before ...} & 上一行漏分号 \\
\code{expected declaration or statement at end of input} & 花括号不配对 \\
\code{'printf' undeclared} & 忘了 \texttt{\#include <stdio.h>} \\
\code{implicit declaration of function} & 函数没声明或头文件没包含 \\
\code{undefined reference to 'main'} & 没有 \code{main} 函数 \\
\code{undefined reference to 'sqrt'} & 链接时忘了 \code{-lm} \\
\code{ISO C90 forbids mixed declarations and code} & 声明写在了语句后面 \\
\texttt{stray '\textbackslash 357' in program} & 用了中文标点 \\
\code{variable-sized type used in sizeof} & C89 不允许变长数组 \\
\code{assignment makes pointer from integer} & 指针类型不匹配 \\
\bottomrule
\end{tabular}
\end{center}

\section{用 gdb 调试}

如果程序能编译但运行结果不对，或者崩溃，可以用 gdb。

\begin{outbox}{编译时加 -g，然后启动 gdb}
gcc bug.c -o bug -g -std=c89 -Wall
gdb bug.exe
\end{outbox}

常用命令：

\begin{center}
\begin{tabular}{@{}ll@{}}
\toprule
命令 & 作用 \\
\midrule
\code{run} & 运行程序 \\
\code{break main} & 在 main 函数处下断点 \\
\code{break 文件名:行号} & 在指定行下断点 \\
\code{next} & 执行下一行（不进入函数） \\
\code{step} & 执行下一行（进入函数） \\
\code{print 变量} & 打印变量的值 \\
\code{continue} & 继续运行到下一个断点 \\
\code{quit} & 退出 \\
\bottomrule
\end{tabular}
\end{center}

\begin{tipbox}[没有 gdb 怎么办]
用 \code{printf} 打印中间变量，同样能定位大部分问题。第 2 章介绍过这种方法。关键是：\textbf{在可疑的位置打印输入、中间结果和边界值}。
\end{tipbox}

\section{VS Code 配置}

安装 VS Code 和 C/C++ 扩展后，在项目目录新建 \file{.vscode/tasks.json}：

\begin{codebox}{tasks.json —— 编译任务}
{
    "version": "2.0.0",
    "tasks": [
        {
            "label": "build",
            "type": "shell",
            "command": "gcc",
            "args": [
                "-std=c89", "-Wall", "-Wextra", "-pedantic", "-g",
                "${file}", "-o", "${fileDirname}\\${fileBasenameNoExtension}.exe"
            ],
            "group": { "kind": "build", "isDefault": true }
        }
    ]
}
\end{codebox}

按 \code{Ctrl+Shift+B} 即可编译。运行可以在终端手动执行生成的 \file{.exe}。

\section{多文件项目}

\begin{outbox}{编译多个源文件}
gcc main.c math_utils.c -o app -std=c89 -Wall -Wextra
\end{outbox}

如果文件较多，可以写 Makefile 或使用 CMake。初学阶段先把“一个文件一个程序”练熟，再学分文件。

\section{编码问题}

\begin{itemize}
  \item 源文件统一保存为 \textbf{UTF-8}（无 BOM）；
  \item Windows 控制台默认可能是 GBK，输出中文可能乱码，可以执行 \code{chcp 65001} 切换到 UTF-8；
  \item 不要在代码里混用中文全角标点和英文半角标点，这是编译错误的常见来源。
\end{itemize}

\section{推荐的学习流程}

\begin{enumerate}
  \item 先手写代码，不要复制粘贴；
  \item 编译，打开所有警告；
  \item 有错先读第一条错误信息，按行号定位；
  \item 改完再编译，直到零警告；
  \item 运行，用边界值测试；
  \item 结果不对就打印中间变量或上调试器；
  \item 最后把代码整理干净，加上必要注释。
\end{enumerate}

\chapter{习题参考答案与提示}
\label{app:answers}

这里给出各章习题的参考答案。编程题只给关键思路或核心代码，完整程序请自己补全；如果和答案不同但能通过测试，也是好答案。

\section{第 1 章}

\begin{enumerate}
  \item $(1010)_2=10$，$(1111\,0000)_2=240$，$(1000\,0001)_2=129$。
  \item $13=(0000\,1101)_2=\mathrm{0D}_{16}$；$100=(0110\,0100)_2=\mathrm{64}_{16}$；$255=(1111\,1111)_2=\mathrm{FF}_{16}$。
  \item $+17=0001\,0001$，取反 $1110\,1110$，加 1 得 \code{1110 1111}。
  \item 取反得 $0001\,0010$，加 1 得 $0001\,0011=19$，所以原值为 $-19$。
  \item \code{1000 0000} 被规定为 $-128$；8 位没有足够的位表示 $+128$，所以正数最大只到 127。
  \item 用 \texttt{printf("\%d", (int)sizeof(int));} 等语句逐个打印，注意 \code{sizeof} 返回 \code{size_t}。
  \item 输出为 0（不相等）。因为 0.1 和 0.2 在二进制中无法精确表示，相加结果与 0.3 有微小误差，不能用 \code{==} 比较。
  \item 地址是内存单元的编号；字节序是多字节数据在内存中的排列顺序。小端把低字节放低地址，大端相反。
  \item 预处理（展开头文件和宏）→ 编译（生成汇编）→ 汇编（生成目标文件）→ 链接（生成可执行文件）。
  \item \code{unsigned char} 从 255 加 1 得到 0（模 256 回绕，行为明确）；\code{signed char} 从 127 加 1 是未定义行为，常见实现得到 $-128$，但不可依赖。
\end{enumerate}

\section{第 2 章}

\begin{enumerate}
  \item 用三次 \code{printf}，注意换行 \texttt{\textbackslash n}。
  \item \texttt{printf("\%d \%d \%d \%d \%d", a+b, a-b, a*b, a/b, a\%b);}，除法是整数除法。
  \item \code{c = (f - 32) * 5.0 / 9.0;}，用 \texttt{\%.2f} 输出。
  \item 周长 \code{2 * 3.1415926 * r}，面积 \code{3.1415926 * r * r}。
  \item 千位 \code{n / 1000}，百位 \texttt{n / 100 \% 10}，十位 \texttt{n / 10 \% 10}，个位 \texttt{n \% 10}，相加即可。
  \item 三步交换：\code{temp = a; a = b; b = temp;}。
  \item 错误：第一行 \code{int a = 1} 少分号；\code{sum} 未声明；\code{printf} 少分号。另外 \code{sum} 应先声明为 \code{int}。
  \item \texttt{\%d} 要求 \code{int}，而 3.14 是 \code{double}，属于未定义行为；应写 \texttt{printf("\%f", 3.14);}。
  \item 各位和：\texttt{n/100 + n/10\%10 + n\%10}；逆序：\texttt{(n\%10)*100 + (n/10\%10)*10 + n/100}。
  \item \code{scanf} 需要把输入写到变量所在的内存，所以必须传地址；\code{printf} 只需要读取变量的值，传值即可。
\end{enumerate}

\section{第 3 章}

\begin{enumerate}
  \item \code{7/2=3}，\texttt{7\%2=1}，\code{-7/2=-3}（向零截断），\texttt{-7\%2=-1}（常见实现；C89 未强制规定）。
  \item \texttt{n \% 2 != 0}，不要写 \code{== 1}（负数取余可能为 $-1$）。
  \item \texttt{n \% 3 == 0 \&\& n \% 5 != 0}。
  \item 判断第 0 位：\texttt{n \& 1}；第 2 位置 1：\code{n |= (1 << 2)}；第 3 位清零：\texttt{n \&= \textasciitilde{}(1 << 3)}。
  \item 赋值运算符右结合，等价于 \code{a = (b = (c + 1))}：先算 \code{c+1}，赋给 \code{b}，再赋给 \code{a}。
  \item \code{a = i++ + ++i;} 在同一个表达式中多次修改 \code{i}，属于未定义行为，结果不可预测。
  \item \texttt{c >= 'a' \&\& c <= 'z'}（或 \code{islower(c)}）。
  \item 比较时 \code{-1} 被转换成 \code{unsigned int}，变成 4294967295，所以 \code{-1 < 1U} 为假。
  \item \code{max = (a > b) ? ((a > c) ? a : c) : ((b > c) ? b : c);}。
  \item 左移超出类型宽度是未定义行为；有符号数左移溢出也可能是未定义行为。例如 \code{1 << 31} 对 32 位 \code{int} 是未定义的。
\end{enumerate}

\section{第 4 章}

\begin{enumerate}
  \item 用 \code{if (n > 0)}、\code{else if (n < 0)}、\code{else} 三种情况。
  \item \texttt{if ((year \% 4 == 0 \&\& year \% 100 != 0) || (year \% 400 == 0))}。
  \item 用字符范围判断：\texttt{c >= 'A' \&\& c <= 'Z'} 等；也可以包含 \file{ctype.h} 用 \code{isupper}、\code{islower}、\code{isdigit}。
  \item 用三个 \code{if} 和临时变量两两比较交换，最后按 \code{a,b,c} 输出。
  \item 先判断 \code{score < 0 || score > 100}，再用 \code{else if} 分档。
  \item 用 \code{switch (month)}，把 3\textasciitilde5 归为春、6\textasciitilde8 归为夏、9\textasciitilde11 归为秋、12/1/2 归为冬。
  \item 先处理 \code{a == 0}，再算判别式，分大于 0、等于 0、小于 0 三种情况。
  \item 先判断 \texttt{a+b>c \&\& a+c>b \&\& b+c>a}，再判断等边、等腰、普通。
  \item \code{if (a = 5)} 是赋值，条件永远为真；\code{if (a == 5)} 才是比较。
  \item \code{case} 必须是编译期常量，变量不是常量。可以用 \code{if-else} 链，或用查表/函数。
\end{enumerate}

\section{第 5 章}

\begin{enumerate}
  \item 三种循环都能写：\code{for (i=1;i<=100;i++) sum+=i;}、\code{while (i<=100) {...}}、\code{do {...} while (i<=100);}。
  \item \code{fact = fact * i;}。10! 能放进 \code{int}；13! = 6227020800 超出 32 位 \code{int}，会溢出。
  \item 读第一个数初始化 \code{max}、\code{min}，循环读其余数并更新，同时累加求平均。
  \item \texttt{if (i \% 3 == 0 || i \% 5 == 0)}，每输出一个计数，满 5 个换行。
  \item 试除到 \code{i*i <= n}；外层循环枚举 2\textasciitilde100。
  \item \code{f[0]=1; f[1]=1; for(i=2;i<30;i++) f[i]=f[i-1]+f[i-2];}。
  \item 辗转相除得到 gcd，lcm = \code{a / gcd * b}（先除后乘，避免溢出）。
  \item 枚举 100\textasciitilde999，拆出百、十、个位，判断立方和。
  \item 用 \code{do-while} 数位数；用 \texttt{\%10} 求各位和；用 \texttt{rev = rev*10 + n\%10} 求逆序。
  \item 数组下标范围是 \code{0..n-1}，\code{i <= n} 会访问 \code{a[n]}，越界。应改为 \code{i < n}。
\end{enumerate}

\section{第 6 章}

\begin{enumerate}
  \item 读入循环 \texttt{for (i=0;i<10;i++) scanf("\%d", \&a[i]);}，输出时用 \texttt{if (i\%5==4) printf("\textbackslash n");} 控制换行。
  \item 同第 5 章第 3 题，只是数据存在数组里。
  \item 双下标 \code{i=0, j=n-1}，\code{while (i<j)} 交换后向中间靠拢。
  \item 遍历数组，\texttt{if (a[i]\%2==0) even++; else odd++;}。
  \item \texttt{count[n\%10]++; n/=10;} 直到 \code{n==0}。
  \item 维护 \code{low}、\code{high}，每次比较 \code{a[mid]}，并统计循环次数。
  \item 用三重循环读入、相加、输出；转置输出 \code{a[j][i]}。
  \item 先求平均，再遍历输出 \code{a[i] > average} 的元素及其下标。
  \item \code{f[0]=1; f[1]=1;} 递推；第 30 项约为 832040，仍在 \code{int} 范围内，但再往后会溢出。
  \item 数组作为函数参数会退化成指针，\code{sizeof(a)} 得到的是指针大小；应把长度作为参数传入。
\end{enumerate}

\section{第 7 章}

\begin{enumerate}
  \item \code{mid.x = (p1.x + p2.x) / 2;}，\code{mid.y} 同理（注意整数除法）。
  \item \code{area = (r.x2 - r.x1) * (r.y2 - r.y1);}，周长 \code{2 * ((x2-x1) + (y2-y1))}。
  \item 遍历结构体数组累加、比较，输出平均分和最高分学生。
  \item 遍历数组，\code{if (s[i].score >= 60) pass++;}。
  \item 依次比较年、月、日；可以用 \code{if} 嵌套或把日期转成“距某天的天数”。
  \item 实部、虚部分别按复数公式计算。
  \item 因为对齐会产生填充字节，\code{sizeof} 往往大于各成员大小之和；例如 \code{char+int+double} 在典型平台上为 16 字节。
  \item 函数体里先执行 \texttt{p->score += x;}，再用 \texttt{if} 把超过 100 的成绩截断为 100。
  \item 定义位域后直接给字段赋值，打印各字段和 \code{sizeof}。
  \item 结构体可以整体赋值，编译器按成员逐个复制（包括填充字节）；数组名是常量，不能作为赋值目标。
\end{enumerate}

\section{第 8 章}

\begin{enumerate}
  \item \texttt{int *p = \&a;}，用 \code{*p} 输出值，\code{*p = 20;} 修改。
  \item 函数内用临时变量交换 \code{*a} 和 \code{*b}。
  \item \texttt{int *p; for (p = a; p < a + n; p++) printf("\%d ", *p);}。
  \item 双指针 \code{left = a, right = a + n - 1}，交换后向中间移动。
  \item 遍历一次同时维护 \code{max}、\code{min}、\code{sum}。
  \item \code{p->score += x;}，超过 100 时截断。
  \item \code{malloc} 后用冒泡排序，最后 \code{free(p); p = NULL;}。
  \item \code{a} 和 \texttt{\&a[0]} 都是首元素地址（\code{int *}）；\texttt{\&a} 是“整个数组的地址”，加 1 会跳过整个数组。
  \item \code{const int *p}：不能通过 p 改内容；\code{int * const p}：p 本身不能改指向，但内容可改。
  \item 数组参数退化为指针，函数内部无法用 \code{sizeof} 求原数组长度，所以必须额外传 \code{n}。
\end{enumerate}

\section{第 9 章}

\begin{enumerate}
  \item \code{strlen(s)}；自己实现用循环数到 \texttt{\textbackslash0}。
  \item \code{strcat(dst, src);}，确保 \code{dst} 足够大。
  \item \code{strcmp(a, b)}，返回 0 相等，小于 0 表示 a 在前。
  \item 遍历字符，用范围判断分类；注意跳过换行符。
  \item \texttt{if (c >= 'a' \&\& c <= 'z') c = c - 32;}。
  \item 双下标从两端向中间比较，全部相等则是回文。
  \item 用 \code{in_word} 标志，从空白进入非空白时计数加 1。
  \item 用读写双下标，遇到空格跳过，最后补 \texttt{\textbackslash0}。
  \item \code{count[s[i] - 'a']++;}，最后输出非零项。
  \item 字符数组在栈上复制了一份，可以修改；字符指针指向只读常量区，修改是未定义行为。
\end{enumerate}

\section{第 10 章}

\begin{enumerate}
  \item 冒泡排序外层控制轮数，内层相邻比较；每轮结束输出一次数组。
  \item 选择排序把“找最小”改成“找最大”，交换到前面。
  \item 插入排序中 \code{a[j] > key} 比较一次、\code{a[j+1]=a[j]} 移动一次，分别计数。
  \item 希尔排序把间隔从 \code{n/2} 开始每次减半。
  \item 归并排序先递归二分，再合并两个有序段；合并时输出临时数组。
  \item 快速排序选基准划分，递归左右；每趟划分后输出数组。
  \item 写 \code{cmp_student}，按 \code{score} 比较，调用 \code{qsort}。
  \item 先排序，再用 \code{low/high/mid} 二分查找并计数。
  \item 用 \code{clock()} 包住排序调用，计算 \code{(double)(end-start)/CLOCKS_PER_SEC}。
  \item 选择排序会“跨距离交换”，可能把两个相等元素的相对顺序打乱。例如 \code{5a, 5b, 3}，第一轮把 3 和 5a 交换，得到 \code{3, 5b, 5a}，5a 和 5b 的顺序变了。
\end{enumerate}

\section{第 11 章}

\begin{enumerate}
  \item \texttt{\#define PI 3.1415926}，\texttt{\#define E 2.71828}，\texttt{\#define GOLDEN 1.618}。
  \item \texttt{\#define SQUARE(x) ((x)*(x))}；\code{SQUARE(2+3)} 展开为 \code{((2+3)*(2+3))}=25。
  \item \texttt{\#define MAX3(a,b,c) MAX(MAX(a,b),c)}，并保证 \code{MAX} 的参数都加括号。
  \item \texttt{\#define SWAP(a,b) do \{ int t=(a); (a)=(b); (b)=t; \} while(0)}。
  \item \texttt{\#define PRINT\_VAR(x) printf("\%s = \%d\textbackslash n", \#x, (x))}。
  \item \texttt{\#define MAKE\_VAR(n) var\_\#\#n}，用 \code{MAKE_VAR(1)} 生成 \code{var_1}。
  \item 用 \texttt{\#ifdef DEBUG} 包住调试输出，未定义 \code{DEBUG} 时编译为空。
  \item 头文件写 \texttt{\#ifndef XXX\_H}、\texttt{\#define XXX\_H}、内容、\texttt{\#endif}。
  \item 不加括号时 \code{SQUARE(2+3)} 会展开成 \code{2+3*2+3}=11；应写成 \code{((x)*(x))}。
  \item \code{SQUARE(i++)} 展开为 \code{((i++)*(i++))}，\code{i} 被修改两次，属于未定义行为；先用临时变量保存 \code{i} 的值。
\end{enumerate}

\chapter{术语中英对照}
\label{app:glossary}

本附录按中文拼音/英文顺序整理书中出现的常用术语，方便查英文资料。

\begin{longtable}{@{}p{0.18\textwidth}p{0.24\textwidth}p{0.48\textwidth}@{}}
\toprule
中文 & English & 说明 \\
\midrule
\endfirsthead
\toprule
中文 & English & 说明 \\
\midrule
\endhead
\bottomrule
\endfoot
地址 & address & 内存单元的编号，指针里存的就是地址 \\
实参 & argument & 调用函数时实际传入的值 \\
数组 & array & 一段连续存放的同类型元素 \\
ASCII 码 & ASCII & 常用字符的 7 位编码标准 \\
赋值 & assignment & 把右边的值写进左边的变量 \\
二进制 & binary & 以 0 和 1 表示数字 \\
位 & bit & 最小的二进制单位，取值 0/1 \\
字节 & byte & 8 位，内存编址的基本单位 \\
缓冲区 & buffer & 临时存放输入输出数据的区域 \\
调用 & call & 执行一个函数 \\
强制类型转换 & cast & 显式地把一种类型转成另一种 \\
字符 & char & 单个字符，本质是小整数 \\
编译器 & compiler & 把 C 代码翻译成机器指令的程序 \\
常量 & constant & 程序运行中不改变的值 \\
CPU & central processing unit & 中央处理器 \\
悬空指针 & dangling pointer & 指向已释放内存的指针 \\
调试 & debug & 定位和修复程序错误的过程 \\
声明 & declaration & 告诉编译器“有这么个名字和类型” \\
定义 & definition & 为变量/函数分配存储或给出实现 \\
解引用 & dereference & 通过指针访问它指向的数据，用 \code{*} \\
双精度 & double & 双精度浮点数，通常 8 字节 \\
转义字符 & escape sequence & 反斜杠加字符表示特殊字符，如 \texttt{\textbackslash n} \\
表达式 & expression & 由变量、常量、运算符组成、有值的式子 \\
文件 & file & 存储在磁盘上的数据集合 \\
浮点数 & floating point & 带小数点的数，计算机用专门格式存储 \\
函数 & function & 完成特定任务的一段代码 \\
全局变量 & global variable & 定义在所有函数外的变量 \\
堆 & heap & 动态分配的内存区域，用 \code{malloc}/\code{free} 管理 \\
十六进制 & hexadecimal & 以 16 为基数，前缀 \code{0x} \\
标识符 & identifier & 变量、函数、类型等的名字 \\
初始化 & initialization & 定义变量时给初值 \\
整数 & integer & 没有小数部分的数 \\
关键字 & keyword & C 语言保留的特殊单词 \\
库 & library & 预先写好的函数集合 \\
链接器 & linker & 把目标文件和库组合成可执行文件 \\
局部变量 & local variable & 定义在函数或块内的变量 \\
循环 & loop & 重复执行一段代码 \\
宏 & macro & 预处理阶段的文本替换 \\
内存泄漏 & memory leak & 申请的内存没有释放 \\
空指针 & null pointer & 值为 \code{NULL} 的指针 \\
操作数 & operand & 参与运算的数据 \\
运算符 & operator & 表示运算的符号 \\
溢出 & overflow & 数值超出类型能表示的范围 \\
形参 & parameter & 函数定义中声明的参数 \\
指针 & pointer & 存放地址的变量 \\
预处理器 & preprocessor & 编译前处理 \texttt{\#} 指令的程序 \\
函数原型 & prototype & 函数的声明，说明返回值和参数类型 \\
递归 & recursion & 函数调用自身 \\
返回值 & return value & 函数执行后返回给调用者的值 \\
作用域 & scope & 名字有效的范围 \\
段错误 & segmentation fault & 访问非法内存导致的崩溃 \\
栈 & stack & 存放局部变量和函数调用信息的区域 \\
语句 & statement & 程序执行的基本单位，以分号结尾 \\
字符串 & string & 以 \texttt{\textbackslash0} 结尾的字符数组 \\
字符串常量 & string literal & 双引号括起来的字符序列 \\
结构体 & struct & 把不同类型成员组合成的类型 \\
语法 & syntax & 语言的书写规则 \\
补码 & two's complement & 计算机表示负数的一种编码方法（本书不展开） \\
未定义行为 & undefined behavior & 标准没有规定结果的行为，必须避免 \\
联合体 & union & 多个成员共用同一块内存的类型 \\
变量 & variable & 有名字、类型、地址的存储空间 \\
\code{void} & void & 表示“无类型/无返回值” \\
野指针 & wild pointer & 未初始化或指向无效地址的指针 \\
\end{longtable}

\begin{tipbox}[查英文资料的技巧]
遇到不认识的术语，先在本表查中文，再用英文关键词搜索，能找到更准确的结果。例如搜 \code{dangling pointer} 比搜“悬空指针”得到的资料多得多。
\end{tipbox}

\chapter{函数速查（补充）}
\label{app:function}

本书的章节表里没有单独安排“函数”一章，但后面的结构体、指针、排序都用到了函数。本附录把函数的必要语法补齐，够你读懂和编写本书的代码。

\section{函数的定义与调用}

\begin{codebox}{function-basic.c —— 定义一个函数}
#include <stdio.h>

int add(int a, int b)
{
    int sum;

    sum = a + b;
    return sum;
}

int main(void)
{
    int result;

    result = add(3, 5);
    printf("%d\n", result);

    return 0;
}
\end{codebox}

\begin{outbox}{运行结果}
8
\end{outbox}

函数定义由四部分组成：

\begin{itemize}
  \item \textbf{返回类型}：\code{int} 表示函数返回一个整数；不返回任何值写 \code{void}；
  \item \textbf{函数名}：\code{add}，遵循标识符命名规则；
  \item \textbf{参数列表}：\code{(int a, int b)}，没有参数写 \code{(void)}；
  \item \textbf{函数体}：花括号里的语句，用 \code{return} 返回结果。
\end{itemize}

\begin{pitfall}[C89 的参数声明]
C89 支持 \code{int add(int a, int b)} 这种写法，但 \code{return} 的类型和参数类型都必须写清楚。不要写 \code{add(a, b)} 这种“隐式 int”的老式写法。
\end{pitfall}

\section{返回值与 \texttt{void}}

\begin{codebox}{function-void.c —— 不返回值的函数}
#include <stdio.h>

void print_line(int n)
{
    int i;

    for (i = 0; i < n; i++) {
        printf("-");
    }
    printf("\n");
}

int main(void)
{
    print_line(20);
    return 0;
}
\end{codebox}

\code{void} 函数可以用 \code{return;} 提前结束，但不能返回一个值。

\begin{codebox}{function-early.c —— 提前返回}
int safe_divide(int a, int b)
{
    if (b == 0) {
        return 0;       /* 简单处理：除数为 0 返回 0 */
    }
    return a / b;
}
\end{codebox}

\section{函数声明（原型）}

如果函数定义在 \code{main} 之后，或者跨文件使用，需要先声明：

\begin{codebox}{function-prototype.c —— 函数原型}
#include <stdio.h>

int add(int a, int b);      /* 声明：告诉编译器有这样的函数 */

int main(void)
{
    printf("%d\n", add(3, 5));
    return 0;
}

int add(int a, int b)       /* 定义 */
{
    return a + b;
}
\end{codebox}

C89 不要求参数名与声明一致，但类型必须一致。养成“先声明、后使用”的习惯。

\section{参数传递}

\subsection{按值传递}

C 语言的参数默认是\textbf{按值传递}：函数拿到的是实参的副本，修改副本不影响原件。

\begin{codebox}{pass-value.c —— 按值传递}
#include <stdio.h>

void change(int x)
{
    x = 100;
}

int main(void)
{
    int a = 10;

    change(a);
    printf("%d\n", a);      /* 仍然是 10 */
    return 0;
}
\end{codebox}

\subsection{传地址}

要让函数修改调用者的变量，传指针：

\begin{codebox}{pass-pointer.c —— 传地址}
#include <stdio.h>

void change(int *p)
{
    *p = 100;
}

int main(void)
{
    int a = 10;

    change(&a);
    printf("%d\n", a);      /* 100 */
    return 0;
}
\end{codebox}

\subsection{数组参数}

数组作为参数会退化成指针，必须额外传长度：

\begin{codebox}{pass-array.c —— 数组参数}
int sum_array(int a[], int n)
{
    int i;
    int sum = 0;

    for (i = 0; i < n; i++) {
        sum = sum + a[i];
    }
    return sum;
}
\end{codebox}

\section{递归}

函数可以直接或间接调用自己，这叫递归。递归必须有\textbf{终止条件}，否则会无限递归导致栈溢出。

\begin{codebox}{recursion.c —— 递归求阶乘}
int factorial(int n)
{
    if (n <= 1) {
        return 1;               /* 终止条件 */
    }
    return n * factorial(n - 1); /* 递归调用 */
}
\end{codebox}

递归适合描述“问题可以分解成同样形式的小问题”，例如阶乘、斐波那契、归并排序、快速排序、树的遍历。递归的缺点是调用层次深时占用栈空间大。

\section{作用域与存储类别}

\begin{itemize}
  \item \textbf{局部变量}：定义在函数或块内，只在块内可见，函数返回后失效；
  \item \textbf{全局变量}：定义在所有函数之外，整个程序可见，默认初始化为 0；
  \item \code{static} 局部变量：只初始化一次，函数返回后值仍然保留；
  \item \code{static} 全局变量/函数：只在当前源文件可见。
\end{itemize}

\begin{codebox}{scope.c —— 作用域示例}
#include <stdio.h>

int global_count = 0;       /* 全局变量 */

void f(void)
{
    int local = 0;          /* 每次调用都重新初始化 */
    static int calls = 0;   /* 只初始化一次，跨调用保留 */

    local++;
    calls++;
    global_count++;

    printf("local=%d, calls=%d, global=%d\n",
           local, calls, global_count);
}
\end{codebox}

\section{多文件编程}

大型程序通常拆成多个源文件和头文件：

\begin{codebox}{math\_utils.h —— 头文件}
#ifndef MATH_UTILS_H
#define MATH_UTILS_H

int add(int a, int b);
int max_int(int a, int b);

#endif
\end{codebox}

\begin{codebox}{math\_utils.c —— 实现文件}
#include "math_utils.h"

int add(int a, int b)
{
    return a + b;
}

int max_int(int a, int b)
{
    return (a > b) ? a : b;
}
\end{codebox}

\begin{codebox}{main.c —— 使用函数}
#include <stdio.h>
#include "math_utils.h"

int main(void)
{
    printf("%d\n", add(2, 3));
    return 0;
}
\end{codebox}

编译命令：

\begin{outbox}{命令行}
gcc main.c math_utils.c -o app -std=c89 -Wall
\end{outbox}

\begin{tipbox}[头文件里只放声明]
头文件放函数声明、宏、类型定义；函数实现放在 \code{.c} 文件里。头文件用 \texttt{\#ifndef} 保护，避免重复包含。
\end{tipbox}

\section{函数的常见错误}

\begin{center}
\begin{tabular}{@{}p{0.38\textwidth}p{0.52\textwidth}@{}}
\toprule
错误 & 后果 \\
\midrule
声明与定义参数类型不一致 & 编译错误或运行异常 \\
调用函数前没有声明 & 隐式声明警告，C99 起报错 \\
函数没有 \code{return} 却使用返回值 & 未定义行为 \\
递归没有终止条件 & 栈溢出崩溃 \\
数组参数忘记传长度 & 越界访问 \\
返回局部数组的地址 & 悬空指针 \\
\bottomrule
\end{tabular}
\end{center}

\begin{tipbox}[函数设计的三条建议]
\begin{enumerate}
  \item 一个函数只做一件事；
  \item 函数名用动词，参数和返回值类型写清楚；
  \item 函数尽量短，超过一屏就考虑拆分。
\end{enumerate}
\end{tipbox}


\end{document}
